\section{The window bound: assembly, Lemma K and the minus part}\label{sec:pf:vwin}

\providecommand{\vwinQZ}{\Q_\ell/\Z_\ell}
\providecommand{\vwincyc}{\chi_{\mathrm{cyc}}}
\providecommand{\vwinsh}{\mathrm{shaLen}}
\providecommand{\vwinVt}{\widetilde V}
\providecommand{\vwinX}{\mathcal X}
\providecommand{\vwinmin}{\mathrm{minusLen}}
\providecommand{\vwinplus}{\mathrm{plusLen}}
\providecommand{\vwincok}{\lambda^{\mathrm{co}}}
\providecommand{\vwinVL}{\mathrm{VLen}}
\providecommand{\vwinexc}{\mathrm{Exc}}
\providecommand{\vwinpt}{\mathrm{pt}}

This section proves the window bound $\sfname{VWin}$ of \cref{def:window:vwin} from the statement
$\sfname{Row0b}$ on Galois representations attached to polarized regular algebraic cuspidal representations
(\cref{thm:vwin}), together with most of the intermediate results of \cref{sec:window} on which this deduction
rests. The deduction has three layers.
\begin{enumerate}[label=(\arabic*)]
\item \emph{From $\sfname{Row0b}$ to the package $\sfname{HypothesesH'}$} (\cref{sec:pf:vwin:top}). $\sfname{Row0b}$,
a base change theorem from definite unitary groups and a list of formalized published theorems are assembled into
the literature package $\sfname{LiteratureRankNHypothesis}$; together with a descent theorem to definite unitary
groups and further classical theorems on automorphic forms (Clozel's criteria, Galois representations for Hilbert
modular forms, quadratic and cyclic base change and descent, \dots) it gives $\sfname{HypothesesH'}$
(\cref{pf:vwin-he:H}). The main reduction $\sfname{HypothesesH'}\Rightarrow\sfname{VWin}$ is
\lref{UpperTheory.V.HV.DescentVanishingReductions.vwin_of_hypotheses_H_E_alt}, proved in
\cref{sec:pf:vwin-he}.
\item \emph{The final assembly} (\cref{sec:pf:vwin:assembly,sec:pf:vwin:plus,sec:pf:vwin:towers,sec:pf:vwin:Z,sec:pf:vwin:minus}). Given Lemma K and the four counting statements $\sfname{CountI}$, $\sfname{CountII}$,
$\sfname{CountIII}$, $\sfname{CountV}$ of \cref{pf:vwin-he:countdef}, the window bound follows from three partial
bounds, for the $2$-part, the minus parts and the plus parts of the groups $\Sha_D(j)$, and from elementary
estimates (\cref{pf:vwin:offour}). We prove: the numerical assembly (\cref{prop:window:assembly}); the bound for the
plus parts from the counts (\cref{thm:plus-part}); the bound for the $2$-part and the bad layers by Iwasawa theory in
cyclotomic towers and the choice of the window (\cref{thm:window:bab}, \cref{lem:window:tower}); the vanishing
statement for the trivial character (\cref{prop:window:propZ}); and the bound for the minus parts by an Euler system
of Gauss sums (\cref{prop:psibound}, \cref{lem:window:rec}).
\item \emph{Lemma K} (\cref{sec:pf:vwin:lemmaK}): an explicit, doubly exponential bound for the orders of the groups
$H^2(G_{L,S_\ell},\Zl(j+1))$, uniform in $\ell$ (\cref{thm:window:lemmaK}), proved from Poitou--Tate duality and a
rank bound of Borel type by a quantitative version of Soul\'e's method.
\end{enumerate}
Finally, \cref{sec:pf:vwin:patch} treats a Taylor--Wiles patching theorem for universal ordinary determinant rings,
which is used in the proofs of two automorphy statements of the plus part elsewhere in Part~II.

\emph{What is complete and what is cited.} The deductions in \cref{sec:pf:vwin:top,sec:pf:vwin:assembly,sec:pf:vwin:plus,sec:pf:vwin:towers} are given in full, down to statements which are either proved elsewhere in Part~II
(cited by their formal names), classical theorems of which the development contains formal proofs (cited from the
literature), or auxiliary lemmas of this part of the development which we state precisely and cite. In
\cref{sec:pf:vwin:minus} the reciprocity law, the Kolyvagin step and the induction are proved; the construction of the
Euler system of Gauss sums and the comparison of Stickelberger values with Bernoulli numbers are stated precisely and
cited. In \cref{sec:pf:vwin:lemmaK}, Lemma K is deduced from twenty-one intermediate statements and a rank bound of Borel
type; we prove fourteen of them, including all which carry the arithmetic (classes from divisors, counting, detection of
powers, the passage to $\Sha$ and to $H^2$, the numerics), and state precisely and cite the remaining combinatorial statements on the bar complex,
the Steinitz statement and the rank bound. In \cref{sec:pf:vwin:patch} the patching theorem is deduced from six statements
which we describe. \Cref{sec:pf:vwin:plain} contains a local lemma on cuspidal representations used elsewhere in Part~II.

\emph{Conventions.} We use the notation of \cref{sec:window} and of \cref{sec:pf:vwin-he}. For a number field $F$, an
algebraic closure $\Omega$ of $F$ and a prime $\ell$, an $\ell$-datum $D$, the groups $\Sha_D(j)$ and the lengths
$\vwinsh(D,j)\in\N\cup\{\infty\}$, $\vwinVL(D,j)=\len_{\Zl}H^2(G_{F,S_\ell},\Zl(j+1))$ are those of
\cref{def:window:sha}; for $F$ CM, $\vwinplus(D,j)$ and $\vwinmin(D,j)$ are the lengths of the plus and minus parts
(\cref{def:window:pm}). The length $\vwinsh(D,j)$ does not depend on $(\Omega,D)$; we write $\vwinsh(F,\ell,j)$ for it
(\lref{UpperTheory.V.LDataCardLeFinrank.shaLen_eq_shaInv}). A length $n\in\N\cup\{\infty\}$ enters real inequalities
through its natural-number part $n^\natural$, with $\infty^\natural=0$; this is why every bound below comes with a
separate finiteness assertion. All logarithms are natural; $\log_\ell$ denotes the integer part of the logarithm to base
$\ell$ ($\log_\ell 0=0$). For a number field $E$ let $\cond(E)$ be the least $N\ge1$ such that $E$ embeds into
$\Q(\zeta_N)$ ($\cond(E)=0$ if there is none).

% ---------------------------------------------------------------------------
\subsection{The intermediate statements}\label{sec:pf:vwin:stmts}

\begin{definition}[Abelian data and the fields $E_{M,k,q,d}$]\label{pf:vwin:fields}
An \emph{abelian datum} $M=(m_M,H_M)$ is an integer $m_M\ge1$ with a subgroup $H_M\subset(\Z/m_M)^\times$; its degree is
$\deg M=[(\Z/m_M)^\times:H_M]$. For a primitive Dirichlet character $\psi$ of conductor $N_\psi$ let
$M_\psi=(N_\psi,\ker\psi)$; for $m\ge1$ let $M^{\mathrm{cyc}}_m=(m,1)$. For $k\in\N$ and primes $q,d$, a number field $E$ is
an \emph{$E_{M,k,q,d}$} if there are $N$ divisible by $2^{k+2}$, $m_M$, $q$, $d$, an embedding $j:E\to\Q(\zeta_N)$ and
a primitive $N$-th root of unity $\zeta$ such that for every automorphism $\tau$ of $\Q(\zeta_N)$ and every $a\in\N$
with $\tau\zeta=\zeta^a$, $\tau$ fixes $j(E)$ if and only if $a\equiv\pm1\pmod{2^{k+2}}$, $a\bmod m_M\in H_M$, and
$\bigl(\frac aq\bigr)=\bigl(\frac ad\bigr)=1$ (Jacobi symbols). For $M=M_\psi$ this is the field $E_{\psi,k,q,d}$ of
\cref{def:window:fields}. A field $L$ is \emph{of type $M$} if there are $N$ divisible by $m_M$, $j:L\to\Q(\zeta_N)$ and
$\zeta$ as above such that $\tau$ fixes $j(L)$ if and only if $a\bmod m_M\in H_M$. We say that $(m;q,d)$ is a
\emph{prime pair} if $q\equiv3\pmod8$ and $d\equiv7\pmod{16}$ are
primes and $qd$ is prime to $m$.
\formallabel{UpperTheory.V.IsCompField}
\end{definition}

The statements $E=E_{M_\psi,k,q,d}$ and $E=E_{\psi,k,q,d}$ are equivalent
(\lref{UpperTheory.V.isCompField_ofChar}; the subgroup $H_{M_\psi}$ is the kernel of $\psi$ on units). By
\cref{lem:window:cm} every $E_{\psi,k,q,d}$ with $q\equiv3\pmod4$ is a CM field, and $1\le\cond(E)\le2^{k+2}N_\psi qd$
and $[E:\Q]\le\cond(E)$ (\lref{UpperTheory.V.isCMField_of_isChiField},
\lref{UpperTheory.V.conductor_bounds_of_isChiField}).

\begin{definition}[Partial bounds]\label{pf:vwin:partial}
Let $E$ be a number field, $t\in\N$ and $B\in\R$. $\sfname{VVat}(\ell,E,t,B)$ means: for every algebraic closure
$\Omega$ of $E$ and every $\ell$-datum $D$, $\vwinsh(D,j)<\infty$ for $0<j<t$ and
$\sum_{j=1}^{t-1}\vwinsh(D,j)^\natural\le B$. For $E$ CM, $\sfname{VVminusLe}(E,t,B)$ means: for every $\Omega$ and every
choice of $p$-data $D_p$ for all primes $p$, $\vwinmin(D_p,j)<\infty$ for odd $p$ and $0<j<t$, and
$\sum_{\ell\in S}\log\ell\sum_{j=1}^{t-1}\vwinmin(D_\ell,j)^\natural\le B$ for every finite set $S$ of odd primes;
$\sfname{VVplusLe}(E,t,B)$ is defined in the same way with $\vwinplus$.
\formallabel{UpperTheory.V.VVat}
\formallabel{UpperTheory.V.VVminusLe}
\formallabel{UpperTheory.V.VVplusLe}
\end{definition}

These are the statements of \cref{def:window:partial}. The window statements are expressed with \emph{admitted pairs}:
for $A_1,C_t\in\R$, a pair of functions $m,t:\R\to\N$ is $(A_1,C_t)$-admitted if $2\le m(x)$, $2^{m(x)}\le x^{A_1}$,
$2\le t(x)$ and $t(x)\le x^{C_t}$ for all $x\ge2$. For $J,X,Y\in\R$ we write $Y\in\mathrm{W}_J(X)$ if $X^{1/J}\le Y$
and $Y^2\le X$.

\begin{definition}[The partial statements]\label{pf:vwin:stmtdef}
\begin{enumerate}[label=(\alph*)]
\item\label{pf:vwin:stmt-minus} (Minus part.) For $C,o,n,N,t\in\N$ (with $C$ real) put
\[
 \beta(C,o,n,N,t)=3nt^2\log(tN)+C\sum_{\ell\mid o,\ \ell\ne2}\ell\,n^2t^2\log(t\ell N)^2 .
\]
$\sfname{ThmBcFree}$ is the statement: there is $C\in\R$ such that for every abelian datum $M$, all $k,t\ge2$, every
prime pair $(m_M;q,d)$ and every CM field $E=E_{M,k,q,d}$, $\sfname{VVminusLe}(E,t,\beta(C,\deg M,[E:\Q],\cond E,t))$.
Its admitted form $\sfname{ThmBcFree}(C_t)$ is: there is $C$ such that for every $M$, every $A_1$, every
$(A_1,C_t)$-admitted pair $(m,t)$, every $x\ge2$, every prime pair $(m_M;q,d)$ and every CM field $E=E_{M,m(x),q,d}$,
$\sfname{VVminusLe}(E,t(x),\beta(C,\deg M,[E:\Q],\cond E,t(x)))$.
\item\label{pf:vwin:stmt-const} (Constants.) A \emph{constant package} $K$ consists of functions $R_0$ (on abelian data) and
$R_\ell$ (on primes and abelian data) with values in $\N$, and a real $C_K>0$. For a finite set $B'$ of primes and abelian
data $L(\ell)$, $\ell\in B'$, put $R=R_0(M)+\sum_{\ell\in B'}R_\ell(L(\ell))$ and
$J_K(C_t,\eps;M,B',L)=2^{R+2}(2C_KC_t/\eps)^R$.
\item\label{pf:vwin:stmt-bab} (Two-part and bad layers.) For $x\in\R$ let $\mathrm{B}_a(M,\eps,m,t,q,d,x)$ be the statement that
every $E=E_{M,m(x),q,d}$ satisfies $\sfname{VVat}(2,E,t(x),\exp(x^{4\eps/5}))$, and $\mathrm{B}_b(K,B',L,\eps,t,x)$ the statement
that for every $\ell\in B'$, every field $L_0$ of type $L(\ell)$ and every extension $L_1/L_0$ generated by $\ell$-power
roots of unity with $[L_1:L_0]=\ell^{j_\ell}$, $j_\ell=\log_\ell R_\ell(L(\ell))+3$, one has
$\sfname{VVat}(\ell,L_1,t(x),\exp(x^{4\eps/5}))$. A real $Y$ is a \emph{$K$-window} for $(M,B',L,\eps,m,t,q,d,X)$ if
$Y\in\mathrm{W}_J(X)$ with $J=J_K(C_t,\eps;M,B',L)$ and $\mathrm{B}_a$ and $\mathrm{B}_b$ hold at every $x\in[Y,Y^2]$.
$\mathrm{B}(K,C_t)$ is the statement: for every abelian datum $M$, every finite set $B'$ of odd primes with abelian
data $L(\ell)$ such that $\ell\mid m_{L(\ell)}$ and every $a\in H_{L(\ell)}$ satisfies $a\equiv1\pmod\ell$ (so that the
fields of type $L(\ell)$ contain a primitive $\ell$-th root of unity), and all $\eps\in(0,\frac1{52})$ and $A_1\in\R$,
there is $X_0$ such that for every
$(A_1,C_t)$-admitted $(m,t)$, every $X\ge X_0$ and every prime pair $(m_M;q,d)$ with $\log(qd)\le X^{\eps/(4J)}$ there is
a $K$-window $Y$.
\item\label{pf:vwin:stmt-plus} (Plus part.) For a primitive $\psi$ let $B_\psi$ be the bad set of \cref{def:window:cases},
$L_\psi(\ell)=M^{\mathrm{cyc}}_{\ell N_\psi}$ and $J_\psi=J_K(C_t,\eps;M_\psi,B_\psi,L_\psi)$. $\mathrm{P}(K,C_t)$ is the
statement: for every primitive $\psi$ there are $C\in\R$ and primes $r_1,r_2$ such that for all $\eps\in(0,\frac1{52})$,
$A_1$ there is $X_0$ with the following property. For every $(A_1,C_t)$-admitted $(m,t)$, every $X\ge X_0$, every prime pair
$(N_\psi;q,d)$ with $q,d>2X^{C_t}+1$, with $q,d$ larger than each of $N_\psi,o_\psi,p_K(\psi),h^*(\psi),p_1(\psi),r_1,r_2$,
and with $\log(qd)\le X^{\eps/(4J_\psi)}$, every $K$-window $Y$ for $(M_\psi,B_\psi,L_\psi,\eps,m,t,q,d,X)$ and every
$x\in[Y,Y^2]$, every CM field $E=E_{\psi,m(x),q,d}$ satisfies
\begin{equation}\label{pf:vwin:hvat}
\sfname{VVplusLe}\bigl(E,t(x),(\cond(E)\,t(x))^C+C\exp(x^{4\eps/5})\bigr).
\end{equation}
\item\label{pf:vwin:stmt-joint} (Joint window.) $\mathrm{J}(C_t)$ is the statement: for every primitive $\psi$, all
$\eps\in(0,\frac1{52})$, $A_1$ and every $(A_1,C_t)$-admitted $(m,t)$ there are $J>0$ and $C,P,X_0$ such that for all
$X\ge X_0$ and all prime pairs $(N_\psi;q,d)$ with $q,d>\max(2X^{C_t}+1,P)$ and $\log(qd)\le X^{\eps/(4J)}$ there is
$Y\in\mathrm{W}_J(X)$ such that for every $x\in[Y,Y^2]$, $\mathrm{B}_a(M_\psi,\eps,m,t,q,d,x)$ holds and every CM field
$E=E_{\psi,m(x),q,d}$ satisfies \eqref{pf:vwin:hvat}.
\end{enumerate}
\formallabel{UpperTheory.V.minusBound}
\formallabel{UpperTheory.V.ThmBcFree}
\formallabel{UpperTheory.V.ThmBc}
\formallabel{UpperTheory.V.BConst}
\formallabel{UpperTheory.V.ThmBab}
\formallabel{UpperTheory.V.ThmHV}
\formallabel{UpperTheory.V.JointWin}
\end{definition}

The statement $\sfname{ThmBcFree}$ is hypothesis (a) of \cref{prop:window:assembly}, and $\mathrm{J}(C_t)$ is its
hypothesis (b); $\mathrm{B}(K,C_t)$ and $\mathrm{P}(K,C_t)$ are the conclusions of \cref{thm:window:bab,thm:plus-part}.


% ---------------------------------------------------------------------------
\subsection{From \texorpdfstring{$\sfname{Row0b}$}{Row0b} to the window bound}\label{sec:pf:vwin:top}

\begin{theorem}[\cref{thm:vwin}]\label{pf:vwin:main}
$\sfname{Row0b}$, with the local parameter predicate $\sfname{IsRecT}$ of \cref{sec:galois}, implies $\sfname{VWin}$.
\formallabel{UpperTheory.One.vwin_of_row}
\end{theorem}

The proof is a chain of three implications, each of which only assembles results proved elsewhere. We first state the
base change theorem which enters the literature package, then the two packages.

\begin{theorem}[Weak base change from definite unitary groups]\label{pf:vwin:bc}
Let $E$ be a CM field such that $E/E^+$ is unramified at every finite place, $n\ge1$, and let $\pi$ be an
irreducible automorphic representation of the definite unitary group $\Uni_n=\{g\in\GL_n(\A_E):{}^t\bar g\,g=1\}$.
Then there are a partition $n=N_1+\dots+N_c$ with $N_i\ge1$, discrete conjugate self-dual automorphic representations
$V_i$ of $\GL_{N_i}(\A_E)$ with local components $\sigma_{i,w}$ and infinitesimal characters $m_{i,\tau}$
($\tau:E\to\C$), and irreducible smooth representations $P_w$ of $\GL_n(E_w)$ for all finite places $w$, such that:
$P_w$ is a subquotient of the induced representation $\sigma_{1,w}\times\dots\times\sigma_{c,w}$; $P_w$ is spherical
whenever $\pi$ is unramified at $w$; at unramified inert places the Satake parameter of $P_w$ is the base change of that
of $\pi_w$; at split places $P_w$ is the component of $\pi$ at $w$ under $\Uni_n(E^+_v)\cong\GL_n(E_w)$; and
$\sum_im_{i,\tau}$ is the infinitesimal character attached to the weight of $\pi$ at $\tau$.
\formallabel{UpperTheory.BaseChangeIdentities.baseChangeN}
\end{theorem}

This is weak base change from definite unitary groups to $\GL_n$, a theorem of Labesse \cite{Labesse2011} (see also
\cite{ArthurClozel1989} for the cyclic case), proved there with the stabilized trace formula. The development
contains a formal proof which does not follow the trace formula route
(\lref{UpperTheory.BaseChangeIdentities.baseChangeN_of}, with \lref{UpperTheory.HeckeLevelFrameExistence.candFin}:
finiteness of systems of Hecke eigenvalues at bounded level, and a spectral accumulation argument for Eisenstein
contributions); we use only the statement.

\begin{lemma}[The literature package]\label{pf:vwin:litN}
$\sfname{Row0b}$ (with $\sfname{IsRecT}$) implies $\sfname{LiteratureRankNHypothesis}$, i.e.\ component
\ref{pf:vwin-he:H1} of \cref{pf:vwin-he:H}. Moreover $\sfname{LiteratureRankNHypothesis}$ implies $\sfname{VWin}$.
\formallabel{UpperTheory.One.Join.automorphic_window_of_inputs}
\formallabel{UpperTheory.One.Join.v_window_of_depth_zero_residual_bc_cm}
\formallabel{UpperTheory.One.Join.local_inputs_from_literature}
\formallabel{UpperTheory.One.Join.automorphic_inputs_from_literature}
\formallabel{UpperTheory.LiteratureTheorems.galois_cohomology_from_literature}
\formallabel{UpperTheory.One.Join.cm_abelian_inputs_from_literature}
\formallabel{UpperTheory.V.v_window_of_literal_automorphy_n}
\formallabel{UpperTheory.V.HV.DescentVanishingReductions.v_window_of_automorphy_csd_ss_mono_odd}
\end{lemma}

\begin{proof}
The package $\sfname{LiteratureRankNHypothesis}$ consists of four groups of statements (\cref{pf:vwin-he:H}, (a)--(d)).
Groups (a) (representation theory of $p$-adic $\GL_n$), (c) (Galois cohomology) and (d) (complex multiplication) do
not involve $\sfname{Row0b}$: each of their members is a theorem of the development, and most are formalized published
theorems (see \cref{sec:literature}). In group (b) (automorphic forms on $\GL_n$ and Galois representations) two
members are not of this kind: the base change statement, which is \cref{pf:vwin:bc}, and the statement on Galois
representations attached to polarized regular algebraic cuspidal representations, which is $\sfname{Row0b}$ with
$\sfname{IsRecT}$ itself. The other members of (b) are again theorems of the development. This proves the first
assertion. (Formally, $\sfname{Row0b}$ is first packaged as a structure with one field, whose only content is
$\sfname{Row0b}$; the theorems supplying the members of the four groups take this structure as an argument but do not
use it.)

For the second assertion let $\sfname{LiteratureRankNHypothesis}$ hold. The descent statement
$\sfname{DescWeak}$ of \cref{pf:vwin:hprime} below is a theorem (\cref{pf:vwin:descweak}), so \cref{pf:vwin:hprime}
gives $\sfname{HypothesesH'}$, and \lref{UpperTheory.V.HV.DescentVanishingReductions.vwin_of_hypotheses_H_E_alt}
($\sfname{HypothesesH'}$ implies $\sfname{VWin}$; this is \cref{thm:window:HE}, proved in \cref{sec:pf:vwin-he})
gives $\sfname{VWin}$.
\end{proof}

\begin{proof}[Proof of \cref{pf:vwin:main}]
Combine the two assertions of \cref{pf:vwin:litN}.
\end{proof}

\begin{remark}\label{pf:vwin:rem-names}
In this chain the only mathematical inputs besides theorems of the development are those of $\sfname{Row0b}$; in
particular no ``open image'' hypothesis is used, although the structure packaging $\sfname{Row0b}$ is named in this way
in the development. The descriptive names of the intermediate implications (``depth zero'', ``residual'', ``csd'',
``mono'', ``odd'') carry no additional content: each of them is a one-line application of the next.
\end{remark}

\subsubsection{The package \texorpdfstring{$\sfname{HypothesesH'}$}{HypothesesH'}}

We recall that $\sfname{HypothesesH'}$ is the conjunction of $\sfname{LiteratureRankNHypothesis}$ with the statements
listed in \ref{pf:vwin-he:H2} and \ref{pf:vwin-he:H3} of \cref{pf:vwin-he:H}. One of them is the descent statement
$\sfname{DescWeak}$: \emph{let $E$ be a CM field with $E/E^+$ unramified at all finite places, $p$ a prime, $S$ a finite set
of places of $E^+$ which split in $E$ and contains the places above $p$, $2\le n\le4$, and $\pi$ a regular algebraic
cuspidal representation of $\GL_n(\A_E)$ whose contragredient is its complex conjugate and which has a Satake parameter at
every place of $E$ not above $S$; then there is an automorphic representation $\sigma$ of the definite unitary group
$\Uni_n$ which is unramified at all non-split places of $E$ and such that for almost all split places $w$ the component of
$\sigma$ at $w$ (viewed as a representation of $\GL_n(E_w)$) has a Satake parameter equal to that of $\pi$ at $w$.}

\begin{proposition}\label{pf:vwin:descweak}
$\sfname{DescWeak}$ holds.
\formallabel{UpperTheory.V.Z.Tf.SR.descWeak}
\formallabel{UpperTheory.V.Z.Tf.SR.descWeakAll}
\formallabel{UpperTheory.V.Z.Tf.SR.has_of_split}
\end{proposition}

\begin{proof}
It suffices to prove the stronger statement in which $1\le n$ is arbitrary, the set $S$ is omitted and $\pi$ is only
assumed to have Satake parameters at the non-split places. For such $\pi$ the shared theorem
\lref{UpperTheory.V.Z.Tf.GaloisFromTraceFormula.runTh} (a consequence of a trace formula comparison) provides
either the required $\sigma$, or a \emph{decomposition datum}: a finite family $(P_i)$ of representations of groups
$\GL_{N_i}(\A_E)$ with $\sum_iN_i=n$, regular in the sense of the development, such that at almost all split places the
sum of the Satake parameters of the $P_i$ equals that of $\pi$, but such that $\pi$ is not one of the $P_i$. The second
alternative is impossible: by the shared rigidity theorem for split places \lref{UpperTheory.V.Z.Tf.SR.splitRigid},
applied after normalizing $\pi$ to be unitary, a regular decomposition datum which agrees with $\pi$ at almost all split
places contains $\pi$. The statement of $\sfname{DescWeak}$ follows, since the places of $E$ above $S$ are split.
\end{proof}

$\sfname{DescWeak}$ is a descent statement of the type proved by Labesse \cite{Labesse2011} for definite unitary groups.

\begin{proposition}[The package $\sfname{HypothesesH'}$]\label{pf:vwin:hprime}
$\sfname{LiteratureRankNHypothesis}$ and $\sfname{DescWeak}$ imply $\sfname{HypothesesH'}$. More precisely, the
components of \ref{pf:vwin-he:H2} and \ref{pf:vwin-he:H3} are supplied as follows.
\begin{enumerate}
\item[(a)] Clozel's criterion: a cuspidal representation of $\GL_a(\A_E)$, $E$ totally imaginary, whose infinitesimal
character is that of an irreducible algebraic representation, is cohomological for it \cite{Clozel1990}; deduced from
$\sfname{LiteratureRankNHypothesis}$.
\item[(b)] Existence of $\Aut(\C)$-conjugates of regular algebraic cuspidal representations \cite{Clozel1990}, and Borel's
bound for the number of cohomological cuspidal representations of given level by the dimension of the cohomology of the
locally symmetric space \cite{Borel1983}: labelled in the development as published results.
\item[(c)] Galois representations with local--global compatibility for regular algebraic polarized cuspidal
representations of $\GL_2$ over totally real fields \cite{Carayol1986,Taylor1989}; deduced from
$\sfname{LiteratureRankNHypothesis}$ (in particular from $\sfname{Row0b}$) by quadratic base change (e) to a totally
imaginary quadratic extension, using the admissibility statement of (d).
\item[(d)] Strong multiplicity one, admissibility, and Gelfand--Kazhdan duality (the contragredient of a cuspidal
representation of $\GL_a$ over a totally real field is its transpose-inverse) \cite{JacquetShalika1981}: the first two are
labelled as published results, the third is deduced from strong multiplicity one.
\item[(e)] Quadratic base change for $\GL_2$ from a totally real field to a totally imaginary quadratic extension
\cite{Langlands1980}, obtained from a formal proof of quadratic base change.
\item[(f)] The cuspidal representation of $\GL_2(\A_\Q)$ attached to a classical newform of weight $\ge2$; deduced from
the semisimplicity of the space of cusp forms (j) and admissibility (d).
\item[(g)] $\sfname{DescWeak}$: the hypothesis.
\item[(h)] Cyclic base change and descent of prime degree for $\GL_n$ over totally imaginary fields \cite{ArthurClozel1989}.
\item[(i)] Descent to $\Q$ of regular algebraic cuspidal representations of $\GL_2$ over an imaginary quadratic field with
sufficiently regular infinitesimal character and Satake parameters invariant under the Galois conjugation at almost all
places \cite{Langlands1980}; this proves the first alternative of the corresponding component of \ref{pf:vwin-he:H3}.
\item[(j)] Existence of a totally definite quaternion algebra over every totally real field of even degree
\cite{Vigneras1980}, and semisimplicity of the space of cusp forms over totally real fields: labelled as published
results.
\end{enumerate}
\formallabel{UpperTheory.V.HV.DescentVanishingReductions.published_hypotheses_from_patching_alt}
\formallabel{UpperTheory.V.Af.clozelCohC_alt}
\formallabel{UpperTheory.V.Z.Coh.clozelConjC}
\formallabel{UpperTheory.V.Z.Coh.borelInjC}
\formallabel{UpperTheory.Gal.SV.galois_rep_exists_totally_real_rank_two_of_base_change}
\formallabel{UpperTheory.V.Af.strongMultOneReal}
\formallabel{UpperTheory.V.Af.gelfandKazhdanReal}
\formallabel{UpperTheory.ConvolutionTheory.R2.quadBC2}
\formallabel{UpperTheory.V.Br.quadBC0_of_two}
\formallabel{UpperTheory.V.Af.newformRepOf}
\formallabel{UpperTheory.V.Z.Bc.cycC_holds}
\formallabel{UpperTheory.ConvolutionTheory.RankTwoDescent.descS}
\formallabel{UpperTheory.V.Af.defQuatExists}
\formallabel{UpperTheory.V.Af.cuspSemisimpleReal}
\end{proposition}

\begin{proof}
$\sfname{HypothesesH'}$ is a conjunction; each conjunct is the statement named in the corresponding item, supplied by
the theorems of the development whose formal names are listed after the first tag above (in the order of the items),
the admissibility statement being the shared \lref{UpperTheory.V.Af.admissibleReal}. All of these are known
theorems with the cited references; the development contains formal proofs of (a), (c), the third part of (d), (e),
(f), (h) and (i), which we do not reproduce (the descent theorem (i) alone occupies some $19{,}000$ lines). The proof of
the proposition consists in collecting them.
\end{proof}

\begin{remark}\label{pf:vwin:rem-hprime}
(1) In (i) the development proves the first of the two alternatives offered by \ref{pf:vwin-he:H3} (descent to $\Q$);
the alternative Galois-theoretic form is not needed. (2) The formal name of \cref{pf:vwin:hprime} refers to patching,
but its proof is a bundling step: no patching argument enters at this point.
\end{remark}


% ---------------------------------------------------------------------------
\subsection{The final assembly}\label{sec:pf:vwin:assembly}

We now prove the root of the main reduction which is not proved in \cref{sec:pf:vwin-he}: the deduction of
$\sfname{VWin}$ from the literature package, Lemma K and the four counting statements. Recall from
\cref{pf:vwin-he:countdef} the statements $\mathrm{Count}(P)$ for properties $P$ of tuples $(h,\ell,\OO,j,\theta)$, and
that $\sfname{CountZ}=\mathrm{Count}(P)$ for $P$: ``$\ell\notin B_\chi$, $(\ell-1)\mid2j$ and $\theta=1$'' (case Z of
\cref{tab:window:cases}).

\begin{theorem}\label{pf:vwin:offour}
$\sfname{LiteratureRankNHypothesis}$, the existence of $C_K$ with $\sfname{LemmaK}(C_K)$, and $\sfname{CountI}$,
$\sfname{CountII}$, $\sfname{CountIII}$, $\sfname{CountV}$ imply $\sfname{VWin}$.
\formallabel{UpperTheory.V.HV.Core.of_four_alt}
\end{theorem}

\begin{theorem}\label{pf:vwin:counts}
Assume Poitou--Tate duality, the local finiteness and local Tate duality statements over finite extensions of $\Ql$
(members of group (c) of $\sfname{LiteratureRankNHypothesis}$), $\sfname{LemmaK}(C_K)$, $\mathrm{B}(K,401)$ for a
constant package $K$, $\sfname{ThmBcFree}$, $\sfname{CountI}$, $\sfname{CountII}$, $\sfname{CountIII}$,
$\sfname{CountV}$ and $\sfname{CountZ}$. Then $\sfname{VWin}$ holds.
\formallabel{UpperTheory.V.HV.vwin_of_counts}
\formallabel{UpperTheory.V.HV.jointWin_of_counts}
\end{theorem}

\begin{proof}[Proof of \cref{pf:vwin:offour}, given \cref{pf:vwin:counts}]
Poitou--Tate duality in the form used below follows from the normalized form contained in
$\sfname{LiteratureRankNHypothesis}$ (\lref{UpperTheory.V.HV.U.pt_s}); the local statements and the orthogonality of
unramified classes are members of its group (c). By \cref{thm:window:bab} (\cref{pf:vwin:bab} below), Poitou--Tate
duality, the local statements, the theorem of Ferrero--Washington (\cref{lit:window:fw},
\lref{UpperTheory.V.FW.ferreroWashington}) and Lemma K give a constant package $K$ with $\mathrm{B}(K,C_t)$ for
all $C_t\ge1$, in particular for $C_t=401$. By \cref{prop:psibound} (\cref{pf:vwin:psibound} below),
Poitou--Tate duality and the orthogonality of unramified classes give $\sfname{PsiBound}$ and then
$\sfname{ThmBcFree}$. By \cref{prop:window:propZ} (\cref{pf:vwin:propZ} below), $\sfname{CountZ}$ holds. Now
\cref{pf:vwin:counts} applies.
\end{proof}

\begin{proof}[Proof of \cref{pf:vwin:counts}]
By \cref{thm:plus-part} (\cref{pf:vwin:hv} below) the hypotheses give $\mathrm{P}(K,401)$; with $\mathrm{B}(K,401)$
this gives $\mathrm{J}(401)$ (\cref{pf:vwin:jointof}); $\sfname{ThmBcFree}$ gives $\sfname{ThmBcFree}(401)$
(\cref{pf:vwin:bciff}); and \cref{pf:vwin:assembly} gives $\sfname{VWinAt}(\psi,401)$ for every primitive $\psi$, which
is $\sfname{VWin}$.
\end{proof}

\subsubsection{Splitting and elementary estimates}

\begin{lemma}\label{pf:vwin:parts}
Let $E$ be a CM field, $t\in\N$ and $B_2,B_-,B_+\in\R$. If $\sfname{VVat}(2,E,t,B_2)$, $\sfname{VVminusLe}(E,t,B_-)$ and
$\sfname{VVplusLe}(E,t,B_+)$ hold, then $\sfname{VVle}(E,t,B_2\log2+B_-+B_+)$ holds.
\formallabel{UpperTheory.V.vvle_of_parts}
\end{lemma}

\begin{proof}
Fix $\Omega$ and $p$-data $D_p$. For odd $p$ and $j\ge1$, $\vwinsh(D_p,j)=\vwinplus(D_p,j)+\vwinmin(D_p,j)$ by
\cref{lem:window:pm} (\lref{UpperTheory.V.shaLen_eq_add}); hence the finiteness assertions of the three hypotheses
give that of $\sfname{VVle}$, and for $0<j<t$ the natural-number parts add. Let $S$ be a finite set of primes and
$g(\ell)=\log\ell\sum_{j=1}^{t-1}\vwinsh(D_\ell,j)^\natural\ge0$. The contribution of $\ell=2$ (if $2\in S$) is at most
$B_2\log 2$; that of the odd primes of $S$ is the sum of the corresponding expressions for $\vwinmin$ and $\vwinplus$, hence
at most $B_-+B_+$.
\end{proof}

\begin{lemma}[Monotonicity]\label{pf:vwin:vvlemono}
Let $E$ be a number field, $t'\le t$ in $\N$ and $B\le B'$ in $\R$. If $\sfname{VVle}(E,t,B)$ holds, then
$\sfname{VVle}(E,t',B')$ holds.
\formallabel{UpperTheory.Top.VVle.mono}
\end{lemma}

\begin{proof}
Fix $\Omega$ and $p$-data $D_p$ for all primes $p$. The finiteness of $\vwinsh(D_p,j)$ for $0<j<t'$ is part of the
hypothesis, since $t'\le t$. For a finite set $S$ of primes, the terms $\vwinsh(D_\ell,j)^\natural$ with $t'\le j<t$ are
non-negative and $\log\ell\ge0$, so
$\sum_{\ell\in S}\log\ell\sum_{j=1}^{t'-1}\vwinsh(D_\ell,j)^\natural\le\sum_{\ell\in S}\log\ell\sum_{j=1}^{t-1}
\vwinsh(D_\ell,j)^\natural\le B\le B'$.
\end{proof}

\begin{lemma}[Elementary estimates]\label{pf:vwin:est}
Let $\eps\in(0,\frac1{52})$.
\begin{enumerate}[label=(\alph*)]
\item\label{pf:vwin:est-adm} If $C_t\ge1$ and $m:\R\to\N$ satisfies $\sfname{PageN}(\eps,x,m(x))$ for all $x\ge1$, then
$(m,x\mapsto\lfloor x^{C_t}\rfloor)$ is $(4,C_t)$-admitted.
\item\label{pf:vwin:est-big} For $J>0$, $C_t$ and $P$ there is $X_0$ such that for $X\ge X_0$, every pair $(q,d)$ with
$\sfname{WinPair}(\eps,X^{1/J};q,d)$ satisfies $q,d>2X^{C_t}+1$ and $q,d>P$.
\item\label{pf:vwin:est-log} There is $Y_0$ such that $\log(qd)\le Y^{\eps/4}$ whenever $Y\ge Y_0$ and
$\sfname{WinPair}(\eps,Y;q,d)$.
\item\label{pf:vwin:est-cond} For $C_t\in\R$ and $m_0\in\N$ there is $x_0$ such that $\log(Nt)\le2x^{\eps/4}$ whenever $x\ge x_0$,
$1\le N\le2^{k+2}m_0qd$, $\sfname{PageN}(\eps,x,k)$, $1\le t\le x^{C_t}$ and $\log(qd)\le x^{\eps/4}$.
\item\label{pf:vwin:est-beta} For $C\in\R$ and $o\in\N$ put $K=3+|C|\sum_{\ell\mid o,\ \ell\ne2}\ell^3$. If $1\le N$, $1\le t$, $n\le N$
and $\log(Nt)\le2W$, then $\beta(C,o,n,N,t)\le Ke^{8W}$.
\item\label{pf:vwin:est-err} For $K,C\in\R$ there is $x_0$ such that for $x\ge x_0$ and every $u\ge1$ with $\log u\le2x^{\eps/4}$,
\[
 e^{x^{4\eps/5}}\log2+Ke^{8x^{\eps/4}}+u^C+Ce^{x^{4\eps/5}}\ \le\ \sqrt{\exp(x^\eps)} .
\]
\end{enumerate}
\formallabel{UpperTheory.V.admissible_of_pageN}
\formallabel{UpperTheory.V.eventually_winPair_gt_power}
\formallabel{UpperTheory.V.log_mul_le_rpow_of_winPair}
\formallabel{UpperTheory.V.log_conductor_mul_height_le_of_pageN}
\formallabel{UpperTheory.V.minusBound_le_exp_of_log_le}
\formallabel{UpperTheory.V.error_terms_le_sqrt_exp_rpow}
\end{lemma}

\begin{proof}
\ref{pf:vwin:est-adm} Let $x\ge2$. From $2^{m(x)}\ge x^{2+8\eps}\ge4$ we get $m(x)\ge2$; from $2^{m(x)}\le2x^{2+8\eps}$, $2\le x$ and
$8\eps<1$ we get $2^{m(x)}\le x^{3+8\eps}\le x^4$; and $\lfloor x^{C_t}\rfloor\ge\lfloor x\rfloor\ge2$.
\ref{pf:vwin:est-big} By $\sfname{WinPair}$, $\log q>\frac12X^{\eps/(5J)}$, and $\log d>\frac12X^{\eps/(5J)}(\log X^{1/J})^2\ge\frac12
X^{\eps/(5J)}$ once $X\ge e^J$; and $\frac12X^{\eps/(5J)}$ exceeds $\log(2X^{C_t}+1+|P|)$ for large $X$.
\ref{pf:vwin:est-log} $\log(qd)\le Y^{\eps/5}(1+(\log Y)^2)\le Y^{\eps/4}$ for large $Y$.
\ref{pf:vwin:est-cond} By $\sfname{PageN}$, $2^{k+2}\le8x^{2+8\eps}$, so $\log(Nt)\le\log8+(2+8\eps+C_t)\log x+\log m_0+x^{\eps/4}$, which is
at most $2x^{\eps/4}$ for large $x$.
\ref{pf:vwin:est-beta} Put $X=Nt\ge1$; then $X\le e^{2W}$, so $X^4\le e^{8W}$. Since $\log y\le y$ for $y>0$, $\log(tN)\le X$ and
$\log(t\ell N)\le\ell X$. Hence $3nt^2\log(tN)\le3X^3\le3X^4$ and $\ell n^2t^2\log(t\ell N)^2\le\ell X^2(\ell X)^2=\ell^3X^4$, and the
claim follows by summing.
\ref{pf:vwin:est-err} We have $u^C\le e^{2|C|x^{\eps/4}}$. Each of the four terms is at most a constant times
$\exp(\max(x^{4\eps/5},\,8x^{\eps/4},\,2|C|x^{\eps/4}))$, and since $\frac{4\eps}5<\eps$ and $\frac\eps4<\eps$ the sum is at most
$\exp(\frac12x^\eps)=\sqrt{\exp(x^\eps)}$ for large $x$.
\end{proof}

\begin{lemma}\label{pf:vwin:bciff}
For $C_t\ge1$, $\sfname{ThmBcFree}(C_t)$ is equivalent to $\sfname{ThmBcFree}$.
\formallabel{UpperTheory.V.thmBc_iff}
\end{lemma}

\begin{proof}
If $\sfname{ThmBcFree}$ holds with the constant $C$ and $(m,t)$ is admitted, then $m(x),t(x)\ge2$ for $x\ge2$, and
$\sfname{ThmBcFree}$ applies with $k=m(x)$ and $t=t(x)$. Conversely let $\sfname{ThmBcFree}(C_t)$ hold with $C$, and let
$k,t\ge2$, a prime pair $(m_M;q,d)$ and $E=E_{M,k,q,d}$ be given. Put $x=\max(2^k,t)\ge2$ and let $m(y)=k$, $t(y)=t$ for
$y=x$ and $m(y)=t(y)=2$ otherwise. This pair is $(2,C_t)$-admitted: at $y=x$ one has $2^k\le x\le x^2$ and $t\le x\le x^{C_t}$,
and for $y\ne x$, $y\ge2$, one has $4\le y^2$ and $2\le y\le y^{C_t}$. Applying $\sfname{ThmBcFree}(C_t)$ at $x$ gives the
assertion of $\sfname{ThmBcFree}$ for $(M,k,t,q,d,E)$.
\end{proof}

\begin{lemma}\label{pf:vwin:jointof}
Let $C_t>0$ and let $K$ be a constant package. Then $\mathrm{B}(K,C_t)$ and $\mathrm{P}(K,C_t)$ imply $\mathrm{J}(C_t)$.
\formallabel{UpperTheory.V.jointWin_of}
\end{lemma}

\begin{proof}
Let $\psi$, $\eps$, $A_1$ and an admitted pair $(m,t)$ be given. The data $(M_\psi,B_\psi,L_\psi)$ satisfy the hypothesis
of $\mathrm{B}(K,C_t)$: every $\ell\in B_\psi$ is an odd prime, $\ell$ divides $m_{L_\psi(\ell)}=\ell N_\psi$, and
$H_{L_\psi(\ell)}$ is trivial (\lref{UpperTheory.V.primeSetBoundData_of_primitiveChar}). Let $X_0$ be the threshold of
$\mathrm{B}(K,C_t)$ for these data, and $C,r_1,r_2,X_1$ those of $\mathrm{P}(K,C_t)$. Take $J=J_\psi>0$, the constant $C$,
$P=N_\psi+o_\psi+p_K(\psi)+h^*(\psi)+p_1(\psi)+r_1+r_2$ and the threshold $\max(X_0,X_1)$. If $X$ and $(q,d)$ satisfy the
hypotheses of $\mathrm{J}(C_t)$, then $\mathrm{B}(K,C_t)$ provides a $K$-window $Y$; for $x\in[Y,Y^2]$ it gives
$\mathrm{B}_a$, and $\mathrm{P}(K,C_t)$ gives \eqref{pf:vwin:hvat}, because each of the listed integers is at most $P<q,d$.
\end{proof}

\begin{proposition}[\cref{prop:window:assembly}]\label{pf:vwin:assembly}
Let $C_t\ge1$. If $\mathrm{J}(C_t)$ and $\sfname{ThmBcFree}$ hold, then $\sfname{VWinAt}(\psi,C_t)$ holds for every
primitive $\psi$. In particular $\mathrm{J}(401)$ and $\sfname{ThmBcFree}$ imply $\sfname{VWin}$.
\formallabel{UpperTheory.V.vwinAt_of_joint}
\formallabel{UpperTheory.V.vwin_of_joint}
\end{proposition}

\begin{proof}
By \cref{pf:vwin:bciff} we may use $\sfname{ThmBcFree}(C_t)$, with a constant $C_b$. Fix $\psi$, $\eps\in(0,\frac1{52})$ and
$m:\R\to\N$ with $\sfname{PageN}(\eps,x,m(x))$ for $x\ge1$, and put $t(x)=\lfloor x^{C_t}\rfloor$. By
\cref{pf:vwin:est}\ref{pf:vwin:est-adm}, $(m,t)$ is $(4,C_t)$-admitted, so $\mathrm{J}(C_t)$ (with $A_1=4$) gives $J>0$ and
$C,P,X_J$. Choose: $K$ as in \cref{pf:vwin:est}\ref{pf:vwin:est-beta} for $(C_b,\deg M_\psi)$; $x_T$ as in
\ref{pf:vwin:est-err} for $(K,C)$; $x_N$ as in \ref{pf:vwin:est-cond} for $(C_t,N_\psi)$; $Y_0$ as in \ref{pf:vwin:est-log};
$J'=\max(J,1)$ and $P'=\max(P,N_\psi)$; $X_B$ as in \ref{pf:vwin:est-big} for $(J',C_t,P')$; and $L=\max(x_T,x_N,Y_0,2)$.
We show that the definition of $\sfname{VWinAt}(\psi,C_t)$ is satisfied with the constant $J'\ge1$ and the threshold
$X_0=\max(X_J,X_B,L^{J'})$.

Let $X\ge X_0$ and let $(q,d)$ satisfy $\sfname{WinPair}(\eps,X^{1/J'};q,d)$. Then $L\le X^{1/J'}$. By
\ref{pf:vwin:est-big}, $q$ and $d$ exceed $2X^{C_t}+1$, $P$ and $N_\psi$; as they are primes larger than $N_\psi$, they are
prime to $N_\psi$, and with the congruences of $\sfname{WinPair}$, $(N_\psi;q,d)$ is a prime pair. By \ref{pf:vwin:est-log}
applied to $Y'=X^{1/J'}\ge Y_0$,
\begin{equation}\label{pf:vwin:logqd}
 \log(qd)\le X^{\eps/(4J')}\le X^{\eps/(4J)} .
\end{equation}
Hence $\mathrm{J}(C_t)$ gives $Y$ with $X^{1/J}\le Y$ and $Y^2\le X$; then $X^{1/J'}\le X^{1/J}\le Y$, as required. Let
$x\in[Y,Y^2]$; then $x\ge L\ge2$, $x\ge x_T$ and $x\ge x_N$. Let $E=E_{\psi,m(x),q,d}$; it is a CM field (as $q\equiv3\pmod4$) and an
$E_{M_\psi,m(x),q,d}$. By $\mathrm{J}(C_t)$, $\sfname{VVat}(2,E,t(x),e^{x^{4\eps/5}})$ and
$\sfname{VVplusLe}(E,t(x),(\cond(E)t(x))^C+Ce^{x^{4\eps/5}})$ hold, and $\sfname{ThmBcFree}(C_t)$ gives
$\sfname{VVminusLe}(E,t(x),\beta(C_b,o_\psi,[E:\Q],\cond E,t(x)))$. Now $1\le\cond E\le2^{m(x)+2}N_\psi qd$,
$[E:\Q]\le\cond E$, $1\le t(x)\le x^{C_t}$, and $\log(qd)\le x^{\eps/4}$ by the first inequality of \eqref{pf:vwin:logqd} and
$X^{1/J'}\le x$. So \ref{pf:vwin:est-cond} gives $\log(\cond(E)t(x))\le2x^{\eps/4}$, and \ref{pf:vwin:est-beta} with
$W=x^{\eps/4}$ bounds $\beta(\dots)$ by $Ke^{8x^{\eps/4}}$. By \cref{pf:vwin:parts}, $\sfname{VVle}(E,t(x),B)$ holds with
$B=e^{x^{4\eps/5}}\log2+Ke^{8x^{\eps/4}}+(\cond(E)t(x))^C+Ce^{x^{4\eps/5}}$, and \ref{pf:vwin:est-err} with $u=\cond(E)t(x)\ge1$
shows $B\le\sqrt{\exp(x^\eps)}$. By \cref{pf:vwin:vvlemono}, $\sfname{VVle}(E,\lfloor x^{C_t}\rfloor,
\sqrt{\exp(x^\eps)})$ holds, which is the assertion of $\sfname{VWinAt}$ at $x$. The last statement is the case
$C_t=401$, since $\sfname{VWin}$ is $\sfname{VWinAt}(\psi,401)$ for all $\psi$.
\end{proof}


% ---------------------------------------------------------------------------
\subsection{The plus part}\label{sec:pf:vwin:plus}

In this subsection we prove \cref{thm:plus-part}. We use the hosts $h=(\chi,m,q_1,d)$, the integers $N_h$, the subgroups
$H_h\subset(\Z/N_h)^\times$, the coefficient rings $\OO$ with ramification index $e_\OO$, the sets $\vwinX_h(\OO,j)$ of
characters, the lengths $V_\ell(\theta,j)$ and $\vwinVt_\ell(\theta,j)=V_\ell(\theta,j)^\natural/e_\OO$, the characters
$\theta^\flat$, the integers $R_1(h,\ell,\theta)$ and $c_2(\theta)$, the cases and the counting statements of
\cref{pf:vwin-he:coef,pf:vwin-he:V,pf:vwin-he:host,pf:vwin-he:cases,pf:vwin-he:countdef}. Note that
$\vwinX_h(\OO,j)$ depends only on the parity of $j$. Besides the cases I, II, III, V of \cref{pf:vwin-he:cases} and the
case Z we use case IV: $\ell\in B_\chi$, $j\ge2$ and $\ord\theta^\flat$ odd.

\subsubsection{Statements}

\begin{definition}[Layer bounds, partial plus bounds]\label{pf:vwin:plusdefs}
Let $\chi$ be primitive, $\ell$ a prime, $j,t\in\N$ and $B\in\R$.
\begin{enumerate}[label=(\alph*)]
\item The \emph{layer bound} $\Lambda(\chi,\ell,j,t,B)$ is the statement: for every field $L_0$ of type
$M^{\mathrm{cyc}}_{\ell N_\chi}$ and every extension $L_1/L_0$ generated by $\ell$-power roots of unity with $[L_1:L_0]=\ell^j$,
$\sfname{VVat}(\ell,L_1,t,B)$ holds.
\item For a CM field $E$, an algebraic closure $\Omega$, $p$-data $D_p$, a finite set $S$ of primes and $B\in\R$, the
\emph{partial plus bound} $\Pi(S,t,B)$ means: $\vwinplus(D_\ell,j)<\infty$ for $\ell\in S$, $0<j<t$, and
$\sum_{\ell\in S}\log\ell\sum_{j=1}^{t-1}\vwinplus(D_\ell,j)^\natural\le B$.
\item $\mathrm{Count}_{\mathrm{IV}}$ is the statement: for every primitive $\chi$, every $\ell\in B_\chi$ and every $j$
there is $C$ such that for every host $h$ with character $\chi$, every $t\ge2$ for which $h$ is big, every $B$ with
$\Lambda(\chi,\ell,j,t,B)$, every $\Omega$, every coefficient ring $\OO$ for $\ell$ and every character $\theta$ modulo
$N_h$ with values in $\OO$, trivial on $H_h$ and with $\ord\theta^\flat$ odd: $V_\ell(\theta,j')<\infty$ for $0<j'<t$ and
$\sum_{j'=1}^{t-1}\vwinVt_\ell(\theta,j')\le B+(N_ht)^C$.
\item $\sfname{CountV}'=\mathrm{Count}(P)$ for $P$: ``$(\ell,j,\theta)$ is in case V and $2^{c_2(\theta)}<R_1(h,\ell,\theta)$''.
\end{enumerate}
\formallabel{UpperTheory.V.HV.LayerLe}
\formallabel{UpperTheory.V.HV.CountIV}
\formallabel{UpperTheory.V.HV.CountVs}
\end{definition}

\begin{definition}[Bounds at bad and good primes]\label{pf:vwin:sumdefs}
\begin{enumerate}[label=(\alph*)]
\item $\Sigma_{\mathrm{bad}}$: for every primitive $\chi$, every $\ell\in B_\chi$ and every $j$ there is $C$ such that for every host
$h$ with character $\chi$, every $t\ge2$ for which $h$ is big, every $B\ge0$ with $\Lambda(\chi,\ell,j,t,B)$, every CM field
$E=E_{\chi,m,q_1,d}$ (where $h=(\chi,m,q_1,d)$) and every choice of $\Omega$ and $p$-data, $\Pi(\{\ell\},t,(N_ht)^C+CB)$ holds.
\item $\Sigma_{\mathrm{good}}$: for every primitive $\chi$ there are $C$ and primes $r_1,r_2$ such that for every host
$h=(\chi,m,q_1,d)$ with $r_1,r_2<\min(q_1,d)$, every $t\ge2$ for which $h$ is big, every CM field $E=E_{\chi,m,q_1,d}$, every
choice of $\Omega$ and $p$-data and every finite set $S$ of odd primes not in $B_\chi$, $\Pi(S,t,(N_ht)^C)$ holds.
\item $\mathrm{P}_0$: for every primitive $\chi$ and every function $j_\bullet$ on primes there are $C$ and primes $r_1,r_2$ such
that for every host $h=(\chi,m,q_1,d)$ with $r_1,r_2<\min(q_1,d)$, every $t\ge2$ for which $h$ is big and every $B\ge0$ such
that $\Lambda(\chi,\ell,j_\ell,t,B)$ holds for all $\ell\in B_\chi$, every CM field $E=E_{\chi,m,q_1,d}$ satisfies
$\sfname{VVplusLe}(E,t,(N_ht)^C+CB)$.
\end{enumerate}
\formallabel{UpperTheory.V.HV.SumIn}
\formallabel{UpperTheory.V.HV.SumOut}
\formallabel{UpperTheory.V.HV.PointHV}
\end{definition}

\begin{theorem}[\cref{thm:plus-part}]\label{pf:vwin:hv}
Assume Poitou--Tate duality, the local finiteness and local Tate duality statements over finite extensions of $\Ql$,
$\sfname{LemmaK}(C_K)$, and $\sfname{CountI}$, $\sfname{CountII}$, $\sfname{CountIII}$, $\sfname{CountV}$,
$\sfname{CountZ}$. Then $\mathrm{P}(K,C_t)$ holds for every constant package $K$ and every $C_t>0$.
\formallabel{UpperTheory.V.HV.thmHV_of_counts}
\formallabel{UpperTheory.V.HV.pointHV_of_counts}
\end{theorem}

The proof occupies the rest of this subsection: \cref{pf:vwin:sumout,pf:vwin:sumin} prove $\Sigma_{\mathrm{good}}$ and
$\Sigma_{\mathrm{bad}}$, \cref{pf:vwin:point} combines them into $\mathrm{P}_0$, and \cref{pf:vwin:thmhv} passes from
$\mathrm{P}_0$ to $\mathrm{P}(K,C_t)$.

\subsubsection{Inputs}

We use the following facts about hosts, which are shared lemmas of the development. Let $h=(\chi,m,q_1,d)$ be a host which is
big for $t$.
\begin{enumerate}[label=(H\arabic*)]
\item\label{pf:vwin:H-bchi} (\lref{UpperTheory.V.HV.bchiFacts}) Every $\ell\in B_\chi$ is smaller than $q_1$ and $d$; and for every odd
prime $\ell\notin B_\chi$: $\ell\ge5$, $\ell$ divides none of $o_\chi$, $N_\chi$, $h^*(\chi)$, $[(\Z/N_h)^\times:H_h]$, $\ell\ne p_K(\chi),p_1(\chi)$,
and for $0<j<t$ with $(\ell-1)\mid2j$ one has $j\ge2$, $\ell\le2t$ and $\ell\nmid N_h$.
\item\label{pf:vwin:H-cases} (\lref{UpperTheory.V.HV.cases}) Let $\theta$ be a character modulo $N_h$ trivial on $H_h$, $\ell$ an
odd prime and $0<j<t$. If $\ell\notin B_\chi$, then $(\ell,j,\theta)$ is in case I, Z, II or III; if $\ell\in B_\chi$, it is in case
II, IV or V.
\item\label{pf:vwin:H-odd} (\lref{UpperTheory.V.HV.hostOdd}) If moreover $\ell\ne q_1,d$ and $\ord\theta^\flat$ is odd, then
$\cond\theta$ divides $N_\chi$.
\item\label{pf:vwin:H-index} (\lref{UpperTheory.V.HV.hostIndex}) $16\mid N_h$, $q_1d\mid N_h$ and $[(\Z/N_h)^\times:H_h]\le N_h$.
\end{enumerate}
In addition we use the following lemmas of this part of the development.
\begin{enumerate}[label=(R\arabic*)]
\item\label{pf:vwin:R-card} (\lref{UpperTheory.V.HV.xCard}) For $N\ge1$ and a domain $\OO$ there are at most $N$ Dirichlet
characters modulo $N$ with values in $\OO$, and at most $M$ of them have conductor dividing a given $M\ge1$.
\item\label{pf:vwin:R-level} (\lref{UpperTheory.V.HV.hostLevel}, \lref{UpperTheory.V.HV.hostCond}) Every
$E=E_{\chi,m,q_1,d}$ is the fixed field of $H_h$ under an embedding into $\Q(\zeta_{N_h})$ (``$E$ has level $(N_h,H_h)$''), and
$\cond(E)=N_h$.
\item\label{pf:vwin:R-red} (Reduction to characters; \lref{UpperTheory.V.HV.redPrime},
\lref{UpperTheory.V.HV.redAll}.) Let $E$ be a CM field of level $(N,H)$, $\ell$ an odd prime, $\OO$ a coefficient ring
containing the $\varphi(N)$-th roots of unity, $D$ an $\ell$-datum of $E$ and $j\ne0$. Then
\[
 e_\OO\cdot\vwinplus(D,j)\ \le\ \sum_{\theta}\Bigl(V_\ell(\theta,j)+e_\OO\,\lambda(\ell,N,[(\Z/N)^\times:H],j)\Bigr),
\]
the sum over the characters $\theta$ modulo $N$ trivial on $H$ with $\theta(-1)=(-1)^j$, where $\lambda$ is an explicit loss
term with $\lambda(\ell,N,n,j)\le(Nj)^6$ whenever $2\le N$, $n\le N$ and $j\ge1$ (\lref{UpperTheory.V.HV.lossLe}); if
$\ell\nmid[(\Z/N)^\times:H]$ the loss term can be omitted.
\item\label{pf:vwin:R-sum} (\lref{UpperTheory.V.Chr.sumBack}) If for $0<j<t$ finitely many finite lengths $f_i(j)$ ($i\in T_j\subset
T_0$) and a length $a(j)$ satisfy $e_\OO a(j)\le\sum_{i\in T_j}(f_i(j)+b(j)e_\OO)$, then $a(j)<\infty$ and
$\sum_{j=1}^{t-1}a(j)^\natural\le\sum_{j=1}^{t-1}\sum_{i\in T_j}f_i(j)^\natural/e_\OO+|T_0|\sum_{j=1}^{t-1}b(j)$.
\end{enumerate}
The reduction \ref{pf:vwin:R-red} rests on Shapiro's lemma and inflation--restriction: the plus part of $\Sha_D(j)$ is
identified with a complex-conjugation eigenspace of a Tate--Shafarevich group over $\Q$ with coefficients in the
induced module (\lref{UpperTheory.V.Chr.partModel}), which decomposes into the eigenspaces of the characters of
$(\Z/N)^\times/H$ when $\ell$ does not divide the index, the eigenspace of $\theta$ being the group defining $V_\ell(\theta,j)$; in
general the loss term bounds the failure of this decomposition. Finally we use $\sfname{CountZ}$ (\cref{pf:vwin:propZ}) and the following two statements.
\begin{enumerate}[label=(C\arabic*)]
\item\label{pf:vwin:C-IV} $\mathrm{Count}_{\mathrm{IV}}$ holds; this is proved in \cref{pf:vwin:countIV} below.
\item\label{pf:vwin:C-V} Poitou--Tate duality, the local finiteness and duality statements and $\sfname{LemmaK}(C_K)$ imply
$\sfname{CountV}'$; this is proved in \cref{pf:vwin:countVs} below.
\end{enumerate}

\begin{proposition}\label{pf:vwin:countIV}
$\mathrm{Count}_{\mathrm{IV}}$ holds; for $(\chi,\ell,j)$ one can take $C=\ell N_\chi\ell^j+3$.
\formallabel{UpperTheory.V.HV.countIV}
\end{proposition}

\begin{proof}
Let $\ell\in B_\chi$ (an odd prime), $h$ a host with character $\chi$ which is big for $t\ge2$, $B$ such that $\Lambda(\chi,\ell,j,t,B)$ holds, and
$\theta$ trivial on $H_h$ with $\ord\theta^\flat$ odd. By \ref{pf:vwin:H-bchi}, $\ell\ne q_1,d$, so $\cond\theta\mid N_\chi$ by \ref{pf:vwin:H-odd}.
There are a field $L_0$ of type $M^{\mathrm{cyc}}_{\ell N_\chi}$ and a layer $L_1/L_0$ of degree $\ell^j$, abelian over $\Q$ and embedded in $\Omega$,
with $[L_1:\Q]\le\ell N_\chi\ell^j$, such that every character of conductor dividing $N_\chi$ is trivial on $\Gal(\Omega/L_1)$
(\lref{UpperTheory.V.HV.ivField}; $L_0$ is the fixed field of the trivial subgroup of $(\Z/\ell N_\chi)^\times$). Let $D$ be an
$\ell$-datum of $L_1$. By the layer bound, $\vwinsh(D,e)<\infty$ for $0<e<t$ and $\sum_{e<t}\vwinsh(D,e)^\natural\le B$. For $0<e<t$,
$V_\ell(\theta,e)$ is at most the length of $\Sha_\ell(\Q,W_{\theta^{-1}\vwincyc^{-e}})$ (the conditions at the primes dividing $N_h$ only shrink
the group; \lref{UpperTheory.V.Chr.V'_le_V}), and restriction to $\Gal(\Omega/L_1)$, on which $\theta$ is trivial, gives
\[
 V_\ell(\theta,e)\le e_\OO\,\vwinsh(D,e)+[L_1:\Q]\,\nu(\ell,N_h,e)\,e_\OO ,
\]
by a comparison of $\Sha$ over $\Q$ with $\Sha$ over the abelian extension $L_1$ (\lref{UpperTheory.V.Chr.upV}), whose error
term is $[L_1:\Q]$ times a bound $\nu(\ell,N_h,e)$ for the lengths of certain local groups $H^0$
(\lref{UpperTheory.V.Chr.H0_local_card_bounded_of_cls_ker_le_ker}); moreover $\nu(\ell,N,e)\le2Ne$ (\lref{UpperTheory.V.HV.lossLe}). In particular $V_\ell(\theta,e)<\infty$, and dividing by $e_\OO$ and summing over $0<e<t$
(\lref{UpperTheory.V.Chr.nuSum}) gives $\sum_{e<t}\vwinVt_\ell(\theta,e)\le B+\sum_{e<t}[L_1:\Q]\cdot2N_he\le B+(N_ht)^{\ell N_\chi\ell^j+3}$.
\end{proof}

\begin{proposition}\label{pf:vwin:countVs}
Poitou--Tate duality, the local finiteness and duality statements and $\sfname{LemmaK}(C_K)$ imply $\sfname{CountV}'$.
\formallabel{UpperTheory.V.HV.countVs}
\end{proposition}

\begin{proof}
Fix $\chi$; we may replace $C_K$ by $\lceil C_K\rceil$. Let $M_\chi=\max B_\chi$ and
$n_0=2o_\chi\cdot16o_\chi(M_\chi N_\chi)^2M_\chi^{2\varphi(4o_\chi)}\cdot M_\chi$, and take $C=(2n_0)^{C_K}(n_0+1)+n_0+3$. Let $h$ be big for $t>1$, $\ell$ odd, and
$\theta\in\vwinX_h(\OO,1)$ such that $(\ell,1,\theta)$ is in case V and $2^{c_2(\theta)}<R_1(h,\ell,\theta)$; so $\ell\in B_\chi$ and $\ell\le M_\chi$. The condition on
$c_2$ and $R_1$ bounds the order of $\theta$ by $2o_\chi\cdot16o_\chi(\ell N_\chi)^2\ell^{2\varphi(4o_\chi)}$
(\lref{UpperTheory.V.HV.vsOrd}), hence $\ord\theta\cdot\ell\le n_0$. There is a field $L\subset\Omega$, abelian over $\Q$, containing a
primitive $\ell$-th root of unity, of degree at most $\ell\ord\theta\le n_0$, embedding into $\Q(\zeta_{\ell N_h})$, with $\theta$ trivial on
$\Gal(\Omega/L)$ (\lref{UpperTheory.V.HV.vsField}). Let $D$ be an $\ell$-datum of $L$. As in the proof of \cref{pf:vwin:countIV},
$V_\ell(\theta,1)\le e_\OO\vwinsh(D,1)+[L:\Q]\,\nu(\ell,N_h,1)\,e_\OO$; by Poitou--Tate duality $\vwinsh(D,1)\le\vwinVL(D,1)$ ($L$ is totally
imaginary; \lref{UpperTheory.V.PropA.shaLeV}); and by Lemma K, $\vwinVL(D,1)=\log\#H^2/\log\ell\le2\kappa(C_K,1,[L:\Q],|d_L|)$. Since $L$
embeds into $\Q(\zeta_{\ell N_h})$, $\log|d_L|\le[L:\Q]\log(\ell N_h)$, so $\kappa(C_K,1,[L:\Q],|d_L|)\le(\ell N_h)^{(2n_0)^{C_K}(n_0+1)}$. With
$\nu(\ell,N_h,1)\le2N_h$ this gives $V_\ell(\theta,1)<\infty$ and $\vwinVt_\ell(\theta,1)\le2(\ell N_h)^{(2n_0)^{C_K}(n_0+1)}+2n_0N_h$, and, as
$\log\ell\le\ell$, $\vwinVt_\ell(\theta,1)\log\ell\le(\ell N_h)^C$.
\end{proof}
\subsubsection{Good primes}

\begin{proposition}\label{pf:vwin:sumout}
$\sfname{CountI}$, $\sfname{CountZ}$, $\sfname{CountII}$ and $\sfname{CountIII}$ imply $\Sigma_{\mathrm{good}}$.
\formallabel{UpperTheory.V.HV.sumOut}
\end{proposition}

\begin{proof}
Fix $\chi$. Let $C_1,r_1,r_2$ be the constants of $\sfname{CountI}$ and $C_Z,C_2,C_3$ those of $\sfname{CountZ}$,
$\sfname{CountII}$, $\sfname{CountIII}$; put $K=\max(C_Z,C_2,C_3)$ and $C=C_1+3K+4$. Let $h$, $t$, $E$, $\Omega$, $(D_p)$ and $S$
be as in $\Sigma_{\mathrm{good}}$, and put $M=N_ht\ge2$. Fix an algebraic closure $\Omega_0$ of $\Q$, and for each prime $\ell$ a
coefficient ring $\OO_\ell$ containing the $\varphi(N_h)$-th roots of unity. Let $\ell\in S$, $0<j<t$ and
$\theta\in\vwinX_h(\OO_\ell,j)$.

\emph{Step 1.} If $(\ell,j)$ is not in case I, then $(\ell-1)\mid2j$, so $(\ell,j,\theta)$ is in case Z, II or III by
\ref{pf:vwin:H-cases}, and $\ell\le2t$ by \ref{pf:vwin:H-bchi}. Hence $(\theta,j)$ is bounded by $(\ell N_hj)^{C'}$ with
$C'\le K$, and $1\le\ell N_hj\le2t\cdot N_h\cdot t\le M^3$; so $V_\ell(\theta,j)<\infty$ and
$\vwinVt_\ell(\theta,j)\log\ell\le M^{3K}$.

\emph{Step 2.} If $(\ell,j)$ is in case I, $\sfname{CountI}$ applied to $L=\{\ell\}$ gives $V_\ell(\theta,j)<\infty$. Thus all
$V_\ell(\theta,j)$ are finite.

\emph{Step 3.} Since $\ell\nmid[(\Z/N_h)^\times:H_h]$ (\ref{pf:vwin:H-bchi}) and $E$ has level $(N_h,H_h)$, \ref{pf:vwin:R-red}
gives $e_{\OO_\ell}\vwinplus(D_\ell,j)\le\sum_{\theta\in\vwinX_h(\OO_\ell,j)}V_\ell(\theta,j)$, and \ref{pf:vwin:R-sum} (with
$b=0$) gives $\vwinplus(D_\ell,j)<\infty$ and
\[
\textstyle\sum_{j=1}^{t-1}\vwinplus(D_\ell,j)^\natural\le\sum_{j=1}^{t-1}\sum_{\theta\in\vwinX_h(\OO_\ell,j)}\vwinVt_\ell(\theta,j).
\]

\emph{Step 4.} Fix $0<j<t$ and split $S=S_{\mathrm I}\sqcup S'$, where $S_{\mathrm I}$ consists of the $\ell$ with $(\ell,j)$ in case I.
By $\sfname{CountI}$, $\sum_{\ell\in S_{\mathrm I}}\log\ell\sum_\theta\vwinVt_\ell(\theta,j)\le M^{C_1}$. Every $\ell\in S'$ is at most
$2t$, so $|S'|\le2t$; and for $\ell\in S'$, Step 1 and \ref{pf:vwin:R-card} give
$\log\ell\sum_\theta\vwinVt_\ell(\theta,j)\le N_hM^{3K}$. Hence
$\sum_{\ell\in S}\log\ell\sum_\theta\vwinVt_\ell(\theta,j)\le M^{C_1}+2t\,N_hM^{3K}=M^{C_1}+2M\cdot M^{3K}$.

\emph{Step 5.} By Steps 3 and 4,
$\sum_{\ell\in S}\log\ell\sum_{j=1}^{t-1}\vwinplus(D_\ell,j)^\natural\le t\,(M^{C_1}+2M\cdot M^{3K})\le M^{C_1+3K+4}$, using
$t\le M$ and $M\ge2$.
\end{proof}

\subsubsection{Bad primes}

\begin{proposition}\label{pf:vwin:sumin}
$\sfname{CountII}$, $\sfname{CountV}$, $\sfname{CountV}'$ and $\mathrm{Count}_{\mathrm{IV}}$ imply $\Sigma_{\mathrm{bad}}$.
\formallabel{UpperTheory.V.HV.sumIn}
\end{proposition}

\begin{proof}
Fix $\chi$, $\ell\in B_\chi$ and $j$. Let $C_2,C_5,C_6$ be the constants of $\sfname{CountII}$, $\sfname{CountV}$,
$\sfname{CountV}'$ for $\chi$, $C_4$ that of $\mathrm{Count}_{\mathrm{IV}}$ for $(\chi,\ell,j)$, $C_m=\max(C_2,C_5,C_6)$, and
$C=C_4+2C_m+9+\ell N_\chi$. Let $h,t,B,E,\Omega,(D_p)$ be as in $\Sigma_{\mathrm{bad}}$, put $M=N_ht$ and choose a coefficient ring
$\OO$ for $\ell$ containing the $\varphi(N_h)$-th roots of unity. By \ref{pf:vwin:H-bchi} and \ref{pf:vwin:H-index},
$\ell<q_1,d$, $N_\chi\le N_h$, $16\le N_h$, $\ell\le N_h$ and $M\ge3$.

Let $0<j'<t$ and $\theta\in\vwinX_h(\OO,j')$. By \ref{pf:vwin:H-cases}, $(\ell,j',\theta)$ is in case II, IV or V. In case II
(by $\sfname{CountII}$) and in case V (by $\sfname{CountV}$ or $\sfname{CountV}'$, according as $R_1(h,\ell,\theta)\le2^{c_2(\theta)}$
or not), $V_\ell(\theta,j')<\infty$ and $\vwinVt_\ell(\theta,j')\log\ell\le(\ell N_hj')^{C'}\le(M^2)^{C_m}$, since $\ell N_hj'\le M^2$.
In case IV, $\theta$ belongs to the set $\mathcal O$ of characters $\theta$ trivial on $H_h$ with $\ord\theta^\flat$ odd and
$\theta\in\vwinX_h(\OO,j'')$ for some $0<j''<t$; by $\mathrm{Count}_{\mathrm{IV}}$ (with the layer bound
$\Lambda(\chi,\ell,j,t,B)$), $V_\ell(\theta,j')<\infty$, and every $\theta\in\mathcal O$ satisfies
$\sum_{j'=1}^{t-1}\vwinVt_\ell(\theta,j')\le B+M^{C_4}$. By \ref{pf:vwin:H-odd} (note $\ell\ne q_1,d$) the characters in
$\mathcal O$ have conductor dividing $N_\chi$, so $|\mathcal O|\le N_\chi$ by \ref{pf:vwin:R-card}.

By \ref{pf:vwin:R-red}, for $0<j'<t$,
$e_\OO\vwinplus(D_\ell,j')\le\sum_{\theta\in\vwinX_h(\OO,j')}(V_\ell(\theta,j')+e_\OO\lambda_{j'})$ with
$\lambda_{j'}=\lambda(\ell,N_h,[(\Z/N_h)^\times:H_h],j')\le(N_hj')^6\le M^6$. Since all $\vwinX_h(\OO,j')$ lie in
$\vwinX_h(\OO,0)\cup\vwinX_h(\OO,1)$, which has at most $N_h$ elements, \ref{pf:vwin:R-sum} gives the finiteness of
$\vwinplus(D_\ell,j')$ and
\[
\textstyle\sum_{j'=1}^{t-1}\vwinplus(D_\ell,j')^\natural\le\sum_{j'=1}^{t-1}\sum_{\theta}\vwinVt_\ell(\theta,j')+N_h\,t\,M^6 .
\]
Multiply by $\log\ell$ and split the double sum according as $\theta\in\mathcal O$ or not: the first part is at most
$\log\ell\cdot N_\chi(B+M^{C_4})$, the second at most $t\cdot N_h(M^2)^{C_m}$ (at most $N_h$ characters for each $j'$). Using
$\log\ell\le\ell\le M$, $N_\chi\le M$, $t,N_h\le M$, $M\ge3$ and $\ell N_\chi\le C$, the total is at most $M^C+CB$; this is an
elementary verification.
\end{proof}

\subsubsection{Assembly}

\begin{proposition}\label{pf:vwin:point}
$\Sigma_{\mathrm{bad}}$ and $\Sigma_{\mathrm{good}}$ imply $\mathrm{P}_0$.
\formallabel{UpperTheory.V.HV.point}
\end{proposition}

\begin{proof}
Fix $\chi$ and $j_\bullet$. Let $C_0,r_1,r_2$ be the constants of $\Sigma_{\mathrm{good}}$, let $C_\ell$ ($\ell\in B_\chi$) be those
of $\Sigma_{\mathrm{bad}}$ for $(\chi,\ell,j_\ell)$, $D=\sum_{\ell\in B_\chi}C_\ell$, $b=|B_\chi|$ and $C=C_0+D+b+bD$. Let $h,t,B,E$ be as
in $\mathrm{P}_0$, fix $\Omega$ and $p$-data, and put $M=N_ht\ge2$. Finiteness of $\vwinplus(D_p,j)$ for odd $p$ follows from
$\Sigma_{\mathrm{bad}}$ applied to $\{p\}$ if $p\in B_\chi$ and from $\Sigma_{\mathrm{good}}$ applied to $\{p\}$ otherwise. For a finite
set $S$ of odd primes, split $S$ into $S\cap B_\chi$ and $S\smallsetminus B_\chi$: the second part contributes at most $M^{C_0}$,
and each $\ell$ in the first part at most $M^{C_\ell}+C_\ell B\le M^D+DB$. Hence the sum is at most
$M^{C_0}+b(M^D+DB)\le M^C+CB$.
\end{proof}

\begin{proposition}\label{pf:vwin:thmhv}
$\mathrm{P}_0$ implies $\mathrm{P}(K,C_t)$ for every constant package $K$ and every $C_t>0$.
\formallabel{UpperTheory.V.HV.thmHV}
\end{proposition}

\begin{proof}
Fix $\psi$ and apply $\mathrm{P}_0$ to $\chi=\psi$ and $j_\ell=\log_\ell R_\ell(L_\psi(\ell))+3$; this gives $C,r_1,r_2$. Let
$\eps\in(0,\frac1{52})$ and $A_1$ be given, and put $X_0=2^{J_\psi}$ ($J_\psi>0$). Let $(m,t)$, $X\ge X_0$, $(q,d)$, a $K$-window
$Y$ and $x\in[Y,Y^2]$ be as in $\mathrm{P}(K,C_t)$. Then $x\ge X^{1/J_\psi}\ge2$ and $x\le Y^2\le X$, so $m(x)\ge2$, $t(x)\ge2$ and
$t(x)\le x^{C_t}\le X^{C_t}$. Hence $h=(\psi,m(x),q,d)$ is a host, and it is big for $t(x)$: $2t(x)+1<q,d$ because
$q,d>2X^{C_t}+1$, and the other conditions are hypotheses of $\mathrm{P}(K,C_t)$; moreover $r_1,r_2<q,d$. With
$B=e^{x^{4\eps/5}}\ge0$, the layer bounds $\Lambda(\psi,\ell,j_\ell,t(x),B)$ for $\ell\in B_\psi$ are exactly the statement
$\mathrm{B}_b$ at $x$, which holds because $Y$ is a $K$-window. So $\mathrm{P}_0$ gives, for every CM field
$E=E_{\psi,m(x),q,d}$, $\sfname{VVplusLe}(E,t(x),(N_ht(x))^C+CB)$, and $N_h=\cond(E)$ by \ref{pf:vwin:R-level}. This is
\eqref{pf:vwin:hvat}.
\end{proof}

\begin{proof}[Proof of \cref{pf:vwin:hv}]
By \ref{pf:vwin:C-IV} and \ref{pf:vwin:C-V}, $\mathrm{Count}_{\mathrm{IV}}$ and $\sfname{CountV}'$ hold. \Cref{pf:vwin:sumin,pf:vwin:sumout} give
$\Sigma_{\mathrm{bad}}$ and $\Sigma_{\mathrm{good}}$, \cref{pf:vwin:point} gives $\mathrm{P}_0$, and \cref{pf:vwin:thmhv} gives
$\mathrm{P}(K,C_t)$.
\end{proof}

\begin{remark}\label{pf:vwin:rem-plus}
The hypothesis $\sfname{CountV}$ concerns only the sub-case $R_1\le2^{c_2}$ of case V; the complementary sub-case is
covered by \ref{pf:vwin:C-V}, given Lemma K. Case IV needs no counting hypothesis: it is controlled by the layer bounds, i.e.\
by \cref{thm:window:bab}(b).
\end{remark}


% ---------------------------------------------------------------------------
\subsection{Iwasawa towers, the window, and the two-part}\label{sec:pf:vwin:towers}

In this subsection we prove \cref{lem:window:tower} and \cref{thm:window:bab}. For a prime $\ell$ put $q_\ell=\ell$ if $\ell$
is odd and $q_\ell=4$ if $\ell=2$. For $R\in\N$ let $j_0(\ell,R)=\log_\ell R+3$. For number fields $L_0\subset L_1$, $L_1/L_0$
is a \emph{layer of degree $\ell^m$} if $L_1$ is generated over $L_0$ by $\ell^k$-th roots of unity for some $k$ and
$[L_1:L_0]=\ell^m$.

\subsubsection{Tower bounds}

\begin{definition}[Tower bounds]\label{pf:vwin:abound}
Let $\ell$ be a prime, $R\in\N$, $j_0=j_0(\ell,R)$, and $L_0$ a field.
\begin{enumerate}[label=(\alph*)]
\item The \emph{second bound} $\mathrm{S}(\ell,R,L_0)$: for every $j\ge1$, every layer $L_1/L_0$ of degree $\ell^m$, every
layer $L_2/L_0$ of degree $\ell^{j_0}$, every algebraic closure of $L_2$ and every $\ell$-datum $D$ of $L_2$,
$\vwinsh(L_1,\ell,j)\le\vwinVL(D,j)+R\cdot\max(m-j_0,0)$.
\item $\sfname{ABound}(\ell,R,L_0)$: there are a normed field $K$ which is a normed $\Ql$-algebra with ultrametric norm,
integers $c\ge0$ and $R'\le R$ with $\mu_{\ell^c}\subset L_0$, and points $s_1,\dots,s_{R'}\in K$ such that for every $j\ge1$,
every layer $L_1/L_0$ of degree $\ell^m$ and every $B\in\N$ with $\|s_i-j\|\ge\ell^{-B}$ for all $i$,
\[
 \vwinsh(L_1,\ell,j)\le R'\bigl(\max(m,j_0)+c+B\bigr);
\]
and $\mathrm{S}(\ell,R,L_0)$ holds.
\item The \emph{first bound} $\mathrm{F}(\ell,R,c,L_0,Z,\vwinexc)$, for $Z\in\N$ and a finite set $\vwinexc\subset\N$: $|\vwinexc|\le R$, and
for every $1\le j\le Z$ with $j\notin\vwinexc$ and every layer $L_1/L_0$ of degree $\ell^m$,
$\vwinsh(L_1,\ell,j)\le R(\max(m,j_0)+c+\log_\ell Z)$.
\item For $\eta,C_t,x\in\R$ and a finite set $\vwinexc\subset\N$, $x$ is \emph{$(\eta,C_t)$-good for $\vwinexc$} if every
$e\in\vwinexc$ with $e\le x^{C_t}$ satisfies $e\le x^\eta$.
\item $\mathrm{Tow}(\ell,L,R)$: $L$ is a totally imaginary number field containing a primitive $q_\ell$-th root of unity,
$R\ge1$, and the $\Zl$-module $\mathfrak X_\ell(G)$ of \cref{pf:vwin-he:dinf} (for $L_0=L$ and an algebraic closure of $L$) is
generated by $R$ elements.
\end{enumerate}
\formallabel{UpperTheory.V.Win.SecondBound}
\formallabel{UpperTheory.V.Win.ABound}
\formallabel{UpperTheory.V.Win.FirstBound}
\formallabel{UpperTheory.V.Win.Good}
\formallabel{UpperTheory.V.Tower.IsTower}
\end{definition}

This is the notion of \cref{def:window:abound}; by \cref{pf:vwin-he:dinf}, $\mathfrak X_\ell(G)$ is the module $X'_\infty$ of
\cref{lit:window:fw}, so $\mathrm{Tow}(\ell,L,R)$ says that $X'_\infty$ is generated by $R$ elements.

\begin{theorem}[\cref{lem:window:tower}]\label{pf:vwin:tower}
Assume Poitou--Tate duality and the local finiteness and duality statements over finite extensions of $\Ql$. If
$\mathrm{Tow}(\ell,L_0,R)$ holds, then $\sfname{ABound}(\ell,R,L_0)$ holds, with $K=\Ql$, $R'=R$, points $s_i\in\Zl$, and $c$ the
largest integer such that $\mu_{\ell^c}\subset L_0$.
\formallabel{UpperTheory.V.PropA.aBound_of_isTower}
\formallabel{UpperTheory.V.PropA.aBound_of_descends}
\end{theorem}

The proof is a combination of the descent formula of \cref{sec:pf:vwin-he}, which expresses $\vwinsh$ over a layer as the
length of a cokernel on $\mathfrak X_\ell(G)$, with two estimates for such cokernels, one linear in the layer index (growth
along the tower) and one in terms of the $\ell$-adic distance of the twist to finitely many ``centres''. The centres play the
role of the zeros, in the variable of the twist, of the characteristic power series of $X'_\infty$ in Iwasawa's theory of
$\Zl$-extensions \cite{Iwasawa1973}; no structure theory of Iwasawa modules is used, only that $\mathfrak X_\ell(G)$ is finitely
generated over $\Zl$. The descent formula requires an element of the Galois group whose cyclotomic image topologically generates the
image of the cyclotomic character; such elements exist by the following lemma.

\begin{lemma}[Topological generators]\label{pf:vwin:gengal}
Let $k\subset L$ be number fields, $\Omega$ an algebraic closure of $L$ (hence of $k$), $\ell$ a prime, $G_k=\Gal(\Omega/k)$ and
$G_L=\Gal(\Omega/L)\subset G_k$, and let $\vwincyc$ be the cyclotomic character of $G_k$.
\begin{enumerate}[label=(\alph*)]
\item\label{pf:vwin:gengal-a} $G_L$ is closed in $G_k$.
\item\label{pf:vwin:gengal-b} $L$ is procyclic at $\ell$ (in the sense of \cref{sec:pf:vwin-he}) if and only if $\ell$ is odd or
$\vwincyc(G_L)\subset1+4\Z_2$.
\item\label{pf:vwin:gengal-c} If $U\subset G_k$ is a closed subgroup and $\ell$ is odd or $\vwincyc(U)\subset1+4\Z_2$, there is $\gamma\in U$
such that $\vwincyc(U)$ is contained in the closure of the group generated by $\vwincyc(\gamma)$. In particular, if $L$ is procyclic
at $\ell$, such an element $\gamma$ exists in $U=G_L$.
\end{enumerate}
\formallabel{UpperTheory.V.Tower.genGal}
\formallabel{UpperTheory.V.Tower.existsGen}
\end{lemma}

\begin{proof}
\ref{pf:vwin:gengal-a} $G_L$ is the stabilizer of finitely many generators of $L$ over $k$, hence an open subgroup of $G_k$ for
the Krull topology; an open subgroup is closed, its complement being a union of cosets.
\ref{pf:vwin:gengal-b} The cyclotomic character of $G_L$ is the restriction of that of $G_k$, since both describe the action on the
$\ell$-power roots of unity in $\Omega$; so the condition defining procyclicity of $L$ is the stated one.
\ref{pf:vwin:gengal-c} The group $U$ is compact, being closed in the compact group $G_k$, and $\vwincyc$ is continuous, so
$V=\vwincyc(U)$ is a compact, hence closed, subgroup of $\Zl^\times$, contained in $1+4\Z_2$ if $\ell=2$. Such a subgroup contains
an element $u$ such that $V$ is contained in the closure of the group generated by $u$ (\lref{UpperTheory.V.Kum.existsGen}: for
$\ell$ odd, $\Zl^\times\cong\mu_{\ell-1}\times(1+\ell\Zl)$ with $1+\ell\Zl\cong\Zl$, and for $\ell=2$, $1+4\Z_2\cong\Z_2$, so the
closed subgroups in question are topologically cyclic). Write $u=\vwincyc(\gamma)$ with $\gamma\in U$. The last assertion follows
from \ref{pf:vwin:gengal-a} and \ref{pf:vwin:gengal-b}.
\end{proof}

\paragraph{Unipotent operators.} Let $X$ be a $\Zl$-module and $f$ an endomorphism. We write $f\in\ell^k$ if
$f(X)\subset\ell^kX$, and $\vwincok(f)=\len_{\Zl}X/f(X)\in\N\cup\{\infty\}$. We say that $X$ is \emph{generated by $R$
elements} if there is a surjection $\Zl^R\to X$, and that $g\in\End(X)$ is \emph{unipotent} if $(g-1)^N\in\ell$ for some $N$.

\begin{lemma}\label{pf:vwin:unip}
Let $X$ be generated by $R$ elements and let $g,h\in\End(X)$.
\begin{enumerate}[label=(\alph*)]
\item\label{pf:vwin:unip-a} If $g$ is unipotent, then $(g-1)^R\in\ell$.
\item\label{pf:vwin:unip-b} For $k\ge0$, $(g-1)^{\ell^k}\in\ell$ if and only if $g^{\ell^k}-1\in\ell$.
\item\label{pf:vwin:unip-c} If $h-1\in\ell^k$ then $h^b-1\in\ell^k$ for all $b\ge1$; if $h-1\in\ell$ then $h^\ell-1\in\ell^2$.
\item\label{pf:vwin:unip-d} If $h-1\in\ell^2$, then $\ell X\subset(1+h+\dots+h^{\ell-1})(X)$ and
$\vwincok(h^\ell-1)\le\vwincok(h-1)+R$.
\item\label{pf:vwin:unip-e} If $m\le m'$, then $\vwincok(g^{\ell^m}-1)\le\vwincok(g^{\ell^{m'}}-1)$.
\end{enumerate}
\formallabel{UpperTheory.V.PropA.unipPow}
\formallabel{UpperTheory.V.PropA.frob}
\formallabel{UpperTheory.V.PropA.powIn}
\formallabel{UpperTheory.V.PropA.phi}
\formallabel{UpperTheory.V.PropA.step}
\formallabel{UpperTheory.V.PropA.mono}
\end{lemma}

\begin{proof}
\ref{pf:vwin:unip-a} $g-1$ preserves $\ell X$ and induces a nilpotent endomorphism $F$ of $X/\ell X$, a module of length at most
$R$. A nilpotent endomorphism of a module of length $\le R$ satisfies $F^R=0$ (the chain $\ker F\subset\ker F^2\subset\cdots$
increases strictly until it becomes stationary at the whole module). Hence $(g-1)^R(X)\subset\ell X$.

\ref{pf:vwin:unip-b} In $\F_\ell[T]$ one has $(T-1)^{\ell^k}=T^{\ell^k}-1$, so $(T-1)^{\ell^k}-(T^{\ell^k}-1)=\ell Q(T)$ with
$Q\in\Z[T]$, and $(g-1)^{\ell^k}-(g^{\ell^k}-1)=\ell Q(g)\in\ell$. The assertion follows since $f\in\ell$ and $f'-f\in\ell$ imply
$f'\in\ell$.

\ref{pf:vwin:unip-c} $h^b-1=(h-1)(1+h+\dots+h^{b-1})$ and the second factor preserves $\ell^kX$. If $h-1\in\ell$, then
$1+h+\dots+h^{\ell-1}=\ell+\sum_{i<\ell}(h^i-1)\in\ell$, and the composite of two endomorphisms in $\ell$ lies in $\ell^2$.

\ref{pf:vwin:unip-d} Let $S=1+h+\dots+h^{\ell-1}$. By \ref{pf:vwin:unip-c}, $h^i x\in x+\ell^2X$ for all $i$, so $Sx\in\ell x+\ell^2X$.
Hence $\ell X\subset S(X)+\ell(\ell X)$, so the image of $\ell X$ in $X/S(X)$ equals $\ell$ times itself; being a finitely
generated $\Zl$-module, it vanishes by Nakayama's lemma, i.e.\ $\ell X\subset S(X)$. Now $h^\ell-1=S(h-1)$, and the length of the
cokernel of a composite is at most the sum of the lengths of the cokernels, so
$\vwincok(h^\ell-1)\le\vwincok(S)+\vwincok(h-1)\le\len(X/\ell X)+\vwincok(h-1)\le R+\vwincok(h-1)$.

\ref{pf:vwin:unip-e} $g^{\ell^{m'}}-1=(g^{\ell^m}-1)\circ P(g)$ for a polynomial $P$, so the image of $g^{\ell^{m'}}-1$ is contained
in that of $g^{\ell^m}-1$.
\end{proof}

\begin{lemma}[Growth]\label{pf:vwin:growth}
Let $X$ be generated by $R$ elements and $g\in\End(X)$ unipotent, and put $j_0=j_0(\ell,R)$. For all $m\ge0$,
\[
 \vwincok(g^{\ell^m}-1)\le\vwincok(g^{\ell^{j_0}}-1)+R\cdot\max(m-j_0,0) .
\]
\formallabel{UpperTheory.V.PropA.growth}
\end{lemma}

\begin{proof}
Let $k=\log_\ell R+1$, so that $R<\ell^k$ and $j_0=k+2$. By \cref{pf:vwin:unip}\ref{pf:vwin:unip-a},
$(g-1)^{\ell^k}=(g-1)^{\ell^k-R}(g-1)^R\in\ell$; by \ref{pf:vwin:unip-b}, $g^{\ell^k}-1\in\ell$; by \ref{pf:vwin:unip-c},
$g^{\ell^{k+1}}-1\in\ell^2$ and then $g^{\ell^n}-1\in\ell^2$ for all $n\ge k+1$. By \ref{pf:vwin:unip-d} applied to
$h=g^{\ell^n}$, $\vwincok(g^{\ell^{n+1}}-1)\le\vwincok(g^{\ell^n}-1)+R$ for $n\ge k+1$, and induction gives
$\vwincok(g^{\ell^{j_0+d}}-1)\le\vwincok(g^{\ell^{j_0}}-1)+Rd$ for $d\ge0$. If $m\le j_0$ use \ref{pf:vwin:unip-e}.
\end{proof}

\paragraph{Centres.} Let $c\ge1$ ($c\ge2$ if $\ell=2$) and $u\in\Zl$ with $v_\ell(u-1)=c$. For an endomorphism $\gamma$ of $X$ and
$e,m\in\N$ put $\tau(\gamma,u,e,m)=\vwincok((u^e\gamma)^{\ell^m}-1)$.

\begin{lemma}\label{pf:vwin:centres}
Let $X$ be generated by $R$ elements, $\gamma\in\End(X)$ unipotent, and $u$, $c$ as above. There are $s_1,\dots,s_R\in\Zl$ such
that for all $e,m,B\in\N$ with $|s_i-e|_\ell\ge\ell^{-B}$ for all $i$,
\[
 \tau(\gamma,u,e,m)\le R\bigl(\max(m,j_0)+c+B\bigr);
\]
and for all $e,m\in\N$, $\tau(\gamma,u,e,m)\le\tau(\gamma,u,e,j_0)+R\max(m-j_0,0)$.
\formallabel{UpperTheory.V.PropA.exists_centres_norm}
\formallabel{UpperTheory.V.PropA.exists_centres}
\formallabel{UpperTheory.V.PropA.exists_centres_level}
\formallabel{UpperTheory.V.PropA.twLen_le_growth}
\formallabel{UpperTheory.V.PropA.detBound}
\formallabel{UpperTheory.V.PropA.ball}
\formallabel{UpperTheory.V.PropA.lift}
\end{lemma}

\begin{proof}
\emph{(i) The second inequality.} Since $\ell\mid u-1$, $\ell\mid u^e-1$, so $u^e\gamma$ is again unipotent
($(u^e\gamma-1)^N$ is a combination of $(\gamma-1)^i$ and multiples of $u^e-1$; \lref{UpperTheory.V.PropA.twist}), and
\cref{pf:vwin:growth} applies to $g=u^e\gamma$.

\emph{(ii) Lifting.} Choose a surjection $\pi:\Zl^R\to X$ and a matrix $S\in M_R(\Zl)$ with $\pi\circ S=\gamma^{\ell^j}\circ\pi$
(lift the images of the basis vectors). For $a\in\Zl$, $\pi\circ(aS-1)=(a\gamma^{\ell^j}-1)\circ\pi$, so
$X/(a\gamma^{\ell^j}-1)X$ is a quotient of $\Zl^R/(aS-1)\Zl^R$, whose length is $v_\ell(\det(aS-1))$ (elementary divisors; both
sides are infinite if the determinant vanishes). Hence $\vwincok(a\gamma^{\ell^j}-1)\le v_\ell(\det(aS-1))$.

\emph{(iii) Factorization.} Let $\Qlbar$ carry the extension of $|\cdot|_\ell$. There are $z_1,\dots,z_R\in\Qlbar$ (the
eigenvalues of $S$) with $\det(aS-1)=\prod_i(az_i-1)$ for all $a$.

\emph{(iv) One eigenvalue.} The map $e\mapsto u^{\ell^je}$ satisfies $|u^{\ell^ja}-u^{\ell^jb}|=\ell^{-(j+c)}|a-b|$ for $a,b\in\N$
(\lref{UpperTheory.V.PropA.lteNorm}; via the $\ell$-adic logarithm, $1+\ell^c\Zl\cong\ell^c\Zl$). For units $a,b$ of $\Zl$ and
$z\in\Qlbar$ one has $|a-b|\le\max(|az-1|,|bz-1|)$: if $|z|\ne1$ both terms on the right are $\ge1$, and if $|z|=1$,
$|a-b|=|az-bz|$. Hence the function $r(e)=\ell^{j+c}|u^{\ell^je}z-1|$ satisfies $|e-e'|\le\max(r(e),r(e'))$, i.e.\ the closed
balls of radii $r(e)$ around the integers $e$ in $\Zl$ pairwise intersect. In the ultrametric compact space $\Zl$ such a
family has a common point $s$ (\lref{UpperTheory.V.PropA.helly}): $|s-e|\le r(e)$, i.e.\
$v_\ell(u^{\ell^je}z-1)\le j+c+v_\ell(s-e)$ for all $e\in\N$.

\emph{(v) Level $j$.} Since $(u^e\gamma)^{\ell^j}=(u^{\ell^j})^e\gamma^{\ell^j}$, (ii)--(iv) give $s_1,\dots,s_R$ with
$\tau(\gamma,u,e,j)\le\sum_i\bigl(j+c+v_\ell(s_i-e)\bigr)$ for all $e$.

\emph{(vi) Conclusion.} Take $j=j_0$ in (v) and combine with (i): $\tau(\gamma,u,e,m)\le\sum_i(\max(m,j_0)+c+v_\ell(s_i-e))$.
If $|s_i-e|\ge\ell^{-B}$ then $v_\ell(s_i-e)\le B$.
\end{proof}

\begin{proof}[Proof of \cref{pf:vwin:tower}]
Let $G=\Gal(\overline{L_0}/L_0)$, $\mathfrak X=\mathfrak X_\ell(G)$, and let $c$ be the largest integer with $\mu_{\ell^c}\subset
L_0$; then $c\ge1$, and $c\ge2$ if $\ell=2$. By \cref{pf:vwin-he:layers}, $\vwincyc(G)=1+\ell^c\Zl$, and every $u\in\Zl$ with
$v_\ell(u-1)=c$ generates it topologically; hence there is $\gamma\in G$ such that $\vwincyc(G)$ is contained in the closure of
the group generated by $u=\vwincyc(\gamma)$ and $v_\ell(u-1)=c$ (\lref{UpperTheory.V.PropA.genLevel}). Let $\gamma_{\mathfrak X}$
be the action of $\gamma$ on $\mathfrak X$ (the action of \cref{pf:vwin-he:dinf} with $e=0$). Since $\mathfrak X$ is finitely generated
and $\gamma\in G(1)$, $\gamma_{\mathfrak X}$ is unipotent (\lref{UpperTheory.V.PropA.xUnip}). By the descent formula
\cref{pf:vwin-he:step0} (\lref{UpperTheory.V.Tower.step0}), for
$e\ge1$ and every layer $L_1/L_0$ of degree $\ell^m$,
\begin{equation}\label{pf:vwin:descent}
 \vwinsh(L_1,\ell,e)=\vwincok\bigl((u^e\gamma_{\mathfrak X})^{\ell^m}-1\bigr)=\tau(\gamma_{\mathfrak X},u,e,m),
\end{equation}
because $\gamma$ acts on $\mathfrak X$ through the twisted action by $u^e\gamma_{\mathfrak X}$
(\lref{UpperTheory.V.PropA.descends}). The first bound of $\sfname{ABound}$ now follows from \cref{pf:vwin:centres}, with
$K=\Ql$, $R'=R$ and the points $s_i\in\Zl$ (for $s\in\Zl$ the norm in $\Ql$ is $|s-e|_\ell$). For the second bound let $L_2/L_0$ be a
layer of degree $\ell^{j_0}$ and $D$ an $\ell$-datum of $L_2$. By \cref{pf:vwin:centres} and \eqref{pf:vwin:descent},
$\vwinsh(L_1,\ell,e)\le\vwinsh(L_2,\ell,e)+R\max(m-j_0,0)$, and $\vwinsh(L_2,\ell,e)=\vwinsh(D,e)\le\vwinVL(D,e)$ by Poitou--Tate
duality, since $L_2$ is totally imaginary (\lref{UpperTheory.V.PropA.shaLeV}).
\end{proof}

\subsubsection{Fixed fields and the two-part}

\begin{definition}\label{pf:vwin:twobase}
Let $M$ be an abelian datum and $(m_M;q,d)$ a prime pair. A number field $L_0$ is a \emph{$2$-base} for $(M,q,d)$ if $L_0$ contains
$\mu_4$, $[L_0:\Q]\le8\deg M$, $L_0$ embeds into $\Q(\zeta_{4m_Mqd})$, and for every $k\ge2$ and every $E=E_{M,k,q,d}$ there is a layer
$L_1/L_0$ of degree $2^{m'}$ with $m'\le k+2$ such that $\vwinsh(E,2,j)\le1+\vwinsh(L_1,2,j)$ for all $j\ge1$.
\formallabel{UpperTheory.V.Win.TwoBase}
\end{definition}

\begin{proposition}\label{pf:vwin:fixedtwo}
Assume the theorem of Ferrero--Washington in the form of \cref{lit:window:fw}, and that $\mathrm{Tow}(\ell,L,R)$ implies
$\sfname{ABound}(\ell,R,L)$ for all $\ell,L,R$.
\begin{enumerate}[label=(\alph*)]
\item For every odd prime $\ell$, every abelian datum $A$ with $\ell\mid m_A$ and $H_A\subset\{a\equiv1\bmod\ell\}$, and every field
$L_0$ of type $A$, there is $R$ with $\sfname{ABound}(\ell,R,L_0)$.
\item For every abelian datum $M$ there is $R$ such that for every prime pair $(m_M;q,d)$ there is a $2$-base $L_0$ for $(M,q,d)$
with $\sfname{ABound}(2,R,L_0)$.
\end{enumerate}
\formallabel{UpperTheory.V.Two.fixedIn}
\formallabel{UpperTheory.V.Two.twoIn}
\formallabel{UpperTheory.V.Win.FixedIn}
\formallabel{UpperTheory.V.Win.TwoIn}
\end{proposition}

\begin{proof}
(a) Fields of type $A$ contain $\mu_\ell$ and are totally imaginary. By Ferrero--Washington, the $\ell$-ranks of the class groups of
the fields $L_0(\mu_{\ell^n})$ are bounded; this bounds the number of generators of $\mathfrak X_\ell(G)$ (the Galois group of the
maximal unramified abelian pro-$\ell$ extension of $L_0(\mu_{\ell^\infty})$ in which the places above $\ell$ split completely) by some
$r$, uniformly for all fields of type $A$ (\lref{UpperTheory.V.Two.towerFixed}). Hence $\mathrm{Tow}(\ell,L_0,\max(1,r))$, and the
hypothesis gives $\sfname{ABound}$.

(b) Let $L_0$ be a field of type $M\vee M^{\mathrm{cyc}}_4\vee Q_q\vee Q_d$, the compositum of a field of type $M$ with $\Q(i)$ and
the quadratic subfields $\Q(\sqrt{-q})$, $\Q(\sqrt{-d})$ of $\Q(\zeta_q)$, $\Q(\zeta_d)$ (the abelian datum $Q_p=(p,\text{squares})$). By
\cref{pf:vwin:twobasethm} below, $L_0$ is a $2$-base and $\mathrm{Tow}(2,L_0,R)$ holds with $R=4(g+6\deg M)$, $g$ depending only on $M$; the
hypothesis gives $\sfname{ABound}(2,R,L_0)$.
\end{proof}

For an abelian datum $A$ let $U_A\subset G_\Q$ be the subgroup fixing the fields of type $A$ in $\Omega$, $U_{A,\infty}=U_A\cap\ker\vwincyc$
(for $\ell=2$), and for $N\ge1$ let $\mathcal C(A,N)$ be the $\F_2$-vector space of continuous homomorphisms $U_{A,\infty}\to\Z/2$ vanishing
on the inertia groups at the places not dividing $N$ and on the decomposition groups above $2$. Let $A^{(i)}=M\vee M^{\mathrm{cyc}}_4$ and
$A^{(q,d)}=A^{(i)}\vee Q_q\vee Q_d$.

\begin{proposition}[The $2$-base]\label{pf:vwin:twobasethm}
Let $M$ be an abelian datum.
\begin{enumerate}[label=(\alph*)]
\item\label{pf:vwin:twobasethm-a} If $\#\mathcal C(A^{(i)},1)\le2^g$, then $\#\mathcal C(A^{(q,d)},1)\le2^{4(g+6\deg M)}$ for every prime pair $(m_M;q,d)$.
\item\label{pf:vwin:twobasethm-b} Assume Ferrero--Washington. There is $R$ such that for every prime pair $(m_M;q,d)$ and every field $L$ of type
$A^{(q,d)}$, $\mathfrak X_2(G)$ (for $L_0=L$) is generated by $R$ elements and $\mathrm{Tow}(2,L,R)$ holds.
\item\label{pf:vwin:twobasethm-c} Every field of type $A^{(q,d)}$ is a $2$-base for $(M,q,d)$.
\end{enumerate}
\formallabel{UpperTheory.V.Two.baseRk}
\formallabel{UpperTheory.V.Two.towerTwo}
\formallabel{UpperTheory.V.Two.twoBase}
\end{proposition}

\begin{proof}
\ref{pf:vwin:twobasethm-a} Recall $q\equiv3\pmod8$, $d\equiv7\pmod{16}$, and $q,d$ are primes not dividing $2m_M$. Two steps enlarge
$\mathcal C$. \emph{Allowing ramification at $p\in\{q,d\}$:} the inertia groups at $p$ lie in $U_{A^{(i)},\infty}$
(\lref{UpperTheory.V.Two.ineTwo}), and $p$ has at most $c\deg M$ places in the fixed field of $U_{A^{(i)},\infty}$, with $c=2$ for $p=q$
and $c=4$ for $p=d$: the decomposition group of $p$ in the cyclotomic $\Z_2$-extension has small index because of the congruences
(\lref{UpperTheory.V.Two.orbTwo}). A homomorphism tamely ramified at these places is determined, modulo the unramified ones, by its
values on one generator of each tame inertia group, so allowing ramification at $p$ multiplies the bound by at most $2^{c\deg M}$
(\lref{UpperTheory.V.Two.charStep}). Hence $\#\mathcal C(A^{(i)},qd)\le2^{g+6\deg M}$. \emph{Adjoining $\Q(\sqrt{-p})$:} this replaces
$U_{A,\infty}$ by a subgroup of index at most $2$, and the number of homomorphisms to $\Z/2$ with ramification allowed at $p$ is at most
squared (\lref{UpperTheory.V.Two.quadStep}; a subgroup of index $2$ of a group with $r$ topological generators modulo squares
has at most $2r$). Two such steps give $\#\mathcal C(A^{(q,d)},qd)\le2^{4(g+6\deg M)}$, and $\mathcal C(A^{(q,d)},1)\subset\mathcal C(A^{(q,d)},qd)$.

\ref{pf:vwin:twobasethm-b} Fields of type $A^{(i)}$ contain $\mu_4$; by Ferrero--Washington there is $g$ with $\#\mathcal C(A^{(i)},1)\le2^g$ for
all $\Omega$ (\lref{UpperTheory.V.Two.abRk}). By \ref{pf:vwin:twobasethm-a}, $\#\mathcal C(A^{(q,d)},1)\le2^R$ with $R=4(g+6\deg M)$.
$\mathcal C(A^{(q,d)},1)$ is the $\F_2$-dual of $\mathfrak X_2(G)/2\mathfrak X_2(G)$, so by Nakayama's lemma $\mathfrak X_2(G)$ is generated by
$R$ elements (\lref{UpperTheory.V.Two.genLE_of_abRk}); and $L$ is totally imaginary and contains $\mu_4$, so $\mathrm{Tow}(2,L,R)$ holds.

\ref{pf:vwin:twobasethm-c} A field $L$ of type $A^{(q,d)}$ contains $\mu_4$, has degree at most $8\deg M$ and embeds into
$\Q(\zeta_{4m_Mqd})$ (\lref{UpperTheory.V.Two.abBasic}). Let $k\ge2$, $E=E_{M,k,q,d}$ and $L_1=L(\mu_{2^{k+2}})$, a layer of degree $2^{m'}$
with $m'\le k+2$ (\lref{UpperTheory.V.Two.layerFn}). The Galois group of $\Omega/L_1$ is contained in that of $\Omega/E$ with index dividing
$2$ (\lref{UpperTheory.V.Two.compIndex}: both are described by conditions on the cyclotomic character and the characters defining $M$,
$Q_q$, $Q_d$), and restriction to a subgroup of index at most $2$ has kernel of length at most $1$ on $\Sha_D(j)$, $\ell=2$
(\lref{UpperTheory.V.Two.shaTwo}). Hence $\vwinsh(E,2,j)\le1+\vwinsh(L_1,2,j)$.
\end{proof}

\subsubsection{The window}

\begin{lemma}[\cref{lem:window:window}]\label{pf:vwin:window}
Let $R\in\N$, $0<\eta\le C_t$, $X>1$ and $\vwinexc\subset\N$ a finite set with $|\vwinexc|\le R$. There is
$Y\in\mathrm{W}_J(X)$ with $J=2^{R+2}(C_t/\eta)^R$ such that every $x\in[Y,Y^2]$ is $(\eta,C_t)$-good for $\vwinexc$.
\formallabel{UpperTheory.V.Win.window}
\formallabel{UpperTheory.V.Win.window1}
\end{lemma}

\begin{proof}
We first show by induction on $R$: \emph{for $\rho\ge1$, $U>0$ and a finite set $T\subset\R$ with $|T|\le R$ there is $y$ with
$U/(2^{R+2}\rho^R)\le y$, $2y\le U$, and such that $a\in T$, $y\le s\le2y$ and $a\le s$ imply $\rho a\le s$.} For $R=0$ take
$y=U/2$. For the inductive step, if some $a_0\in T$ satisfies $U/(2\rho)<a_0\le U$, apply the induction hypothesis to
$U/(2\rho)$ and $T\smallsetminus\{a_0\}$; the resulting $y$ satisfies $2y\le U/(2\rho)<a_0$, so $a_0$ never satisfies $a_0\le s\le2y$,
and $U/(2^{R+3}\rho^{R+1})=(U/(2\rho))/(2^{R+2}\rho^R)\le y$. Otherwise take $y=U/2$: if $a\in T$ and $a\le s\le U$, then
$a\le U/(2\rho)$, so $\rho a\le U/2\le s$.

Apply this with $\rho=C_t/\eta$, $U=\log X$ and $T=\{\log e/C_t:e\in\vwinexc\}$ (with $\log 0$ read as any real; $e=0$ is harmless
below), and put $Y=e^y$. Then $Y^2=e^{2y}\le X$ and $X^{1/J}=e^{U/J}\le e^y=Y$. Let $x\in[Y,Y^2]$ and $e\in\vwinexc$ with
$0<e\le x^{C_t}$; with $s=\log x\in[y,2y]$ and $a=\log e/C_t\le s$ we get $\rho a\le s$, i.e.\ $\log e/\eta\le\log x$, i.e.\
$e\le x^\eta$. For $e=0$ the conclusion $0\le x^\eta$ is trivial.
\end{proof}

\subsubsection{Bounds at a point and the proof of the two-part theorem}

\begin{lemma}\label{pf:vwin:point2}
\begin{enumerate}[label=(\alph*)]
\item\label{pf:vwin:p2-exc} If $\sfname{ABound}(\ell,R,L_0)$ holds, there is $c$ with $\mu_{\ell^c}\subset L_0$ such that for every
$Z\in\N$ there is a finite set $\vwinexc$ with $\mathrm{F}(\ell,R,c,L_0,Z,\vwinexc)$.
\item\label{pf:vwin:p2-pt} Let $C_K>0$ with $\sfname{LemmaK}(C_K)$, $L_0$ a field containing $\mu_{q_\ell}$ with
$\mathrm{S}(\ell,R,L_0)$ and $\mathrm{F}(\ell,R,c,L_0,Z,\vwinexc)$, $L_2/L_0$ a layer of degree $\ell^{j_0}$ with $[L_2:\Q]\le n_0$, and
$\eps>0$, $c_0\ge0$, $x\ge1$ such that $x$ is $(\eps/(2C_K),C_t)$-good for $\vwinexc$ and $\log(2|d_{L_2}|)\le c_0x^{\eps/4}$.
Then for $1\le e\le Z$ with $e\le x^{C_t}$ and every layer $L_1/L_0$ of degree $\ell^m$, $\vwinsh(L_1,\ell,e)$ is finite and
\[
 \vwinsh(L_1,\ell,e)\le R\bigl(m+j_0+c+\log_\ell Z\bigr)+2\exp\bigl((2n_0)^{C_K}c_0x^{3\eps/4}\bigr).
\]
\end{enumerate}
\formallabel{UpperTheory.V.Win.exc}
\formallabel{UpperTheory.V.Win.point}
\formallabel{UpperTheory.V.Win.point2}
\end{lemma}

\begin{proof}
\ref{pf:vwin:p2-exc} Take the data $K,c,R',s_i$ of $\sfname{ABound}$. Let $\vwinexc_i$ be the set of $1\le e\le Z$ with
$\|s_i-e\|<\ell^{-\log_\ell Z}$ and $\vwinexc=\bigcup_i\vwinexc_i$. Each $\vwinexc_i$ has at most one element: two distinct
$e,e'\le Z$ satisfy $\|e-e'\|=|e-e'|_\ell\ge\ell^{-\log_\ell Z}$, since $|e-e'|<\ell^{\log_\ell Z+1}$, while two elements of
$\vwinexc_i$ would be at distance $<\ell^{-\log_\ell Z}$ by the ultrametric inequality (\lref{UpperTheory.V.Win.near1}). So
$|\vwinexc|\le R'\le R$, and for $e\notin\vwinexc$ the first bound of $\sfname{ABound}$ applies with $B=\log_\ell Z$.

\ref{pf:vwin:p2-pt} If $e\notin\vwinexc$, use the first bound. If $e\in\vwinexc$, then $e\le x^{\eps/(2C_K)}$ by goodness. Let $D$ be an
$\ell$-datum of $L_2$ (it exists, and $L_2$ contains $\mu_{q_\ell}$ and is totally imaginary). By Lemma K, $H^2(G_{L_2,S_\ell},\Zl(e+1))$
is finite with $\log\#H^2\le\exp((2e[L_2:\Q])^{C_K}\log(2|d_{L_2}|))\le\exp((2n_0)^{C_K}x^{\eps/2}\cdot c_0x^{\eps/4})$, so
$\vwinVL(D,e)=\log\#H^2/\log\ell\le2\exp((2n_0)^{C_K}c_0x^{3\eps/4})$. The second bound gives
$\vwinsh(L_1,\ell,e)\le\vwinVL(D,e)+R\max(m-j_0,0)$.
\end{proof}

For an abelian datum $M$ let $\mathrm{B}_a$ be as in \cref{pf:vwin:stmtdef}\ref{pf:vwin:stmt-bab}.

\begin{lemma}\label{pf:vwin:parts-ab}
Let $C_K>0$ with $\sfname{LemmaK}(C_K)$, $J>0$, $C_t\ge0$, $\eps>0$ and $A_1\in\R$.
\begin{enumerate}[label=(\alph*)]
\item\label{pf:vwin:ab-two} Let $M$ be an abelian datum and $R\in\N$. There is $x_0$ such that for every $(A_1,C_t)$-admitted $(m,t)$, every
$X$, every prime pair $(m_M;q,d)$ with $\log(qd)\le X^{\eps/(4J)}$ and every $2$-base $L_0$ for $(M,q,d)$ with $\sfname{ABound}(2,R,L_0)$,
there is a set $\vwinexc$ with $|\vwinexc|\le R$ such that $\mathrm{B}_a(M,\eps,m,t,q,d,x)$ holds for every $x\ge x_0$ with
$X^{1/J}\le x\le X$ which is $(\eps/(2C_K),C_t)$-good for $\vwinexc$.
\item\label{pf:vwin:ab-bad} Let $\ell$ be an odd prime, $R\in\N$, $A$ an abelian datum with $\ell\mid m_A$ and $H_A\subset\{a\equiv1\bmod\ell\}$, and
suppose that some field $L_0$ of type $A$ satisfies $\sfname{ABound}(\ell,R,L_0)$. There is $x_0$ such that for every admitted $(m,t)$ and
every $X$ there is a set $\vwinexc$ with $|\vwinexc|\le R$ such that for every $x\ge x_0$ with $X^{1/J}\le x\le X$ which is
$(\eps/(2C_K),C_t)$-good for $\vwinexc$, every field $L_0'$ of type $A$ and every layer $L_1/L_0'$ of degree $\ell^{j_0(\ell,R)}$ satisfy
$\sfname{VVat}(\ell,L_1,t(x),\exp(x^{4\eps/5}))$.
\end{enumerate}
\formallabel{UpperTheory.V.Win.partA}
\formallabel{UpperTheory.V.Win.partB}
\end{lemma}

\begin{proof}
Both parts rest on \cref{pf:vwin:point2}\ref{pf:vwin:p2-pt} with $Z=\lfloor X^{C_t}\rfloor$ and the following elementary facts
(\lref{UpperTheory.V.Win.logs}, \lref{UpperTheory.V.Win.final}, \lref{UpperTheory.V.Win.vvat}): if $X^{1/J}\le x$, then
$\log_\ell\lfloor X^{C_t}\rfloor\le2C_tJ\log x$ and $X^{\eps/(4J)}\le x^{\eps/4}$; if $2^m\le x^{A_1}$ then $m+2\le(2|A_1|+2)(1+\log x)$; for
every $C$ there is $x_0$ with $x^{C_t}\bigl(C(1+\log x)+2\exp(Cx^{3\eps/4})\bigr)\le\exp(x^{4\eps/5})$ for $x\ge x_0$; and if
$\vwinsh(E,\ell,e)$ is finite and at most $b$ for $1\le e<t$, then $\sfname{VVat}(\ell,E,t,Tb)$ for every $T\ge t$.

\ref{pf:vwin:ab-two} Let $L_0$ be a $2$-base. Its degree is at most $8\deg M$, so $\mu_{2^c}\subset L_0$ implies $2^c\le2[L_0:\Q]$ and $c\le16\deg M$
(\lref{UpperTheory.V.Win.rootdeg}). Let $L_2/L_0$ be a layer of degree $2^{j_0}$. Since $L_0$ embeds into $\Q(\zeta_{4m_Mqd})$ and has bounded
degree, there are $n_0$ and $c_0$, depending only on $M$ and $R$, with $[L_2:\Q]\le n_0$ and $\log(2|d_{L_2}|)\le c_0w$ whenever
$\log(qd)\le w$ and $w\ge1$ (\lref{UpperTheory.V.Win.layerdisc}); we use $w=x^{\eps/4}$. Let $\vwinexc$ be given by
\cref{pf:vwin:point2}\ref{pf:vwin:p2-exc} for $Z=\lfloor X^{C_t}\rfloor$. For $E=E_{M,m(x),q,d}$, the $2$-base property gives a layer $L_1/L_0$ of
degree $2^{m'}$, $m'\le m(x)+2\le(2|A_1|+2)(1+\log x)$, with $\vwinsh(E,2,e)\le1+\vwinsh(L_1,2,e)$. By \cref{pf:vwin:point2}\ref{pf:vwin:p2-pt}, for
$1\le e<t(x)$ (note $e\le x^{C_t}\le X^{C_t}$),
$1+\vwinsh(L_1,2,e)\le C(1+\log x)+2\exp(Cx^{3\eps/4})$ for a constant $C$ depending on $M,R,J,C_t,\eps,A_1,C_K$. Summing over $e<t(x)\le x^{C_t}$
and using the third fact gives $\sfname{VVat}(2,E,t(x),\exp(x^{4\eps/5}))$ for $x\ge x_0$.

\ref{pf:vwin:ab-bad} The same argument with the layer $L_2=L_1$ of degree $\ell^{j_0}$ over the fixed field $L_0$: $n_0=[L_2:\Q]$ and
$c_0=\log(2|d_{L_2}|)$ are constants, and $m=j_0$. Fields of type $A$ are isomorphic, and so are their layers of a given degree
(\lref{UpperTheory.V.Win.abiso}, \lref{UpperTheory.V.Win.layeriso}); since $\vwinsh$ is invariant under isomorphism, the bound
for $L_0$ gives the bound for every $L_0'$.
\end{proof}

\begin{theorem}[\cref{thm:window:bab}]\label{pf:vwin:bab}
Poitou--Tate duality, the local finiteness and duality statements over finite extensions of $\Ql$, the theorem of
Ferrero--Washington (\cref{lit:window:fw}) and the existence of $C_K$ with $\sfname{LemmaK}(C_K)$ imply: there is a constant package
$K$ such that $\mathrm{B}(K,C_t)$ holds for every $C_t\ge1$.
\formallabel{UpperTheory.V.Win.thmBab_of_lemmaK}
\formallabel{UpperTheory.V.Win.thmBab_of}
\end{theorem}

\begin{proof}
By \cref{pf:vwin:tower}, $\mathrm{Tow}(\ell,L,R)$ implies $\sfname{ABound}(\ell,R,L)$, so \cref{pf:vwin:fixedtwo} applies. Replacing $C_K$ by
$\max(C_K,1)$ we may assume $C_K\ge1$, since $\sfname{LemmaK}(C)$ implies $\sfname{LemmaK}(C')$ for $C\le C'$
(\lref{UpperTheory.V.Win.kmono}). Define the constant package $K$ by: $R_0(M)$ an integer $R$ as in
\cref{pf:vwin:fixedtwo}(b); $R_\ell(A)$ an integer $R$ with $\sfname{ABound}(\ell,R,L_0)$ for some field $L_0$ of type $A$ if there is one
(and $0$ otherwise); and $C_K$.

Let $C_t\ge1$, and let $M$, $B'$, $L(\cdot)$, $\eps$ and $A_1$ be as in $\mathrm{B}(K,C_t)$. Put $R=R_0(M)+\sum_{\ell\in B'}R_\ell(L(\ell))$ and
$J=2^{R+2}(2C_KC_t/\eps)^R>0$. Let $x_A$ be given by \cref{pf:vwin:parts-ab}\ref{pf:vwin:ab-two} for $(M,R_0(M))$, and for $\ell\in B'$ let $x_\ell$ be
given by \ref{pf:vwin:ab-bad} for $(\ell,R_\ell(L(\ell)),L(\ell))$ if a field of type $L(\ell)$ exists (if none exists, the statement $\mathrm{B}_b$
for $\ell$ is empty). Put $x_1=\max(1,x_A,\sum_\ell|x_\ell|)$ and $X_0=\max(2,x_1^J)$.

Let $(m,t)$ be admitted, $X\ge X_0$, and $(m_M;q,d)$ a prime pair with $\log(qd)\le X^{\eps/(4J)}$. Then $x_1\le X^{1/J}$. Let $L_0$ be a $2$-base
for $(M,q,d)$ with $\sfname{ABound}(2,R_0(M),L_0)$, and let $\vwinexc_A$ and $\vwinexc_\ell$ ($\ell\in B'$) be the exceptional sets of
\cref{pf:vwin:parts-ab}. Their union $\vwinexc$ has at most $R$ elements. Put $\eta=\eps/(2C_K)$; then $0<\eta\le1\le C_t$. By
\cref{pf:vwin:window}, applied with $J'=2^{R+2}(C_t/\eta)^R=J$, there is $Y\in\mathrm{W}_J(X)$ such that every $x\in[Y,Y^2]$ is
$(\eta,C_t)$-good for $\vwinexc$, hence for each of the sets $\vwinexc_A,\vwinexc_\ell$. For such $x$ we have $X^{1/J}\le x\le X$ and
$x\ge x_1$, so \cref{pf:vwin:parts-ab} gives $\mathrm{B}_a$ and $\mathrm{B}_b$ at $x$. Hence $Y$ is a $K$-window.
\end{proof}

\begin{remark}\label{pf:vwin:rem-window}
The window is forced by the exceptional twists: for $e\in\vwinexc$ only the second bound, which goes through Lemma K and is
doubly exponential in $[L_2:\Q]$ and polynomial in $e^{C_K}$, is available; it is harmless only when $e\le x^{\eps/(2C_K)}$. The
pigeonhole argument of \cref{pf:vwin:window} produces scales $x$ at which no exceptional twist lies in $(x^{\eta},x^{C_t}]$.
\end{remark}


% ---------------------------------------------------------------------------
\subsection{The trivial character}\label{sec:pf:vwin:Z}

For a prime $\ell$, a coefficient ring $\OO$, an algebraic closure $\Omega$ of $\Q$ and $j\ge1$, let $W_j=\OO\otimes\vwinQZ$ with
$G_\Q$ acting through $\vwincyc^{-j}$, and let $\Sha_\ell(\Q,W_j)\subset H^1(G_\Q,W_j)$ be the group of classes vanishing on all
inertia groups and on the decomposition groups of the places above $\ell$; thus $V_\ell(1,j)\le\len_\OO\Sha_\ell(\Q,W_j)$ for every
host, since the conditions at the primes dividing $N_h$ only shrink the group (\lref{UpperTheory.V.Chr.V'_le_V}).

\begin{proposition}[\cref{prop:window:propZ}]\label{pf:vwin:propZ}
Let $\ell\ge5$ be prime and $j\ge2$ even with $(\ell-1)\mid2j$. Then $\Sha_\ell(\Q,W_j)=0$ for every coefficient ring $\OO$ and every
$\Omega$. Consequently $\sfname{CountZ}$ holds, with the constant $C=0$.
\formallabel{UpperTheory.V.X.propZ}
\formallabel{UpperTheory.V.HV.countZ}
\end{proposition}

The proof distinguishes the cases $(\ell-1)\mid j$ and $(\ell-1)\nmid j$; in the second case $j\equiv\frac{\ell-1}2\pmod{\ell-1}$. Both use
the following control theorem. For a number field $k$, a closed subgroup $H\subset G_k$ and a discrete $G_k$-module $M$, let
$H^1_\Sha(H,M)\subset H^1(H,M)$ be the group of classes whose restrictions to $H\cap I_A$ for all valuation rings $A$ of $\Omega$, and to
$H\cap D_A$ for $A$ above $\ell$, vanish; thus $H^1_\Sha(G_\Q,W_j)=\Sha_\ell(\Q,W_j)$.

\begin{proposition}[Control]\label{pf:vwin:control}
Let $k$ be a number field, $\OO$ a coefficient ring, $U\subset G_k$ an open normal subgroup such that $\ell$ is odd or
$\vwincyc(U)\subset1+4\Z_2$, $\psi:G_k\to\OO^\times$ a character trivial on $U$, $n\ne0$, and $W=\OO\otimes\vwinQZ$ with $G_k$ acting through
$\psi\vwincyc^n$. Let $U_\infty=U\cap\ker\vwincyc$. Then restriction $H^1_\Sha(U,W)\to H^1_\Sha(U_\infty,W)$ is injective, and its image is the
group of classes fixed by the conjugation action of $U$.
\formallabel{UpperTheory.V.Base.control}
\formallabel{UpperTheory.V.Base.controlGamma}
\end{proposition}

\begin{proof}
\emph{Case 1: $\vwincyc(U)\subset1+q_\ell\Zl$.} Since $U$ is open, $\vwincyc(U)$ is open in $\Zl^\times$, so $U\not\subset\ker\vwincyc$, and the
same holds for $U\cap D_A$ for every $A$ above $\ell$ (\lref{UpperTheory.V.Base.decInf}). There are a continuous surjection
$\kappa:U\to\Zl$ with kernel $U_\infty$ and $\gamma\in U$ with $\kappa(\gamma)=1$ and $\vwincyc(\gamma)$ of infinite order
(\lref{UpperTheory.V.Base.gammaGen}; as in \cref{pf:vwin:gengal}, $\vwincyc(U)\subset1+q_\ell\Zl\cong\Zl$ is procyclic and torsion-free).
\emph{Injectivity.} $U$ acts on $W$ through $\vwincyc^n$, and $U_\infty$ acts trivially. If the restriction of a cocycle $z$ to $U_\infty$ is a
coboundary, we may assume that it vanishes on $U_\infty$. As $\vwincyc(\gamma)^n\ne1$, $m\mapsto\gamma m-m$ is multiplication by a non-zero
element of $\OO$, hence surjective on the divisible module $W$; and $\vwincyc(U)$ is contained in the closure of the group generated by
$\vwincyc(\gamma)$. By \cref{pf:vwin-he:cob} (for $\Gamma=U$ and $\omega=\vwincyc|_U$), $z$ is a coboundary.
\emph{Image.} Restrictions of classes of $U$ are fixed by $U$. Conversely, a $U$-invariant class of $U_\infty$ is the restriction of a class $y$
of $U$, because $U/U_\infty\cong\Zl$ is free pro-$\ell$ and $W$ is $\ell$-power torsion (\lref{UpperTheory.V.Base.resZl}). If the
given class satisfies the local conditions, so does $y$: for $A$ not above $\ell$, $\vwincyc$ is trivial on $I_A$, so $U\cap I_A\subset U_\infty$
and the condition is inherited; for $A$ above $\ell$, the restriction of $y$ to $U\cap D_A$ vanishes on $U_\infty\cap D_A$, and the
injectivity argument applied to the open subgroup $U\cap D_A$ of $D_A$ (whose image under $\vwincyc$ is again open in $1+q_\ell\Zl$) shows
that it vanishes.

\emph{Case 2: $\ell$ odd.} Let $V=U\cap\ker(\vwincyc\bmod\ell)$, an open normal subgroup of $G_k$ with $\vwincyc(V)\subset1+\ell\Zl$, $V\cap\ker
\vwincyc=U_\infty$, and $[U:V]$ dividing $\ell-1$, hence a unit in $\OO$. By restriction and corestriction, restriction from $U$ to $V$ is
injective on $H^1_\Sha$, with image the $U$-invariant classes (\lref{UpperTheory.V.Base.selUnit}: the composite of corestriction and
restriction is multiplication by $[U:V]$ on $U$-invariant classes, and the local conditions are inherited because the relative indices of
$V\cap D$ in $U\cap D$ divide $[U:V]$). Composing with Case 1 for $V$ gives the claim for $U$.
\end{proof}

\begin{lemma}\label{pf:vwin:zone}
Let $\ell$ be odd and $j\ge1$ with $(\ell-1)\mid j$. Then $\Sha_\ell(\Q,W_j)=0$.
\formallabel{UpperTheory.V.X.zOne}
\end{lemma}

\begin{proof}
Write $\rho$ for the action on $W=W_j$. Since $\vwincyc(g)^{-j}\equiv1\pmod\ell$ for all $g$ (as $(\ell-1)\mid j$), $\rho$ is trivial on
the $\ell$-torsion $W[\ell]$, and it is trivial on all of $W$ on the kernel of $\vwincyc$. We use two inputs.
\begin{enumerate}[label=(\roman*)]
\item (Control; \cref{pf:vwin:control} for $k=\Q$, $U=G_\Q$, $\psi=1$, $n=-j$.) Restriction from $G_\Q$ to the closed subgroup
$\ker\vwincyc$ is injective on $\Sha_\ell(\Q,W)$.
\item (Kronecker--Weber; \lref{UpperTheory.V.X.kwl}.) Every continuous homomorphism $f:G_\Q\to\Z/\ell$ which vanishes on the inertia
groups of all places not above $\ell$ vanishes on $\Gal(\Omega/\Q(\mu_{\ell^2}))$: the field cut out by $f$ is abelian and unramified
outside $\ell$, hence contained in some $\Q(\mu_{\ell^{n}})$, and a cyclic extension of degree $\ell$ inside $\Q(\mu_{\ell^\infty})$ is
contained in $\Q(\mu_{\ell^2})$.
\end{enumerate}
Let $x\in\Sha_\ell(\Q,W)$ with $\ell x=0$, represented by a cocycle $z$. Then $\ell z=\rho(\cdot)m-m$ for some $m\in W$; choose $m'$ with
$\ell m'=m$ ($W$ is divisible) and put $z'(g)=z(g)-(\rho(g)m'-m')$. Then $\ell z'=0$, so $z'$ takes values in $W[\ell]$, on which $\rho$
is trivial; hence $z'$ is a continuous homomorphism. At a place $A$ not above $\ell$, $\vwincyc$ and hence $\rho$ are trivial on the
inertia group, and $z$ is a coboundary there, so $z$ and $z'$ vanish on the inertia group. Applying (ii) to the compositions of $z'$
with all homomorphisms $W[\ell]\to\Z/\ell$ (which separate points of $W[\ell]$) shows that $z'$ vanishes on
$\Gal(\Omega/\Q(\mu_{\ell^2}))\supset\ker\vwincyc$. On $\ker\vwincyc$ the action is trivial, so $z$ restricted to $\ker\vwincyc$ is zero,
and $x=0$ by (i). Thus $\Sha_\ell(\Q,W)$ has no non-zero element killed by $\ell$. Every element is killed by a power of $\ell$, since
cocycles are locally constant and $W$ is $\ell$-power torsion; so $\Sha_\ell(\Q,W)=0$.
\end{proof}

\begin{proof}[Proof of \cref{pf:vwin:propZ}]
If $(\ell-1)\mid j$, apply \cref{pf:vwin:zone}. Otherwise $(\ell-1)\mid2j$ and $(\ell-1)\nmid j$, $j$ is even and $\ell\ge5$, and the
vanishing is \cref{pf:vwin:zquad} below.

For $\sfname{CountZ}$: let $h$ be a host with character $\chi$, big for $t$, $\ell$ odd, $0<j<t$ and $\theta\in\vwinX_h(\OO,j)$ in case Z,
i.e.\ $\ell\notin B_\chi$, $(\ell-1)\mid2j$ and $\theta=1$. Then $\ell\ne3$ because $3\in B_\chi$, so $\ell\ge5$; $j$ is even because
$\theta=1$ satisfies $\theta(-1)=(-1)^j$ and $2\ne0$ in $\OO$; and $j\ge2$. Hence $V_\ell(1,j)\le\len_\OO\Sha_\ell(\Q,W_j)=0$, and $(\theta,j)$
is bounded by $(\ell N_hj)^0$.
\end{proof}

\begin{lemma}[The quadratic case]\label{pf:vwin:zquad}
Let $\ell\ge5$ be prime and $j\ge2$ even with $(\ell-1)\mid2j$ and $(\ell-1)\nmid j$. Then $\Sha_\ell(\Q,W_j)=0$ for every coefficient
ring $\OO$ and every $\Omega$.
\formallabel{UpperTheory.V.X.zQuad}
\end{lemma}

\begin{proof}
Let $\alpha=\vwincyc^{-j}:G_\Q\to\OO^\times$ be the action on $W=W_j$, and write $\bar x$ for the reduction modulo $\ell$ of $x\in\Zl$.

\emph{(1) Arithmetic of $j$.} Write $\ell=2n+1$. Then $n\mid j$, say $j=nk$; $k$ is odd since $(\ell-1)\nmid j$, so $n$ is even as $j$ is,
i.e.\ $\ell\equiv1\pmod4$ and $\frac{\ell-1}2\mid j$. Let $u\in\Zl^\times$ reduce to a generator of $\F_\ell^\times$; then $\bar u^j\ne1$.

\emph{(2) The quadratic field.} As $\ell\equiv1\pmod 4$, a square root $s\in\Omega$ of $\ell$ lies in $\Q(\mu_\ell)$ (quadratic Gauss
sum), and $g\in G_\Q$ fixes $s$ iff $\overline{\vwincyc(g)}$ is a square in $\F_\ell^\times$, i.e.\
$\overline{\vwincyc(g)}^{(\ell-1)/2}=1$ (Euler's criterion); otherwise $g(s)=-s$. By (1), $\overline{\vwincyc(g)}^{\,j}=1$ whenever
$g(s)=s$. Let $K=\Q(s)$, a real quadratic field of discriminant $\ell$ with $G_K=\{g:g(s)=s\}$; its class number is smaller than
$\ell$ (\lref{UpperTheory.V.X.hQuad}, from Minkowski's bound), hence prime to $\ell$.

\emph{(3) Two consequences.} (i) If $g(s)=s$, then $\alpha(g)$ acts trivially on $W[\ell]$, because $\vwincyc(g)^j\equiv1\pmod\ell$.
(ii) The cyclotomic character maps every inertia group $I_A$ at a place $A$ above $\ell$ onto $\Zl^\times$; choose $\delta\in I_{A_0}$, for
one such $A_0$, with $\vwincyc(\delta)=u$. Then $\alpha(\delta)-1$ is a unit of $\OO$, since $u^j-1$ is a unit by (1); in particular
$\delta(s)\ne s$ by (i).

\emph{(4) Reduction to $\ell$-torsion cocycles.} Since $\alpha(\delta)-1$ is a unit, every class of $H^1(G_\Q,W)$ has a unique cocycle
representative $c$ with $c(\delta)=0$; and every class of $\Sha_\ell(\Q,W)$ is killed by some power of $\ell$. Hence it suffices to
show: every cocycle $c$ with $c(\delta)=0$, $\ell c=0$, vanishing on the inertia groups at places not above $\ell$ and whose
restriction to each decomposition group $D_A$ above $\ell$ is a coboundary, is zero (\lref{UpperTheory.V.ML.shaToShT}; the
statement for $\ell^n c=0$ follows by induction on $n$, applying it to $\ell^{n-1}c$).

\emph{(5) Vanishing on $G_K$.} Let $c$ be such a cocycle. For $A$ above $\ell$ write $c(g)=\alpha(g)m-m$ on $D_A$, and choose
$t\in I_A$ with $\vwincyc(t)=u$. Then $(\alpha(t)-1)\ell m=\ell c(t)=0$, so $\ell m=0$ by (3)(ii), and for $g\in D_A\cap G_K$,
$c(g)=\alpha(g)m-m=0$ by (3)(i). On $G_K$, $c$ takes values in $W[\ell]$, on which $G_K$ acts trivially, so $c|_{G_K}$ is a continuous
homomorphism vanishing on all inertia groups and on the decomposition groups above $\ell$. By class field theory it factors through the
Artin map onto the class group of $K$ (\lref{UpperTheory.V.Kum.pullOnto}; $W$ has no $2$-torsion, which takes care of the real
places). As $\ell c=0$ and the class number of $K$ is prime to $\ell$, $c|_{G_K}=0$.

\emph{(6) Conclusion.} If $g\notin G_K$, then $g(s)=-s=\delta(s)$, so $\delta^{-1}g\in G_K$ and
$c(g)=c(\delta)+\alpha(\delta)c(\delta^{-1}g)=0$. Hence $c=0$.
\end{proof}


% ---------------------------------------------------------------------------
\subsection{The minus part}\label{sec:pf:vwin:minus}

In this subsection we prove \cref{prop:psibound}. Throughout, $\Omega$ is an algebraic closure of $\Q$, $\ell$ an odd prime, $\OO$ a
discrete valuation ring which is a finite $\Zl$-algebra, with valuation $v$, and $W=\OO\otimes\vwinQZ$. For a continuous character
$\alpha:G_\Q\to\OO^\times$ let $W_\alpha$ be $W$ with $G_\Q$ acting through $\alpha$, and $\alpha^*=\vwincyc\alpha^{-1}$. A
\emph{cocycle for $\alpha$} is a locally constant $c:G_\Q\to W$ with $c(gh)=c(g)+\alpha(g)c(h)$. For a valuation ring $A$ of $\Omega$
(a place) we write $D_A$, $I_A$ for its decomposition and inertia groups, and say $A$ is \emph{over $n$} if its maximal ideal contains
$n$. For $n\ge1$, $\Sha_n(\Q,W_\alpha)$ is the group of classes vanishing on all inertia groups and on the decomposition groups of the
places over $n$, as in \cref{pf:vwin-he:V}. Recall $\alpha_k(\chi)=N\cdot\mathrm{primorial}(k+1)\cdot B_{k,\chi^{\mathrm{prim}}}$ for a
Dirichlet character $\chi$ modulo $N$ (\cref{sec:window:minus}).

\begin{proposition}[\cref{prop:psibound}]\label{pf:vwin:psibound}
\begin{enumerate}[label=(\roman*)]
\item Poitou--Tate duality and the orthogonality of unramified classes under local duality imply $\sfname{PsiBound}$
(\cref{def:window:psibound}).
\item $\sfname{PsiBound}$ implies $\sfname{ThmBcFree}$.
\end{enumerate}
\formallabel{UpperTheory.V.ML.iv}
\formallabel{UpperTheory.V.ML.psiOf}
\formallabel{UpperTheory.V.Odd.thmBcFree_of}
\formallabel{UpperTheory.V.Odd.thmBc_of}
\end{proposition}

\begin{definition}[$\sfname{PsiBound}$, formal form]\label{pf:vwin:psidef}
$\sfname{PsiBound}$ is the statement of \cref{def:window:psibound} in the following form: for every abelian CM field $E$, every
$N\ge1$ and every homomorphism $\rho:(\Z/N)^\times\to\Gal(E/\Q)$ compatible with the cyclotomic action (for every $b$, $\rho(b)$ is the
restriction of every automorphism of $\Omega$ acting on $\mu_N$ by $b$), every prime $\ell\nmid2[E:\Q]$, every subring $\OO\subset\Qlbar$
which is a discrete valuation ring finite over $\Zl$, every $j\ge1$, every character $\theta:\Gal(E/\Q)\to\OO^\times$ with
$(\theta\circ\rho)(-1)=(-1)^{j+1}$, and every $a\in\OO$ equal to $\alpha_{j+1}(\theta\circ\rho)$ in $\Qlbar$, the $\OO$-length of the
$\theta$-eigenspace of $\Sha^1(G_{E,S_\ell},\OO\otimes\vwinQZ(-j))$ is at most $v_\OO(a)$.
\formallabel{UpperTheory.V.Odd.PsiBound}
\end{definition}

\subsubsection{Reciprocity at an auxiliary prime}

\begin{lemma}[\cref{lem:window:rec}]\label{pf:vwin:rec}
Assume Poitou--Tate duality and the orthogonality of unramified classes. Let $\alpha:G_\Q\to\OO^\times$ be a character, $0\ne a\in\OO$,
and $w_0\in W$ with $\{y\in\OO:yw_0=0\}=a\OO$; assume that $\{g:\alpha(g)\equiv1\bmod a\}$ is open. Let $N\ne0$ with $\ell\mid N$ and
$q\ne\ell$ a prime. Let $c$ and $n$ be cocycles for $\alpha$ and $\alpha^*$ killed by $a$ such that:
\begin{enumerate}[label=(\alph*)]
\item at every place $A$ not over $Nq$, $\alpha\equiv1\bmod a$ on $I_A$, and $c$ and $n$ vanish on $I_A$;
\item at every place $A$ over $N$ and not over $q$, $c|_{D_A}$ is the coboundary of an element killed by $a$;
\item at a place $A_0$ over $q$, $\alpha\equiv\alpha^*\equiv1\bmod a$ on $D_{A_0}$, and $c$ vanishes on $I_{A_0}$.
\end{enumerate}
If $g\in D_{A_0}$, $\tau\in I_{A_0}$ and $y\in\OO$ satisfy $n(\tau)=yw_0$, then $y\,c(g)=0$.
\formallabel{UpperTheory.V.ML.recip}
\formallabel{UpperTheory.V.ML.Rec}
\end{lemma}

\begin{proof}
Let $r$ be such that $a\mid\ell^r$ (possible since $w_0$ is killed by a power of $\ell$), $B=\OO/a$, and $\iota:B\to W$, $\iota(\bar u)=uw_0$;
$\iota$ is injective and its image is the $a$-torsion $W[a]$ (which is cyclic). Write $c=\iota\circ c'$ and $n=\iota\circ n'$; then $c'$ and $n'$
are cocycles with values in $B$ for $\bar\alpha=\alpha\bmod a$ and $\bar\alpha^*$, and $\bar\alpha\bar\alpha^*=\vwincyc\bmod a$ factors
through the cyclotomic character modulo $\ell^r$. Since $y\,c(g)=\iota(n'(\tau)c'(g))$, we must show $n'(\tau)c'(g)=0$ in $B$. Suppose not.
The finite $\Z/\ell^r$-module $B$ has a $\Z/\ell^r$-linear form $\lambda$ with $\lambda(n'(\tau)c'(g))\ne0$ ($\Z/\ell^r$ is self-injective).

Consider the function $(g_1,g_2)\mapsto\lambda\bigl(n'(g_1)\,\bar\alpha(g_1)\,c'(g_2)\bigr)$, the cup product of $n'$ and $c'$ composed with
$\lambda$; it is a $2$-cocycle of $G_\Q$ with values in $\Z/\ell^r(1)$, the Tate twist, because $\bar\alpha^*\bar\alpha=\vwincyc$. Let $S$
be the set of primes dividing $Nq$; by (a) the cocycle factors through the Galois group of the maximal extension unramified outside $S$
(the kernel of $\vwincyc\bmod\ell^r$, $\alpha\bmod a$, $c'$, $n'$ on the inertia groups outside $S$ is normal and closed). Its local
restrictions are coboundaries: at the primes of $S$ dividing $N$ and not $q$ by (b); at the places outside $S$ by the orthogonality of
unramified classes; and the archimedean place does not intervene as $\ell$ is odd. By Poitou--Tate duality (the exactness of
$H^2(G_{\Q,S},\Z/\ell^r(1))\to\bigoplus_{v\in S}\mathrm{Br}(\Q_v)[\ell^r]\to\Z/\ell^r$, the second map being the sum of the local invariants,
which are isomorphisms), the restriction to $D_{A_0}$ is then also a coboundary. Finally, a local computation at $q$: on $D_{A_0}$ both
characters are trivial modulo $a$, $c'$ is a homomorphism vanishing on $I_{A_0}$ and $n'$ is a homomorphism; the restriction of the cup product
$n'\cup c'$ is a coboundary only if $\lambda(n'(\tau)c'(g))=0$ for all $\tau\in I_{A_0}$ and $g\in D_{A_0}$ (the local pairing of the unramified
class $c'$ with $n'$ is the value of $n'$ on a generator of tame inertia paired with the value of $c'$ at a Frobenius element). This
contradicts the choice of $\lambda$.
\end{proof}

\subsubsection{The Euler system of Gauss sums and the Kolyvagin induction}

Fix $N_0$ prime to $\ell$, $m\ge1$, $F=N_0\ell^m$ and a character $\alpha$. For $t\in(\Z/F)^\times$ let $\sigma_t\in G_\Q$ act on $\mu_F$ by $t$.

\begin{definition}\label{pf:vwin:mldefs}
\begin{enumerate}[label=(\alph*)]
\item $\mathrm{Dat}(N_0,m,F,\alpha)$ means: $F=N_0\ell^m$ with $(N_0,\ell)=1$ and $m\ge1$; $\alpha\equiv1\bmod\ell^m$ on
$\Gal(\Omega/\Q(\mu_F))$; $\alpha(g)\equiv-1\bmod\ell^m$ for every $g$ acting on $\mu_F$ by $-1$; and there is $g$ with
$\alpha(g)-\vwincyc(g)\in\OO^\times$.
\item For $\delta\in G_\Q$ and $a\in\OO$, $\mathrm{T}(\alpha,\delta,F,a)$ is the $\OO$-module of cocycles $c$ for $\alpha$ such that
$c(\delta)=0$, $ac=0$, $c$ vanishes on $I_A$ for every place $A$ not over $F$, and for every place $A$ over $F$, $c|_{D_A}$ is the coboundary of
an element killed by $a$. For $n\ge1$, $\mathrm{T}_n(\alpha,\delta,F,a)$ is its submodule of cocycles vanishing on $D_A$ for all $A$ over $n$.
\item For $b<F$ and $t\in(\Z/F)^\times$ (represented in $[0,F)$) let $\vartheta_b(t)=\lfloor tb/F\rfloor-\lfloor(F-t)b/F\rfloor$, and
$z_b=z_b(\alpha,\sigma)=\sum_{t\in(\Z/F)^\times}\vartheta_b(t)\,\alpha(\sigma_t)^{-1}\in\OO$.
\item A prime $q$ is \emph{admissible} (for $r$) if $q\equiv1\bmod F$, $q\ne r$, and $\ell\nmid\ord(r\bmod q)$; an admissible list is a
finite list of distinct admissible primes.
\item For $M,n\ge1$ and a character $\beta$, a function $\nu:G_\Q\to W$ is a \emph{Kummer class} of type $(M,F,n,\beta)$ if on
$\Gal(\Omega/\Q(\mu_F))$ it is a locally constant homomorphism killed by $M$, satisfies $\nu(gug^{-1})=\beta(g)\nu(u)$ there, and $\nu$
vanishes on $I_A$ for every place $A$ not over $n$.
\item An \emph{Euler system} for $(r,\beta,A_0,w_1)$, where $r$ is a prime, $A_0$ a place over $r$ and $w_1\in W$, is a family $\nu(L)$,
indexed by admissible lists, such that: $\nu(L)$ is a Kummer class of type $(\ell^m,F,Fr\prod L,\beta)$; for $q\notin L$ with $q::L$
admissible, every place $A$ over $q$ and every Frobenius element $\varphi$ at $A$, there are $\tau\in I_A$ and $u\in\OO^\times$ with
$\nu(q::L)(\tau)=u\,\nu(L)(\varphi)$; and there are $\tau_0\in I_{A_0}$ and $u_0\in\OO^\times$ with $\nu(\emptyset)(\tau_0)=u_0w_1$.
\end{enumerate}
\formallabel{UpperTheory.V.ML.Dat}
\formallabel{UpperTheory.V.ML.zVal}
\formallabel{UpperTheory.V.ML.NuSys}
\end{definition}

The modules $\mathrm{T}$, admissible primes and Kummer classes are the shared definitions \lref{UpperTheory.V.ML.ShT},
\lref{UpperTheory.V.ML.AdmP} and \lref{UpperTheory.V.ML.IsNu}.

The normalization $c(\delta)=0$, for $\delta$ with $\alpha(\delta)-1\in\OO^\times$, selects one cocycle in each class: a coboundary
$g\mapsto\alpha(g)w-w$ vanishing at $\delta$ is zero. The elements $z_b$ are images of Stickelberger elements. The Euler system is
provided by Gauss sums:

\begin{proposition}[Euler system of Gauss sums]\label{pf:vwin:nusys}
Assume $\mathrm{Dat}(N_0,m,F,\alpha)$. Let $r\equiv1\bmod F$ be a prime, $b<F$, $A_0$ a place over $r$, $\sigma$ a choice of lifts
$\sigma_t$, and $w\in W$ with $\ell^mw=0$. Then there is an Euler system for $(r,\alpha^*,A_0,z_b\,w)$
(\lref{UpperTheory.V.ML.nuSysEx}).
\end{proposition}

The construction (which we do not reproduce) follows Thaine, Kolyvagin and Rubin \cite{Thaine1988,Kolyvagin1990,Rubin1990}: elements
$X_d\in\Q(\zeta_{Fd})$, for squarefree products $d$ of admissible primes, built from Gauss and Jacobi sums at a prime above $r$, with
$X_1$ generating the $\vartheta_b$-twisted ideal of a prime above $r$ by Stickelberger's theorem; Kolyvagin derivatives along the cyclic
groups $(\Z/q)^\times$; Kummer theory over $\Q(\mu_F)$; and projection to the $\alpha^{-1}$-eigenspace (\lref{UpperTheory.V.ML.laxEx},
\lref{UpperTheory.V.ML.derInv}).

\begin{lemma}[The Kolyvagin step]\label{pf:vwin:key}
Assume the conclusion of \cref{pf:vwin:rec} and $\mathrm{Dat}(N_0,m,F,\alpha)$. Let $n\ge1$, $q\equiv1\bmod F$ prime, $\nu$ a Kummer class of type
$(\ell^m,F,Fnq,\alpha^*)$, $\delta$ with $\alpha(\delta)-1\in\OO^\times$, and $a_0\mid a$. Let $a'$ generate the annihilator of the values of $\nu$
on $\Gal(\Omega/\Q(\mu_F))$, let $A$ be a place over $q$, $\varphi$ a Frobenius element at $A$, and $\tau\in I_A$ such that the annihilator of
$\nu(\tau)$ is $a\OO$. Then there is $e\in\OO$ with $a'=ae$ such that $e\,c(\varphi)=0$ for all $c\in\mathrm{T}_n(\alpha,\delta,F,a_0)$, and
\[
 \len_\OO\mathrm{T}_n(\alpha,\delta,F,a_0)\le\len_\OO\mathrm{T}_{nq}(\alpha,\delta,F,a_0)+v(e).
\]
\formallabel{UpperTheory.V.ML.step}
\end{lemma}

\begin{proof}
As $q\equiv1\bmod F$, $A$ is not over $F$, so $\tau$ lies in $\Gal(\Omega/\Q(\mu_F))$ and $a\mid a'$; write $a'=ae$. Let $w_0$ generate
$W[a']$ and write $\nu(\tau)=yw_0$. For $c\in\mathrm{T}_n$, \cref{pf:vwin:rec}, applied to $c$, to a cocycle for $\alpha^*$ obtained from $\nu$,
with $N=Fn$ and the place $A$ (the hypotheses (a)--(c) follow from the definitions of $\mathrm{T}_n$ and of Kummer classes), gives
$y\,c(\varphi)=0$. Since the annihilator of $\nu(\tau)=yw_0$ is $a\OO$, we have $v(y)+v(a)=v(a')$ if
$v(y)<v(a')$; then $e$ is a unit multiple of $y$ and $e\,c(\varphi)=0$. If $v(y)\ge v(a')$, then $\nu(\tau)=0$, so $a$ is a unit, hence so is
$a_0$, and $c=0$ since $a_0c=0$. So $e$ kills all values $c(\varphi)$. Now the map $c\mapsto c(\varphi)$ sends $\mathrm{T}_n$ to the
$e$-torsion of the cyclic module $W[\ell^m]$, whose length is at most $v(e)$, and its kernel is contained in $\mathrm{T}_{nq}$: a cocycle in
$\mathrm{T}_n$ with $c(\varphi)=0$ vanishes on the decomposition groups of the places over $q$ (it is unramified there, and the decomposition
group modulo inertia is generated by $\varphi$; \lref{UpperTheory.V.ML.shtNu}). This gives the length inequality.
\end{proof}

\begin{proposition}[Kolyvagin's induction]\label{pf:vwin:count}
Assume the conclusion of \cref{pf:vwin:rec} and $\mathrm{Dat}(N_0,m,F,\alpha)$. Let $b<F$, $\delta$ with $\alpha(\delta)-1\in\OO^\times$, $\sigma$ a
choice of lifts, and $a_0\in\OO$ with $a_0z_b\mid\ell^m$. Then $\len_\OO\mathrm{T}(\alpha,\delta,F,a_0)\le v(z_b)$.
\formallabel{UpperTheory.V.ML.count}
\end{proposition}

\begin{proof}
Write $\ell^m=z_bk$; then $z_b\ne0$, $k\ne0$ and $a_0\mid k$. Choose a prime $r\equiv1\bmod F$, a place $A_0$ over $r$ and a generator $w$ of
$W[\ell^m]$, and let $\nu$ be an Euler system for $(r,\alpha^*,A_0,z_bw)$ (\cref{pf:vwin:nusys}). For an admissible list $L$ let $a_L$ generate
the annihilator of the values of $\nu(L)$ on $\Gal(\Omega/\Q(\mu_F))$; then $a_L\mid\ell^m$. Put $\mathrm{T}(L)=\mathrm{T}_{r\prod L}(\alpha,\delta,F,a_0)$
and $j(L)=v(a_L)\le J=v(\ell^m)$.

\emph{Base.} The annihilator of $\nu(\emptyset)(\tau_0)=u_0z_bw$ is $k\OO$. \Cref{pf:vwin:key} with $n=1$, $q=r$ gives $e_0$ with $a_\emptyset=ke_0$ and
$\len\mathrm{T}_1\le\len\mathrm{T}(\emptyset)+v(e_0)$; and $\mathrm{T}(\alpha,\delta,F,a_0)=\mathrm{T}_1$ (no place is over $1$).

\emph{Step.} Let $L$ be admissible with $a_0\mid a_L$ and $\mathrm{T}(L)\ne0$, and let $0\ne c\in\mathrm{T}(L)$. Then $c$ does not vanish on
$\Gal(\Omega/\Q(\mu_F))$ (\lref{UpperTheory.V.ML.shtNu}). By Chebotarev's density theorem and the condition $\alpha\not\equiv\vwincyc$ of
$\mathrm{Dat}$, there are an admissible prime $q\notin L$, a place $A$ over $q$ and a Frobenius element $\varphi\in\Gal(\Omega/\Q(\mu_F))$ at $A$ such
that $c(\varphi)\ne0$ and $\nu(L)(\varphi)$ generates the $\OO$-module of values of $\nu(L)$ on $\Gal(\Omega/\Q(\mu_F))$
(\lref{UpperTheory.V.ML.pick}). By the Euler system relation, $\nu(q::L)(\tau)=u\,\nu(L)(\varphi)$ for some $\tau\in I_A$, so the annihilator
of $\nu(q::L)(\tau)$ is $a_L\OO$. \Cref{pf:vwin:key} gives $e$ with $a_{q::L}=a_Le$, $e\,c(\varphi)=0$ and
$\len\mathrm{T}(L)\le\len\mathrm{T}(q::L)+v(e)$. As $c(\varphi)\ne0$, $e$ is not a unit. Hence $a_0\mid a_{q::L}$, $j(q::L)>j(L)$ and
$\len\mathrm{T}(L)+j(L)\le\len\mathrm{T}(q::L)+j(q::L)$.

\emph{Induction.} Since $j\le J$ and $j$ increases strictly at each step, descending induction on $J-j(L)$ gives
$\len\mathrm{T}(L)+j(L)\le J$ for every admissible $L$ with $a_0\mid a_L$ (if $\mathrm{T}(L)=0$ this is $j(L)\le J$). For $L=\emptyset$ (note
$a_0\mid k\mid a_\emptyset$): $\len\mathrm{T}(\emptyset)+v(k)+v(e_0)\le v(z_b)+v(k)$. With the base,
$\len\mathrm{T}(\alpha,\delta,F,a_0)\le\len\mathrm{T}(\emptyset)+v(e_0)\le v(z_b)$.
\end{proof}

\subsubsection{Proof of the bound for the minus components}

We use the following statements of the development, which we do not reproduce.
\begin{enumerate}[label=(M\arabic*)]
\item\label{pf:vwin:M-zeval} (Stickelberger values; \lref{UpperTheory.V.ML.zEval}.) Let $\psi$ be a primitive Dirichlet character of
conductor $f$ with values in $\OO$, $j\ge1$, $\psi(-1)=(-1)^{j+1}$, and $\alpha=(\vwincyc^j\psi)^{-1}$; assume there is $g$ with
$\alpha(g)-\vwincyc(g)\in\OO^\times$, and let $a\in\OO$ equal $\alpha_{j+1}(\psi)$. Then for every $a_0\ne0$ there are $m$, $b$ and lifts $\sigma$
such that, with $f'$ the prime-to-$\ell$ part of $f$ and $F=f'\ell^m$: $\mathrm{Dat}(f',m,F,\alpha)$, $b<F$, $a_0z_b\mid\ell^m$ and $z_b\mid a$. (A
congruence between the Stickelberger values $z_b$ and generalized Bernoulli numbers.)
\item\label{pf:vwin:M-exc} (The exceptional case; \lref{UpperTheory.V.ML.exc}.) If $\alpha(g)-\vwincyc(g)\notin\OO^\times$ for all $g$, and $\alpha$
is unramified at all places not over $\ell$, then $\Sha_\ell(\Q,W_\alpha)=0$.
\item\label{pf:vwin:M-sha} (Normalization; \lref{UpperTheory.V.ML.shaToShT}.) Let $\delta$ with $\alpha(\delta)-1\in\OO^\times$ and $F\ne0$ such that
$\alpha$ is trivial on $I_A$ for $A$ not over $F$, and for every $A$ over $F$ and not over $\ell$ some $g\in I_A$ has $\alpha(g)-1\in\OO^\times$. If
$\len\mathrm{T}(\alpha,\delta,F,a)\le n$ for all $a\ne0$, then $\len\Sha_\ell(\Q,W_\alpha)\le n$.
\item\label{pf:vwin:M-prim} (\lref{UpperTheory.V.ML.primRam}.) If $\psi$ is primitive of conductor $f$ and $\ell\nmid\ord\psi$, then $\psi$ (as a
character of $G_\Q$) is unramified at the places not over $f$, at every place over $f$ some inertia element $g$ has $\psi(g)-1\in\OO^\times$, and
$\psi(g)-1\in\OO^\times$ whenever $\psi(g)\ne1$.
\item\label{pf:vwin:M-eig} (Eigenspaces; \lref{UpperTheory.V.Chr.eigV}.) For an abelian field $E$ with $\ell\nmid[E:\Q]$ and a character $\theta$ of
$\Gal(E/\Q)$, the length of the $\theta$-eigenspace of $\Sha^1(G_{E,S_\ell},\OO\otimes\vwinQZ(-j))$ equals the length of
$\Sha_\ell(\Q,W_{\theta^{-1}\vwincyc^{-j}})$ (inflation--restriction).
\end{enumerate}

\begin{lemma}\label{pf:vwin:vchar}
Assume the conclusion of \cref{pf:vwin:rec}. Let $\psi$ be a Dirichlet character of level $N$ with values in $\OO$, $j\ge1$, $\ell\nmid\ord\psi$,
$\psi(-1)=(-1)^{j+1}$, and $a\in\OO$ equal to $\alpha_{j+1}(\psi)$. Then $\len_\OO\Sha_\ell(\Q,W_\alpha)\le v(a)$ for $\alpha=(\vwincyc^j\psi)^{-1}$.
\formallabel{UpperTheory.V.ML.vChar}
\end{lemma}

\begin{proof}
\emph{Reduction to primitive characters.} Let $\psi_0$ be the primitive character attached to $\psi$, of conductor $f\mid N$. It defines the same
character of $G_\Q$, and $\alpha_{j+1}(\psi)=(N/f)\,\alpha_{j+1}(\psi_0)$, since both are defined with the Bernoulli number of $\psi_0$; so
$v(\alpha_{j+1}(\psi_0))\le v(a)$. Hence we may assume $\psi$ primitive.

\emph{The exceptional case.} Suppose $\alpha(g)-\vwincyc(g)\notin\OO^\times$ for all $g$. At a place $A$ not over $\ell$, $\vwincyc$ is trivial on
$I_A$, so for $g\in I_A$ with $\psi(g)\ne1$ the element $\alpha(g)-\vwincyc(g)=\psi(g)^{-1}-1$ would be a unit by \ref{pf:vwin:M-prim}; hence
$\psi$, and so $\alpha$, is trivial on these inertia groups, and \ref{pf:vwin:M-exc} gives $\Sha_\ell(\Q,W_\alpha)=0$.

\emph{The main case.} Otherwise, by \ref{pf:vwin:M-zeval} applied with $a_0=1$ we obtain $m$ and $F=f'\ell^m$ with $\mathrm{Dat}(f',m,F,\alpha)$. Choose
$\delta$ acting on $\mu_F$ by $-1$; then $\alpha(\delta)\equiv-1\bmod\ell^m$, so $\alpha(\delta)-1$ is a unit ($\ell$ odd). By \ref{pf:vwin:M-prim}, the
hypotheses of \ref{pf:vwin:M-sha} hold with $F$ replaced by $f\ell$. Let $a_0\ne0$. By \ref{pf:vwin:M-zeval} there are $m,b,\sigma$ with
$\mathrm{Dat}(f',m,f'\ell^m,\alpha)$, $a_0z_b\mid\ell^m$ and $z_b\mid a$, and \cref{pf:vwin:count} gives
$\len\mathrm{T}(\alpha,\delta,f'\ell^m,a_0)\le v(z_b)\le v(a)$. Since $f'\ell^m$ and $f\ell$ have the same prime divisors, $\mathrm{T}(\alpha,\delta,
f'\ell^m,a_0)=\mathrm{T}(\alpha,\delta,f\ell,a_0)$. By \ref{pf:vwin:M-sha}, $\len\Sha_\ell(\Q,W_\alpha)\le v(a)$.
\end{proof}

\begin{proof}[Proof of \cref{pf:vwin:psibound}(i)]
By \cref{pf:vwin:rec} the conclusion of that lemma holds, so \cref{pf:vwin:vchar} applies. In the situation of $\sfname{PsiBound}$,
$\psi=\theta\circ\rho$ has order dividing $[E:\Q]$, which is prime to $\ell$, and $\psi(-1)=(-1)^{j+1}$. By \ref{pf:vwin:M-eig} the length of the
$\theta$-eigenspace equals $\len\Sha_\ell(\Q,W_\alpha)$ with $\alpha=\theta^{-1}\vwincyc^{-j}=(\vwincyc^j\psi)^{-1}$, where $\theta$ is viewed as a
character of $G_\Q$ through $\rho$; \cref{pf:vwin:vchar} bounds it by $v(a)$. This proves $\sfname{PsiBound}$.
\end{proof}

\begin{remark}\label{pf:vwin:rem-minus}
The argument is the classical one, for $j=0$, of Rubin \cite{Rubin1990}, transposed to Tate twists; global duality enters only through
\cref{pf:vwin:rec}, and no Iwasawa theory is used. The bound has the strength of the main conjecture \cite{MazurWiles1984} for each component
with $\psi(-1)=(-1)^{j+1}$, up to the explicit factor $N\cdot\mathrm{primorial}(j+2)$ in $\alpha_{j+1}$.
\end{remark}

\subsubsection{Summation: the minus part of the window bound}

\begin{definition}\label{pf:vwin:pbsb}
For $C\in\R$ let $\mathrm{SB}(C)$ be the statement: for every CM field $E$ embedding into $\Q(\zeta_N)$, every odd prime $\ell$ (more generally every
prime $\ell$ which is odd or such that $\vwincyc\equiv1\bmod4$ on $G_E$), every $\ell$-datum $D$ of $E$ and every $j\ge1$, $\vwinmin(D,j)<\infty$ and
\[
 \vwinmin(D,j)^\natural\log\ell\le C\,\ell\,[E:\Q]^2(j+1)\log(j\ell N)^2 .
\]
Let $\mathrm{PB}$ be the statement: for every CM field $E$ embedding into $\Q(\zeta_N)$, every choice of $p$-data and every $j\ge1$,
$\vwinmin(D_p,j)<\infty$ for odd $p\nmid[E:\Q]$, and for every finite set $S$ of odd primes not dividing $[E:\Q]$,
\[
\textstyle\sum_{\ell\in S}\log\ell\cdot\vwinmin(D_\ell,j)^\natural\le[E:\Q]\log\bigl(16N^{j+2}(j+1)!\bigr).
\]
\formallabel{UpperTheory.V.Odd.StickBound}
\formallabel{UpperTheory.V.Odd.PrimeBound}
\end{definition}

\begin{lemma}\label{pf:vwin:pb}
\begin{enumerate}[label=(\alph*)]
\item $\sfname{PsiBound}$ implies $\mathrm{PB}$.
\item There is $C\ge0$ with $\mathrm{SB}(C)$.
\item If $(m_M;q,d)$ is a prime pair, $E=E_{M,k,q,d}$ and $\ell$ is an odd prime dividing $[E:\Q]$, then $\ell\mid\deg M$.
\end{enumerate}
\formallabel{UpperTheory.V.Odd.primeBound_of}
\formallabel{UpperTheory.V.Odd.stickBound}
\formallabel{UpperTheory.V.Odd.deg}
\end{lemma}

\begin{proof}
(a) Let $\ell\nmid2[E:\Q]$ and choose $\OO$ containing the values of all characters of $\Gal(E/\Q)$, and $\rho$ as in $\sfname{PsiBound}$
(\lref{UpperTheory.V.Odd.actEx}, \lref{UpperTheory.V.Odd.actFacts}). The minus part of $\Sha_D(j)$ decomposes into the
$\theta$-eigenspaces for the characters $\theta$ of $\Gal(E/\Q)$ with $\theta(-1)=(-1)^{j+1}$ (at most $[E:\Q]$ of them), so by
$\sfname{PsiBound}$, $\vwinmin(D_\ell,j)\le\sum_\theta v_\ell(\alpha_{j+1}(\theta\circ\rho))$, with $v_\ell$ normalized by $v_\ell(\ell)=1$
(\lref{UpperTheory.V.Odd.atPrime}). By the non-vanishing of $B_{k,\chi}$ for $\chi(-1)=(-1)^k$ and the integrality and size of the products of
the normalized Bernoulli numbers $\alpha_{j+1}$ over Galois orbits of characters (\lref{UpperTheory.V.Bern.b_nzC},
\lref{UpperTheory.V.Odd.prodVal}), there are integers $w(\ell)$ with $\sum_\theta v_\ell(\alpha_{j+1}(\theta\circ\rho))\le w(\ell)$ for all
odd $\ell\nmid[E:\Q]$ and $\sum_{\ell\in S}w(\ell)\log\ell\le[E:\Q]\log(16N^{j+2}(j+1)!)$ for every finite set $S$ of primes. This gives
$\mathrm{PB}$.

(b) By an annihilation argument with higher Stickelberger determinants in the manner of Coates--Sinnott \cite{CoatesSinnott1974}, there is $C$
such that for every such $(E,N,\ell,D,j)$ some power $\ell^r$ with $r\log\ell\le C[E:\Q](j+1)\log(j\ell N)$ kills the minus part
(\cref{pf:vwin:killpow} below); and the $\ell$-torsion of $H^1(G_{E,S_\ell},\vwinQZ(-j))$ has length at most
$g\le8\ell[E:\Q]\log(\ell N)$ (\lref{UpperTheory.V.Odd.gen}). A submodule killed by $\ell^r$ of a module whose $\ell$-torsion has length $\le g$
has length $\le rg$ (\lref{UpperTheory.V.Odd.lenPow}). Hence $\vwinmin(D,j)\log\ell\le r\log\ell\cdot g\le8C\ell[E:\Q]^2(j+1)\log(j\ell N)^2$,
using $\log(\ell N)\le\log(j\ell N)$; so $\mathrm{SB}(8C)$ holds.

(c) $E$ is the fixed field of the subgroup of $(\Z/N)^\times$ defined in \cref{pf:vwin:fields}, whose index divides a power of $2$ times
$\deg M$ (the conditions modulo $2^{k+2}$ and the two Jacobi symbols have $2$-power indices); so an odd prime dividing $[E:\Q]$ divides
$\deg M$.
\end{proof}

\begin{lemma}[Killing the minus part]\label{pf:vwin:killpow}
There is $C\ge0$ such that for every CM field $E$ embedding into $\Q(\zeta_N)$ ($N\ge1$), every prime $\ell$, every algebraic closure $\Omega$
and $\ell$-datum $D$ of $E$ with $E$ procyclic at $\ell$, and every $e\ge1$, there is $r\in\N$ with $r\log\ell\le C[E:\Q](e+1)\log(e\ell N)$ such
that $\ell^r$ kills the minus part of $\Sha_D(e)$.
\formallabel{UpperTheory.V.Odd.killPow}
\end{lemma}

\begin{proof}
Put $M=N\ell$ and $a=2M+1$, which is odd, prime to $M$ and $>1$. The group $(\Z/M)^\times$ acts on $E$ through a surjection $\phi$ onto
$\Gamma=\Gal(E/\Q)$ (\lref{UpperTheory.V.Odd.actEx}), and $c=\phi(-1)$ is complex conjugation, $c\ne1$. The minus part is the part on which
$c$ acts by $-1$.
\emph{(1) Annihilation.} There is $z\in\Zl[\Gamma]$, an integral representative of the higher Stickelberger element of level $M$ attached to
$a$ and the twist $e+1$ (of Coates--Sinnott type \cite{CoatesSinnott1974}), which annihilates $\Sha_D(e)$
(\lref{UpperTheory.V.Odd.zKill}; this is proved by Stickelberger's theorem and Kummer theory in the layers $E(\mu_{\ell^n})$, which is
where procyclicity at $\ell$ is used).
\emph{(2) Determinant.} If $z$ kills $x$ and $cx=-x$, then $\det{}_c(z)\,x=0$, where $\det_c(z)\in\Zl$ is the determinant of multiplication by
$z$ on the free $\Zl$-module $\Zl[\Gamma]/(c+1)$; indeed $\Zl[\Gamma]$ acts on $x$ through this quotient, and the adjugate matrix gives
$\mathrm{adj}(z)z=\det_c(z)$ (\lref{UpperTheory.V.Odd.adjKill}).
\emph{(3) Size of the determinant.} Since $\Gamma$ is abelian, $\det_c(z)$ is the product of the values $\chi(z)$ over the odd characters $\chi$ of
$\Gamma$. These are non-zero (as $\phi(-1)\ne1$; they are generalized Bernoulli numbers up to explicit factors) and bounded explicitly in terms of
$M$, $a$ and $e$. Hence $\det_c(z)\ne0$, and its $\ell$-adic valuation, together with the $\ell$-part of $\varphi(M)$ which enters the
normalization of $z$, gives $r$ such that $\ell^r$ kills the minus part, with $r\log\ell$ at most $[E:\Q]$ times the logarithm of that bound plus
$\log(N\ell)$ (\lref{UpperTheory.V.Odd.detVal}).
\emph{(4) Conclusion.} An elementary estimate of the bound in (3) gives $r\log\ell\le C[E:\Q](e+1)\log(e\ell N)$
(\lref{UpperTheory.V.Odd.bdxLog}).\end{proof}

\begin{proof}[Proof of \cref{pf:vwin:psibound}(ii)]
Let $C\ge0$ with $\mathrm{SB}(C)$; we show $\sfname{ThmBcFree}$ with this $C$. Let $M$, $k,t\ge2$, a prime pair $(m_M;q,d)$ and a CM field $E=E_{M,k,q,d}$ be
given; put $n=[E:\Q]$ and $F=\cond(E)\ge1$ (so $E$ embeds into $\Q(\zeta_F)$). Finiteness of $\vwinmin(D_p,j)$ for odd $p$ follows from $\mathrm{SB}$ if
$p\mid n$ and from $\mathrm{PB}$ otherwise. Let $S$ be a finite set of odd primes, and split it according as $\ell\nmid n$ or $\ell\mid n$.

For $\ell\nmid n$: by $\mathrm{PB}$, for each $0<j<t$ the contribution is at most $n\log(16F^{j+2}(j+1)!)\le n(2t+5)\log(tF)$
(\lref{UpperTheory.V.Odd.log_bnd_le_mul_log}; here $j+1\le t$ and $t\ge2$); summing over the $t-1$ values of $j$ and using
$(t-1)(2t+5)\le3t^2$ gives at most $3nt^2\log(tF)$.

For $\ell\mid n$: then $\ell\mid\deg M$ by \cref{pf:vwin:pb}(c), and by $\mathrm{SB}(C)$ the contribution of $\ell$ is at most
$\sum_{j=1}^{t-1}C\ell n^2(j+1)\log(j\ell F)^2\le C\ell n^2t^2\log(t\ell F)^2$. Summing over the odd primes dividing $\deg M$ gives the second term of
$\beta(C,\deg M,n,F,t)$.
\end{proof}

\begin{remark}\label{pf:vwin:rem-SB}
The bound $\mathrm{SB}(C)$ holds for every odd prime, not only for those dividing $[E:\Q]$, and it does not use $\sfname{PsiBound}$. It is not used for
$\ell\nmid[E:\Q]$ because its sum over all primes $\ell$ would diverge; there the main-conjecture strength of $\sfname{PsiBound}$, summed over $\ell$ by
the product formula for $\Pi_j$, is needed.
\end{remark}


% ---------------------------------------------------------------------------
\subsection{Lemma K}\label{sec:pf:vwin:lemmaK}

For a number field $L$ let $r_2(L)$ be its number of complex places, $d_L$ its discriminant and $h_L$ its class number. For $C\in\R$ and
$e,n,D\in\N$ put $\kappa(C,e,n,D)=\exp\bigl((2en)^C\log(2D)\bigr)$, so that $\sfname{LemmaK}(C_K)$ (\cref{def:window:lemmaK}) asserts: for every
prime $\ell$, every totally imaginary number field $L$ containing a primitive $q_\ell$-th root of unity, every algebraic closure of $L$, every
$\ell$-datum $D$ of $L$ and every $e\ge1$, the group $H^2(G_{L,S_\ell},\Zl(e+1))$ is finite and
$\log\#H^2(G_{L,S_\ell},\Zl(e+1))\le\kappa(C_K,e,[L:\Q],|d_L|)$.

\begin{theorem}[\cref{thm:window:lemmaK}]\label{pf:vwin:lemmaK}
$\sfname{LiteratureRankNHypothesis}$ implies that $\sfname{LemmaK}(C_K)$ holds for some $C_K\in\R$.
\formallabel{UpperTheory.V.ZK.lemmaK_litN_alt}
\formallabel{UpperTheory.V.ZK.lemmaK_of_litN_alt}
\formallabel{UpperTheory.V.ZK.lemmaK}
\formallabel{UpperTheory.V.ZK.lemmaK_of_steps}
\end{theorem}

Finiteness is Soul\'e's theorem \cite{Soule1979}; the point of Lemma K is a bound which is explicit and uniform in $\ell$. The proof is a
quantitative version of Soul\'e's method (see also \cite{DwyerFriedlander1985}): classes in the homology of $\GL_N(\OO_L)$ are attached to
elements of the cyclotomic layers $F=L(\mu_{\ell^\nu})$ whose divisors satisfy an eigenvalue relation under $\Gal(F/L)$, characteristic
classes evaluated at residue fields detect them, and a rank bound of Borel type for the rational homology of $\GL_N(\OO_L)$ together with a
bound for its torsion limits their number; this bounds the eigenspaces of class groups which control $\Sha$, and Poitou--Tate duality passes
from $\Sha$ to $H^2$. The formal proof is organized as a deduction from twenty-one intermediate statements and a rank bound. We state the rank
bound and the intermediate statements and give the deduction; then we prove the statements which carry the arithmetic (the cycle and
the relation of \ref{pf:vwin:K2}, $\mathrm S_{\mathrm{detect}}$, and \ref{pf:vwin:K5}--\ref{pf:vwin:K7}), and the basic properties
$\mathrm S_{\mathrm{chain}}$, $\mathrm S_{\mathrm{prism}}$ of the bar complex. The remaining combinatorial statements
($\mathrm S^\times_{\mathrm{prism}}$, $\mathrm S_\eta$, \ref{pf:vwin:K3}) and the Steinitz statement are stated precisely and cited.

\subsubsection{The rank bound}

For a ring $R$ and $n,d\ge1$ let $\psi_d:\GL_n(R)\to\GL_{dn}(R)$, $A\mapsto1_d\otimes A$ (block diagonal with $d$ copies of $A$), and for
$n\le n'$ let $\mathrm{st}:\GL_n(R)\to\GL_{n'}(R)$ be the standard embedding. For $\ell\ge1$, $m\ge0$ and $x\in H_m(\GL_n(R),\Z)$ define
$z_t(x)\in H_m(\GL_{\ell^tn}(R),\Z)$ by $z_0(x)=x$ and $z_{t+1}(x)=\psi_\ell(z_t(x))-\ell^{t+2}\,\mathrm{st}(z_t(x))$.

\begin{theorem}[A rank bound of Borel type]\label{pf:vwin:bodd}
Let $L$ be a totally imaginary number field, $k$ even with $k+1<r_2(L)$, $\ell\ge1$ and $n\ge2(k+1)+4$. Then the subgroup of
$H_{k+1}(\GL_{\ell^kn}(\OO_L),\Z)$ generated by the elements $z_k(x)$, $x\in H_{k+1}(\GL_n(\OO_L),\Z)$, has rank at most $r_2(L)$.
\formallabel{UpperTheory.V.ZK.bOdd}
\formallabel{UpperTheory.V.ZK.BOdd}
\end{theorem}

This is a statement of the shape of Borel's computation of the rational homology of $\GL_n(\OO_L)$ in the stable range \cite{Borel1974}:
the operations $\psi_\ell-\ell^{t+2}\mathrm{st}$ annihilate the decomposable part, and the primitive part in degree $2m-1$ ($m\ge2$) has
rank $r_2(L)$ for $L$ totally imaginary. The development contains a formal proof of about $10{,}600$ lines
(\lref{PositivityTheory.KB.bOdd_zk}), which we use as a black box.

\subsubsection{The intermediate statements}

We describe the statements; in all of them $\ell$ is a prime and $q$ a positive integer.
\begin{enumerate}[label=(K\arabic*)]
\item\label{pf:vwin:K1} (Chains.) For a commutative ring $A$ and a group $G$ let $C_m(G;A)$ be the free $A$-module on $G^m$ with the
inhomogeneous bar boundary $\partial$, functorial in $G$, and for $k\in G$ let $P_k:C_m\to C_{m+1}$ be the prism operator of conjugation by $k$.
The statements $\mathrm{S}_{\mathrm{chain}}$, $\mathrm{S}_{\mathrm{prism}}$, $\mathrm{S}_{\mathrm{prism}}^{\times}$ assert: $\partial^2=0$,
functoriality, $P$ commutes with homomorphisms, the homology of the trivial group vanishes in positive degree; the homotopy formula
$\partial P_k+P_k\partial=\mathrm{conj}_{k^{-1}}-1$; and $P_{kk'}-P_k-P_{k'}\circ\mathrm{conj}_{k^{-1}}$ maps cycles to boundaries.
(\cref{pf:vwin:schain}, \lref{UpperTheory.V.ZK.sPrismMul}.)
\item\label{pf:vwin:K2} (An explicit cycle.) $\mathrm{S}_{zc}$: an explicit chain $z_{q,e}\in C_{2e+2}(\Z/q;\Z/q)$ on the cyclic group of order $q$ is a
cycle; $\mathrm{S}_{\mathrm{rel}}$: for $q=\ell^\nu$ some combination $\sum_tn_t\,(\mathrm{id},t)_*z_{q,e}$, $t<2e+4$, with one coefficient a unit, is a
boundary in $C_\ast((\Z/q)^2;\Z/q)$; $\mathrm{S}_\eta$: for $M<\ell^c$, $c\le\nu$ and $t_0\le M$ there is an element of the group ring
$\Z[(\Z/\ell^\nu)^2]$ whose image under $(x_1,x_2)\mapsto x_1+tx_2$ vanishes for $t\le M$, $t\ne t_0$, and equals
$\sum_{i\le M}a_i[(1+\ell^c)^i]$ for $t=t_0$, where $\prod_{j<M}(X-(1+\ell^c)^j)=\sum_ia_iX^i$. (\cref{pf:vwin:zc},
\cref{pf:vwin:srel}, \lref{UpperTheory.V.ZK.sEta}.)
\item\label{pf:vwin:K3} (Characteristic classes over finite fields.) $\mathrm{S}_{\mathrm{trans}}$: for a finite field $k$ and a ring $A$ in which
$\#k$ is $1$, every linear form on $C_{m+1}(k^\times;A)$ vanishing on boundaries extends to an \emph{additive family}: a family of linear forms
$\theta_V$ on $C_{m+1}(\GL(V);A)$, for all finite-dimensional $k$-vector spaces $V$, vanishing on boundaries, invariant under isomorphisms and
additive for direct sums on cycles. $\mathrm{S}_{\mathrm{coc}}$: for a primitive $q$-th root of unity $\zeta\in k$ and a homomorphism $v:k^\times\to\Z/q$
there is a linear form on $C_{2e+3}(k^\times;\Z/q)$, vanishing on boundaries, with explicit values $x^{e+1}v(g)$ on the prisms of the images of
$z_{q,e}$. $\mathrm{S}_\theta$ (deduced from the previous statements): there is an additive family $\theta$ of degree $2e+3$ with
$\theta_V\bigl(P_g(\chi_x)_*z_{q,e}\bigr)=x^{e+1}v(\det g)$ for all $V$, $g\in\GL(V)$, $x\in\N$, where $\chi_x$ is the scalar character
$c\mapsto\zeta^{xc}$. $\mathrm{S}_\phi$: additive families pulled back along reduction maps $\GL_n(R)\to\GL_n(k)$ form primitive families
(compatible with stabilization, multiplied by $d$ under $\psi_d$, vanishing on boundaries). (\lref{UpperTheory.V.ZK.sTrans},
\lref{UpperTheory.V.ZK.sCoc}, \lref{UpperTheory.V.ZK.sThetaOf}, \lref{UpperTheory.V.ZK.sPhi}.)
\item\label{pf:vwin:K4} (Arithmetic.) $\mathrm{S}_{\mathrm{Stein}}$: given divisors $a_j,b_j$ ($j\in J$) of a number field $F$ and $\beta\in F^\times$ with
$\mathrm{div}(\beta)=\sum_j(b_j-a_j)$ (and $J\ne\emptyset$ or $\beta=1$), there is $g\in\GL_J(F)$ of determinant $\beta$ carrying
$\bigoplus_j\mathfrak a_j$ onto $\bigoplus_j\mathfrak b_j$, where $\mathfrak a_j=\{x:v_P(x)\ge a_j(P)\ \forall P\}$ and similarly $\mathfrak b_j$
(Steinitz). $\mathrm{S}_{\mathrm{detect}}$: if $\zeta_{\ell^\nu}\in F$ and $\beta\in F^\times$ is an $\ell^\nu$-th power in
the residue fields of almost all primes of degree one over a subfield $L$, then $\beta$ is an $\ell^\nu$-th power in $F$.
(\lref{LateralTheory.Fin.SStein.intersectionNumber_baseChange}, \cref{pf:vwin:sdetect}.)
\item\label{pf:vwin:K5} (Classes from divisors.) Let $F/L$ be a \emph{layer} with data $(\ell,c,\nu,\zeta,\gamma)$: $F/L$ is Galois with group generated by
$\gamma$ of order $\ell^{\nu-c}$, $\zeta\in\OO_F^\times$ is a primitive $\ell^\nu$-th root of unity, $\gamma\zeta=\zeta^{1+\ell^c}$, $1\le c\le\nu$, $c\ge2$ if
$\ell=2$. The divisor relation for $a\in\N$ ($\ord_P(\alpha)=\mathfrak d(\gamma P)-a\,\mathfrak d(P)$ for all $P$) and the twisted
norm $\mathrm{N}_a(\alpha)=\prod_{j<\ord\gamma}\gamma^j(\alpha)^{a^j}$ are those of \cref{pf:vwin:laydef}. $\mathrm{S}_{\mathrm{class}}$ (deduced from \ref{pf:vwin:K1},
\ref{pf:vwin:K2}, $\mathrm{S}_{\mathrm{Stein}}$): if $2e+3<\ell^c$ and $\alpha$ satisfies the divisor relation for $a=(1+\ell^c)^{e+1}$, there are a cycle
$y\in C_{2e+3}(\GL_N(\OO_L);\Z/\ell^\nu)$ and $0\ne b\in\OO_F$ such that for every residue map $r:\OO_F\to k$ to a finite field which is surjective on
$\OO_L$ and with $r(b)\ne0$, and every additive family $\theta$ normalized as in $\mathrm{S}_\theta$, $\theta(r_*y)=p_e\cdot v(u)$, where $u$ is the residue
of $\mathrm{N}_a(\alpha)$ and $p_e=\prod_{j<2e+3,\,j\ne e+1}((1+\ell^c)^{e+1}-(1+\ell^c)^j)$. $\mathrm{S}_{\mathrm{upper}}$ (deduced from
\ref{pf:vwin:K1} and \cref{pf:vwin:bodd}): for $L$ totally imaginary, $k$ even with $k+1<r_2(L)$ and a primitive family $\Phi$ with values in a
$\Z/\ell^\nu$-module, the set of values $\ell^k\Phi(y)$ on cycles $y$ has at most $(\ell^\nu)^{r_2(L)}\,t(L,2k+6,k+1)\,t(L,2k+6,k)$ elements, where
$t(L,N,m)$ is the order of the torsion of $H_m(\GL_N(\OO_L),\Z)$. (\cref{pf:vwin:sclass}, \cref{pf:vwin:supper}.)
\item\label{pf:vwin:K6} (Counting.) $\mathrm{S}_{\mathrm{few}}$ (deduced from $\mathrm{S}_\theta$, $\mathrm{S}_\phi$, $\mathrm{S}_{\mathrm{class}}$,
$\mathrm{S}_{\mathrm{upper}}$, $\mathrm{S}_{\mathrm{detect}}$): for $L$ totally imaginary and a layer $F/L$ as above with $2e+3<\ell^c$ and
$2e+3<r_2(L)$, every set of elements $\alpha$ satisfying the divisor relation for $a=(1+\ell^c)^{e+1}$, pairwise inequivalent in the sense that
$\mathrm{N}_a(\alpha)^s$ and $\mathrm{N}_a(\alpha')^s$ do not differ by an $\ell^\nu$-th power ($s=\ell^{2e+2+(2e+3)(c+2e+3)}$), has at most
$(\ell^\nu)^{r_2(L)}t(L,4e+10,2e+3)\,t(L,4e+10,2e+2)$ elements. $\mathrm{S}_{\mathrm{count}}$: if such a bound $B$ holds (with exponent $s=\ell^b$), then the
group of homomorphisms $\varphi$ from the divisors of $F$ to $\Z/\ell^\nu$ vanishing on principal divisors and with $\varphi\circ\gamma^{-1}=a\varphi$ is finite,
of order at most $B\,\ell^{(b+v_\ell(E))g_L+c+v_\ell(E)}/(\ell^\nu)^{r_2(L)}$, where $a=(1+\ell^c)^E$ and $g_L=2[L:\Q]+\log_2h_L+2$.
$\mathrm{S}_\Sha$: if $\mu_{\ell^c}\subset L$, $\mu_{\ell^{c+1}}\not\subset L$ and these groups have order at most $B$ for all layers $F/L$, then
$\ell^{\vwinsh(L,\ell,E)}\le B$. $\mathrm{S}_{\mathrm{layer}}$: for $L\ni\zeta_{q_\ell}$ and $E\ge1$ there is a totally imaginary extension $L'$ of degree
$d\le([L:\Q]+1)(4E+4)$ with $\mu_{\ell^c}\subset L'$ maximal, $2E+1<\ell^c\le2[L':\Q]$, $2E+1<r_2(L')$, $|d_{L'}|\le|d_L|^d(\ell^c)^{[L':\Q]}$ and
$\vwinsh(L,\ell,E)\le\vwinsh(L',\ell,E)$. $\mathrm{S}_{\mathrm{tors}}$: there is $C_0$ with $\log t(L,4E+6,m)\le\kappa(C_0,E,[L:\Q],|d_L|)$ for all
totally imaginary $L$, $E\ge1$ and $m\le2E+1$. (\cref{pf:vwin:sfew}, \cref{pf:vwin:scount}, \cref{pf:vwin:ssha}, \cref{pf:vwin:slayer}, \cref{pf:vwin:stors}.)
\item\label{pf:vwin:K7} (Numerics.) $\mathrm{S}_{\mathrm{num}}$: the counting statements imply that there is $C$ such that for every prime $\ell$, every
totally imaginary $L\ni\zeta_{q_\ell}$ and $E\ge1$, $\vwinsh(L,\ell,E)$ is finite and $\vwinsh(L,\ell,E)\log\ell\le\kappa(C,E,[L:\Q],|d_L|)$; we denote this
statement by $\Sha\mathrm{K}(C)$. (\cref{pf:vwin:snum}.)
\end{enumerate}

\begin{proof}[Proof of \cref{pf:vwin:lemmaK}]
Poitou--Tate duality is obtained from $\sfname{LiteratureRankNHypothesis}$ as in the proof of \cref{pf:vwin:offour}, and the local finiteness and
duality statements are members of it. All intermediate statements \ref{pf:vwin:K1}--\ref{pf:vwin:K7} are theorems of the development, and
\cref{pf:vwin:bodd} holds. Combine them in the order indicated: \ref{pf:vwin:K1}--\ref{pf:vwin:K3} give $\mathrm{S}_\theta$; \ref{pf:vwin:K1},
\ref{pf:vwin:K2} and $\mathrm{S}_{\mathrm{Stein}}$ give $\mathrm{S}_{\mathrm{class}}$; \ref{pf:vwin:K1} and \cref{pf:vwin:bodd} give
$\mathrm{S}_{\mathrm{upper}}$; these with $\mathrm{S}_\phi$ and $\mathrm{S}_{\mathrm{detect}}$ give $\mathrm{S}_{\mathrm{few}}$; with
$\mathrm{S}_{\mathrm{count}}$, $\mathrm{S}_\Sha$, $\mathrm{S}_{\mathrm{layer}}$ and $\mathrm{S}_{\mathrm{tors}}$, \cref{pf:vwin:snum} gives
$\Sha\mathrm{K}(C)$ for some $C$. Finally \cref{pf:vwin:sfinal} gives $\sfname{LemmaK}(C_K)$ for some $C_K$.
\end{proof}

\begin{remark}\label{pf:vwin:rem-K}
(1) Several of the intermediate statements are supplied by theorems whose formal names refer to intersection numbers and base change; these
names do not reflect their content. (2) The combinatorial statements \ref{pf:vwin:K1}--\ref{pf:vwin:K3} concern the bar complex of a group and
are independent of arithmetic; the arithmetic enters through $\mathrm{S}_{\mathrm{Stein}}$, $\mathrm{S}_{\mathrm{detect}}$, the rank bound and the
counting statements.
\end{remark}

\subsubsection{The numerical step}

\begin{proposition}\label{pf:vwin:snum}
The statements $\mathrm{S}_{\mathrm{few}}$, $\mathrm{S}_{\mathrm{count}}$, $\mathrm{S}_\Sha$, $\mathrm{S}_{\mathrm{layer}}$ and
$\mathrm{S}_{\mathrm{tors}}$ imply $\Sha\mathrm{K}(C)$ with $C=5(\max(C_0,9)+2)+11$, where $C_0$ is the constant of $\mathrm{S}_{\mathrm{tors}}$.
\formallabel{UpperTheory.V.ZK.sNum}
\formallabel{UpperTheory.V.ZK.SNum}
\end{proposition}

\begin{proof}
Let $\ell$, $L$ (totally imaginary, containing $\zeta_{q_\ell}$) and $E=k+1\ge1$ be given, and let $L'$, $c$, $d$ be as in $\mathrm{S}_{\mathrm{layer}}$;
put $n=[L:\Q]$, $n'=[L':\Q]=dn$. Let $F/L'$ be a layer with data $(\ell,c,\nu,\zeta,\gamma)$. Since $2k+3=2E+1<\ell^c$ and $2E+1<r_2(L')$,
$\mathrm{S}_{\mathrm{few}}$ (with $e=k$) gives the hypothesis of $\mathrm{S}_{\mathrm{count}}$ with $B=(\ell^\nu)^{r_2(L')}T_1T_2$, where
$T_1=t(L',4k+10,2k+3)$, $T_2=t(L',4k+10,2k+2)$, and exponent $\ell^b$, $b=2k+2+(2k+3)(c+2k+3)$; so the group of $\mathrm{S}_{\mathrm{count}}$ for
$a=(1+\ell^c)^E$ has order at most $T_1T_2\,\ell^X$ with $X=(b+v_\ell(E))g_{L'}+c+v_\ell(E)$. By $\mathrm{S}_\Sha$ (applicable since $\mu_{\ell^c}\subset
L'$ is maximal), $\ell^{s'}\le T_1T_2\ell^X$ for $s'=\vwinsh(L',\ell,E)$; and $\vwinsh(L,\ell,E)\le s'$ by $\mathrm{S}_{\mathrm{layer}}$, so
$\vwinsh(L,\ell,E)$ is finite and
\[
 \vwinsh(L,\ell,E)\log\ell\le\log T_1+\log T_2+X\log\ell .
\]
Put $u=2En'$ and $Y=\log(2|d_{L'}|)\ge\frac12$. By $\mathrm{S}_{\mathrm{tors}}$, $\log T_i\le\exp(u^{C_0}Y)$. The elementary bounds
$g_{L'}\le2n'+4n'Y+2$ (from Minkowski's bound for the class number), $c\le\ell^c\le2n'$,
$v_\ell(E)<E$, $\log\ell\le u$ and $n'\ge2E+4$ give $X\log\ell\le u^9Y$, hence
$\vwinsh(L,\ell,E)\log\ell\le2\exp(u^{C_0}Y)+u^9Y\le\exp(u^{\max(C_0,9)+2}Y)=\kappa(\max(C_0,9)+2,E,n',|d_{L'}|)$. Finally
$n'=dn$ with $d\le(n+1)(4E+4)$, $|d_{L'}|\le|d_L|^d(\ell^c)^{n'}$ and $\ell^c\le2n'$ bound this by $\kappa(C,E,n,|d_L|)$; these are routine
estimates.
\end{proof}

\subsubsection{From \texorpdfstring{$\Sha$}{Sha} to \texorpdfstring{$H^2$}{H2}}

\begin{proposition}\label{pf:vwin:sfinal}
Assume Poitou--Tate duality and the local finiteness and duality statements over finite extensions of $\Ql$. If $\Sha\mathrm{K}(C)$ holds, then
$\sfname{LemmaK}(C_K)$ holds for some $C_K>0$.
\formallabel{UpperTheory.V.ZK.sFinal}
\formallabel{UpperTheory.V.ZK.SFinal}
\formallabel{UpperTheory.V.ZK.ShaK}
\end{proposition}

\begin{proof}
Let $C_K>0$ be such that for all $e,n,D\ge1$
\begin{equation}\label{pf:vwin:numfin}
 4n\log(2D)\cdot2e(2e+3)\cdot\log(n+1)+\kappa(C,e,n,D)\le\kappa(C_K,e,n,D);
\end{equation}
such $C_K$ exists because $\kappa(C_K,e,n,D)=(2D)^{(2en)^{C_K}}$ grows faster than any fixed polynomial in $e,n,\log D$ times $\kappa(C,e,n,D)$
when $C_K$ is large (\lref{UpperTheory.V.K.numFin}). Let $\ell$, $L$, $D$, $e$ be as in $\sfname{LemmaK}$, $n=[L:\Q]$, and let $s=\vwinsh(L,\ell,e)$,
which is finite with $s\log\ell\le\kappa(C,e,n,|d_L|)$ by $\Sha\mathrm{K}(C)$. By Poitou--Tate duality (\lref{UpperTheory.V.LDataCardLeFinrank.VLen_le_shaLen_add};
$L$ is totally imaginary), $\vwinVL(D,e)$ is at most $\vwinsh(D,e)$ plus the sum over the places $v\mid\ell$ of the lengths of the local groups
$H^0(L_v,\vwinQZ(-e))$, and this sum is at most $n(2+\log_\ell n+e)$ (a local computation: the length of $H^0(L_v,\vwinQZ(-e))$ is bounded in terms of
$[L_v:\Ql]$, $\log_\ell[L_v:\Ql]$ and $e$). Hence $H^2=H^2(G_{L,S_\ell},\Zl(e+1))$ has finite length $k\le s+n(2+\log_\ell n+e)$, so it is finite
with $\#H^2=\ell^k$, and
\[
 \log\#H^2=k\log\ell\le s\log\ell+n(2+e)\log\ell+n\log_\ell n\cdot\log\ell .
\]
Since $\mu_\ell\subset L$, $\ell\le n+1$ (\lref{UpperTheory.V.K.primeLe}), and $\ell^{\log_\ell n}\le n$; so the last two terms are at most
$n(3+e)\log(n+1)\le4n\log(2|d_L|)\cdot2e(2e+3)\log(n+1)$, using $\log(2|d_L|)\ge\log2>\frac12$. By \eqref{pf:vwin:numfin},
$\log\#H^2\le\kappa(C_K,e,n,|d_L|)$.
\end{proof}


\subsubsection{Layers and twisted norms}

We now prove the intermediate statements which carry the arithmetic of Lemma K. Throughout, $\ell$ is a prime.

\begin{definition}[Layers, divisor relations, twisted norms]\label{pf:vwin:laydef}
Let $L$ be a number field. $L$ has \emph{top root} $\ell^c$ if $c\ge1$ ($c\ge2$ if $\ell=2$), $\mu_{\ell^c}\subset L$ and
$\mu_{\ell^{c+1}}\not\subset L$. Let $F/L$ be an extension of number fields. $F/L$ is a \emph{layer with data
$(\ell,c,\nu,\zeta,\gamma)$} if $F/L$ is Galois, $\gamma$ generates $\Gal(F/L)$ and has order $\ell^{\nu-c}$, $\zeta\in\OO_F^\times$
is a primitive $\ell^\nu$-th root of unity with $\gamma\zeta=\zeta^{1+\ell^c}$, $1\le c\le\nu$, and $c\ge2$ if $\ell=2$. For
$\gamma\in\Gal(F/L)$ and $a\in\N$:
\begin{enumerate}[label=(\alph*)]
\item $\alpha\in F^\times$ satisfies the \emph{divisor relation} for $(\gamma,a)$ with respect to a divisor $\mathfrak d$ of $F$ if
$\ord_P(\alpha)=\mathfrak d(\gamma P)-a\,\mathfrak d(P)$ for every prime $P$ of $\OO_F$;
\item the \emph{twisted norm} is $\mathrm N_a(\alpha)=\prod_{j<\ord\gamma}\gamma^j(\alpha)^{a^j}$;
\item for $q\ge1$, $\mathrm{Eq}(\gamma,a,q)$ is the group of functions $\varphi$ from the primes of $\OO_F$ to $\Z/q$ such that
$\sum_P\ord_P(\alpha)\varphi(P)=0$ for all $\alpha\in F^\times$ and $\varphi(\gamma^{-1}P)=a\varphi(P)$ for all $P$;
\item for $s,q,B\in\N$, $\mathrm{Few}(\gamma,a,s,q,B)$ is the statement: every finite set $S$ of elements of $F^\times$ which
satisfy the divisor relation for $(\gamma,a)$ (with respect to some divisor), and such that $\mathrm N_a(\alpha)^s=\mathrm
N_a(\alpha')^s\delta^q$ with $\delta\in F^\times$ and $\alpha,\alpha'\in S$ implies $\alpha=\alpha'$, has at most $B$ elements.
\end{enumerate}
Finally, $g_L=2[L:\Q]+\log_2h_L+2$, where $h_L$ is the class number of $L$.
\formallabel{UpperTheory.V.ZK.TopRoot}
\formallabel{UpperTheory.V.ZK.Lay}
\formallabel{UpperTheory.V.ZK.DivDel}
\formallabel{UpperTheory.V.ZK.nrm}
\formallabel{UpperTheory.V.ZK.EqHom}
\formallabel{UpperTheory.V.ZK.Few}
\formallabel{UpperTheory.V.ZK.gens}
\end{definition}

So $\mathrm{Eq}(\gamma,a,q)$ is the group of homomorphisms $\varphi:\mathrm{Cl}(F)\to\Z/q$ with $\varphi\circ\gamma^{-1}=a\varphi$,
and the statements $\mathrm S_{\mathrm{few}}$, $\mathrm S_{\mathrm{count}}$, $\mathrm S_\Sha$ of \ref{pf:vwin:K6} are
formulated with these notions.

\subsubsection{The bar complex}

For a commutative ring $A$ and a group $G$, $C_m(G;A)$ is the free $A$-module on the symbols $[g_1|\cdots|g_m]$, $g_i\in G$, with the
boundary $\partial[g_1|\cdots|g_m]=[g_2|\cdots|g_m]+\sum_{i=1}^{m-1}(-1)^i[\cdots|g_ig_{i+1}|\cdots]+(-1)^m[g_1|\cdots|g_{m-1}]$, and a
homomorphism $f$ acts by $f_*[g_1|\cdots|g_m]=[f(g_1)|\cdots|f(g_m)]$. For $k\in G$ the prism operator is
\[
 P_k[g_1|\cdots|g_m]=\sum_{i=0}^m(-1)^i[g_1|\cdots|g_i|k|k^{-1}g_{i+1}k|\cdots|k^{-1}g_mk].
\]

\begin{lemma}[$\mathrm S_{\mathrm{chain}}$, $\mathrm S_{\mathrm{prism}}$]\label{pf:vwin:schain}
\begin{enumerate}[label=(\alph*)]
\item $\partial^2=0$; $f_*\partial=\partial f_*$; $(gf)_*=g_*f_*$; $f_*P_k=P_{f(k)}f_*$.
\item If $G$ is trivial, every cycle of positive degree is a boundary.
\item $\partial P_k+P_k\partial=(c_{k^{-1}})_*-1$ on $C_m(G;A)$ for $m\ge1$, where $c_{k^{-1}}(g)=k^{-1}gk$.
\end{enumerate}
\formallabel{UpperTheory.V.ZK.sChain}
\formallabel{UpperTheory.V.ZK.SChain}
\formallabel{UpperTheory.V.ZK.sPrism}
\formallabel{UpperTheory.V.ZK.SPrism}
\formallabel{UpperTheory.V.ZK.prism}
\end{lemma}

\begin{proof}
(a) $\partial$ is the alternating sum of the face maps $d_0,\dots,d_m$ of the bar construction, which satisfy the simplicial identities
$d_id_j=d_{j-1}d_i$ for $i<j$; hence $\partial^2=0$. The other identities are immediate from the formulas, since homomorphisms commute with
products and conjugation. (b) $C_m(1;A)=A$, and $\partial:C_{m+1}\to C_m$ is multiplication by $\sum_{i=0}^{m+1}(-1)^i$, i.e.\ by $1$ if $m$ is
odd and by $0$ if $m$ is even ($m\ge0$). So in degree $m+1\ge1$ the cycles vanish if $m$ is odd, and are all of $A=\partial C_{m+2}$ if $m$ is even.
(c) Pass to the homogeneous bar complex, $[g_1|\cdots|g_m]\mapsto(1,g_1,g_1g_2,\dots,g_1\cdots g_m)$ modulo the diagonal action of $G$. There
$P_k$ becomes the usual prism operator $(x_0,\dots,x_m)\mapsto\sum_i(-1)^i(x_0,\dots,x_i,x_ik,\dots,x_mk)$ between the identity and right
multiplication by $k$, which induces $(c_{k^{-1}})_*$ on the inhomogeneous complex; the standard computation of the prism boundary gives the
formula.
\end{proof}
\subsubsection{The relation in the homology of \texorpdfstring{$(\Z/q)^2$}{(Z/q)2}}

For $q\ge1$ and $e\in\N$ let $z_{q,e}\in C_{2e+2}(\Z/q;\Z/q)$ be the chain $\sum_{j\in(\Z/q)^{e+1}}[j_0|1|j_1|1|\cdots|j_e|1]$
(the group $\Z/q$ written additively), and for $t\in\N$ let $d_t:\Z/q\to(\Z/q)^2$, $x\mapsto(x,tx)$.

\begin{lemma}[$\mathrm S_{zc}$]\label{pf:vwin:zc}
$z_{q,e}$ is a cycle.
\formallabel{UpperTheory.V.ZK.sZc}
\formallabel{UpperTheory.V.ZK.SZc}
\formallabel{UpperTheory.V.ZK.zc}
\formallabel{UpperTheory.V.ZK.dt}
\formallabel{LateralTheory.Q.SZcT.SZc_of}
\end{lemma}

\begin{proof}
The boundary of $[g_1|\cdots|g_n]$ is $[g_2|\cdots|g_n]+\sum_{i=1}^{n-1}(-1)^i[\cdots|g_i+g_{i+1}|\cdots]+(-1)^n[g_1|\cdots|g_{n-1}]$.
Sum over $j$, with $n=2e+2$. The first faces do not depend on $j_0$ and add up to $q$ times a chain, which is $0$ with
coefficients in $\Z/q$. For $k<e$, the face merging $j_k$ with the following entry $1$ (sign $-1$) and the face merging that
entry $1$ with $j_{k+1}$ (sign $+1$) become the same sum after the substitutions $j_k\mapsto j_k-1$, resp.\
$j_{k+1}\mapsto j_{k+1}-1$, and cancel. The face merging $j_e$ with the last entry (sign $-1$) becomes, after
$j_e\mapsto j_e-1$, the last face (sign $+1$), and they cancel.
\end{proof}

\begin{proposition}[$\mathrm S_{\mathrm{rel}}$]\label{pf:vwin:srel}
For $q=\ell^\nu$ and $e\in\N$ there are $n_0,\dots,n_{2e+3}\in\Z/q$ and $t_0\le2e+3$ such that $n_{t_0}$ is a unit and
$\sum_{t=0}^{2e+3}n_t\,(d_t)_*z_{q,e}$ is a boundary in $C_*((\Z/q)^2;\Z/q)$.
\formallabel{LateralTheory.Fin.SRel.intersectionNumber_baseChange}
\formallabel{UpperTheory.V.ZK.SRel}
\formallabel{LateralTheory.Fin.SRel.relation_statement_holds}
\formallabel{LateralTheory.Q.SRelT.SRel_of}
\formallabel{LateralTheory.Q.SRelT.rel_bd}
\end{proposition}

\begin{proof}
Let $H=H_{2e+2}((\Z/q)^2;\Z/q)$. By \cref{pf:vwin:zc} and functoriality of the bar complex, the chains $(d_t)_*z_{q,e}$ are
cycles; let $x_t\in H$ be their classes. The $\Z/q$-module $H$ is generated by $2e+3$ elements
(\lref{LateralTheory.Q.SRelT.gen_two}: by the K\"unneth formula, $H_n((\Z/q)^2;\Z/q)$ is free of rank $n+1$ over $\Z/q$, since
$H_i(\Z/q;\Z/q)\cong\Z/q$ for all $i$; the formal proof is an induction on $e$ using the projection $(\Z/q)^2\to\Z/q$).
Consider $\Phi:(\Z/q)^{2e+4}\to H$, $(n_t)\mapsto\sum_tn_tx_t$, with image $N$ and kernel $K$. If $K\subset\ell(\Z/q)^{2e+4}$,
then $N/\ell N\cong(\Z/q)^{2e+4}/(K+\ell(\Z/q)^{2e+4})$ has dimension $2e+4$ over $\F_\ell$. But a finite $\Z/q$-module is a direct
sum of cyclic modules, the number of which is $\dim_{\F_\ell}$ of its $\ell$-torsion and also of its reduction modulo $\ell$; so
$\dim N/\ell N=\dim N[\ell]\le\dim H[\ell]\le2e+3$, a contradiction. Hence some $(n_t)\in K$ has a coordinate $n_{t_0}$ not
divisible by $\ell$, i.e.\ a unit, and $\sum_tn_tx_t=0$ means that $\sum_tn_t(d_t)_*z_{q,e}$ is a boundary.
\end{proof}

\subsubsection{Classes from divisors}

\begin{proposition}[$\mathrm S_{\mathrm{class}}$]\label{pf:vwin:sclass}
Assume $\mathrm S_{\mathrm{chain}}$, $\mathrm S_{\mathrm{prism}}$, $\mathrm S_{\mathrm{prism}}^\times$, $\mathrm S_{zc}$,
$\mathrm S_{\mathrm{rel}}$, $\mathrm S_\eta$ and $\mathrm S_{\mathrm{Stein}}$. Let $F/L$ be a layer with data
$(\ell,c,\nu,\zeta,\gamma)$, $q=\ell^\nu$, $e\in\N$ with $2e+3<\ell^c$, $a=(1+\ell^c)^{e+1}$, and let $\alpha\in F^\times$ satisfy the
divisor relation for $(\gamma,a)$ with respect to $\mathfrak d$. Then there are $N\ge1$, a cycle $y\in C_{2e+3}(\GL_N(\OO_L);\Z/q)$
and $0\ne b\in\OO_F$ with the following property. Let $k$ be a finite field and $r:\OO_F\to k$ a ring homomorphism which is
surjective on $\OO_L$, with $r(b)\ne0$; let $v:k^\times\to\Z/q$ be a homomorphism and $\theta$ an additive family of degree
$2e+3$ over $k$ with coefficients $\Z/q$, normalized as in $\mathrm S_\theta$ with respect to $r(\zeta)$ and $v$. Then the
residue $u=r(\mathrm N_a(\alpha))\in k^\times$ is defined and
\[
 \theta\bigl(r_*y\bigr)=p_e\,v(u),\qquad p_e=\prod_{j<2e+3,\ j\ne e+1}\bigl((1+\ell^c)^{e+1}-(1+\ell^c)^j\bigr),
\]
where $r_*y$ is the image of $y$ in $C_{2e+3}(\GL_N(k);\Z/q)$.
\formallabel{UpperTheory.V.ZK.sClassOf}
\formallabel{UpperTheory.V.ZK.SClassOf}
\formallabel{UpperTheory.V.ZK.SClass}
\formallabel{UpperTheory.V.ZK.Cl.cJoin}
\formallabel{UpperTheory.V.ZK.Cl.CJoin}
\formallabel{UpperTheory.V.ZK.pe}
\end{proposition}

The idea is that of Soul\'e: the divisor relation makes $\alpha$ the ``determinant'' of an automorphism of a lattice in a vector
space with an action of $\mu_q$, and a characteristic class evaluated on the prism of this automorphism picks up the residue of
$\alpha$.

\begin{proof}
Put $J=(\Z/q)^2$, $\chi_t(x_1,x_2)=x_1+tx_2$ and $\kappa:\Gal(F/L)\to\Z/q$ with $\sigma\zeta=\zeta^{\kappa(\sigma)}$, so
$\kappa(\gamma^j)=(1+\ell^c)^j$. Consider the $L$-vector space $F^J$ and its group $\mathcal G$ of $L$-linear automorphisms; for a family
$\mathfrak a=(\mathfrak a_j)_{j\in J}$ of divisors of $F$ let $\Lambda(\mathfrak a)=\{x\in F^J:\ \ord_P(x_j)\ge\mathfrak a_j(P)\ \forall
j,P\}$, so $\Lambda(0)=\OO_F^J$. Let $\rho:J\to\mathcal G$, $\rho(x)=\mathrm{diag}(\zeta^{x_1j_1+x_2j_2})_j$; then $\rho\circ d_t$ is the
diagonal action $c\mapsto\mathrm{diag}(\zeta^{c\chi_t(j)})$ of $\Z/q$, and $\rho$ preserves every $\Lambda(\mathfrak a)$.
We say that $g\in\mathcal G$ \emph{carries} $(\chi,\mathfrak b,\delta)$, for $\chi:J\to\Z/q$, a family $\mathfrak b$ and
$\delta:\Z/q\to F^\times$, if $g$ maps $\Lambda(0)$ onto $\Lambda(\mathfrak b)$ and has a block form: it maps each coordinate line
$Fe_j$ semilinearly (for some $\sigma_j\in\Gal(F/L)$ with $\sigma_j(\zeta)^{\chi(j)}=\zeta^{\beta(j)}$) into the block $\bigoplus_{\beta(j')=
\beta(j)}Fe_{j'}$, followed by $F$-linear automorphisms $g'_x$ of the blocks with $\det g'_x=\delta(x)$. Such $g$ commutes with
the diagonal action of $\Z/q$ through $\chi$ (\lref{UpperTheory.V.ZK.Cl.cRep}).

\emph{Step 1: the data.} By $\mathrm S_{\mathrm{rel}}$ choose $n_t$ ($t\le2e+3$) and $t_0$ with $n_{t_0}$ a unit and
$\sum_tn_t(d_t)_*z_{q,e}$ a boundary. By $\mathrm S_\eta$ (with $M=2e+3<\ell^c$) choose $\eta:J\to\Z$ with
$\chi_{t*}\eta=0$ for $t\ne t_0$ and $\chi_{t_0*}\eta=\sum_{i\le M}a_i[(1+\ell^c)^i]$, where
$\prod_{j<M}(X-(1+\ell^c)^j)=\sum_ia_iX^i$. Let $P(X)=\prod_{j<M,\,j\ne e+1}(X-(1+\ell^c)^j)=\sum_ip_iX^i$ and
$\beta=\prod_i\gamma^{-i}(\alpha)^{p_i}$. Since $(X-a)P(X)=\sum_ia_iX^i$, the divisor relation gives
$\ord_P(\beta)=\sum_ia_i\,\mathfrak d(\gamma^iP)$ (\lref{UpperTheory.V.ZK.Cl.cBetaDiv}).

\emph{Step 2: the automorphisms.} For every $t\le2e+3$ there is $g_t\in\mathcal G$ carrying $(\chi_t,(\eta_j\mathfrak d)_j,\delta_t)$,
where $\delta_t=1$ for $t\ne t_0$, and $\delta_{t_0}(1)=\beta$, $\delta_{t_0}(x)=1$ for $x\ne1$
(\lref{UpperTheory.V.ZK.Cl.cGtNe}, \lref{UpperTheory.V.ZK.Cl.cGtEq}). Indeed, on the block of $\chi_t$-weight $x$
the lattice $\Lambda((\eta_j\mathfrak d)_j)$ is a sum of fractional ideals whose total divisor is $(\chi_{t*}\eta)(x)\mathfrak d$
(after transport of the blocks of weights $(1+\ell^c)^ix$ to weight $x$ by $\gamma^i$, which changes $\mathfrak d$ into
$\mathfrak d\circ\gamma^i$); for $t\ne t_0$ these divisors vanish, and for $t=t_0$ the only non-trivial block has total divisor
$\sum_ia_i\,\mathfrak d\circ\gamma^i=\mathrm{div}(\beta)$. In both cases $\mathrm S_{\mathrm{Stein}}$ provides an $F$-linear
automorphism of the block of determinant $\delta_t(x)$ carrying $\OO_F$-lattices as required.

\emph{Step 3: the cycle.} Choose $0\ne b_0\in\OO_L$ such that every prime of $L$ not dividing $b_0$ splits completely in $F$ in
the sense needed for reduction (\lref{UpperTheory.V.ZK.Cl.cSplit}), and a multiple $0\ne b\in\OO_L$ of $b_0$ lying in every
prime in the support of $\mathfrak d$ (\lref{UpperTheory.V.ZK.Cl.cSupp}). Let $\mathcal G_b\subset\mathcal G$ be the
stabilizer of $\OO_F[1/b]^J$ and $\mathcal H\subset\mathcal G_b$ that of $\Lambda(0)$. As the lattices $\Lambda((\eta_j\mathfrak
d)_j)$ and $\Lambda(0)$ agree away from $b$, $g_t\in\mathcal G_b$ (\lref{UpperTheory.V.ZK.Cl.cLatLoc}); $\rho$ takes values in
$\mathcal H$; $g_t$ commutes with $\rho(d_t(x))$; and $g_{t_0}^{-1}\rho(x)g_{t_0}$ and $g_t^{-1}g_{t_0}$ lie in $\mathcal H$, since
they preserve $\Lambda(0)$. The chain lemma (\lref{UpperTheory.V.ZK.Cl.cChain}, a consequence of
$\mathrm S_{\mathrm{chain}}$, $\mathrm S_{\mathrm{prism}}$ and $\mathrm S^\times_{\mathrm{prism}}$) then gives a cycle
$y_0\in C_{2e+3}(\mathcal H;\Z/q)$ whose image in $C_{2e+3}(\mathcal G_b;\Z/q)$ is homologous to
$\sum_tn_tP_{g_t}\bigl((\rho\circ d_t)_*z_{q,e}\bigr)$, where $P_g$ is the prism operator of $g$: this sum is a cycle up to the
boundary $\sum_tn_t(d_t)_*z_{q,e}$ transported by $\rho$, and the conditions on $g_t$ allow to move it into $\mathcal H$.
Finally $\mathcal H$ acts $\OO_L$-linearly on the finite projective $\OO_L$-module $\OO_F^J$
(\lref{UpperTheory.V.ZK.Cl.cInt}); adding a complement $P$ with $\OO_F^J\oplus P\cong\OO_L^N$, on which $\mathcal H$ acts
trivially, gives $\lambda:\mathcal H\to\GL_N(\OO_L)$ (\lref{UpperTheory.V.ZK.Cl.cFree}). Put $y=n_{t_0}^{-1}\lambda_*y_0$.

\emph{Step 4: evaluation.} Let $r$, $k$, $v$, $\theta$ be as in the statement. Since $r(b)\ne0$, the prime of $L$ below $r$
splits completely in $F$, so $\OO_F\otimes_{\OO_L}k\cong k^{\Gal(F/L)}$, and reduction defines $\mu:\mathcal G_b\to
\GL(k^{\Gal(F/L)\times J})$, compatible with $\lambda$ up to the trivial summand (\lref{UpperTheory.V.ZK.Cl.cMu}). Let $U$
be the group of $b$-units of $F$ and $f:U\to\Z/q$, $f(x)=v(r(x))$; it is additive and $\Gal(F/L)$-stable
(\lref{UpperTheory.V.ZK.Cl.cUnit}), and $\alpha$, $\beta$, $\mathrm N_a(\alpha)\in U$ because $b$ lies in the primes of the support
of $\mathfrak d$. By additivity and invariance of $\theta$ the trivial summand does not contribute
(\lref{UpperTheory.V.ZK.Cl.cTriv}); boundaries are killed, so
$\theta(r_*\lambda_*y_0)=\sum_tn_t\,\theta\bigl(\mu_*P_{g_t}((\rho\circ d_t)_*z_{q,e})\bigr)$. The value of $\theta$ on the prism of a
block automorphism is computed from the normalization of $\theta$ ($\theta_V(P_g(\chi_x)_*z_{q,e})=x^{e+1}v(\det g)$) and additivity
over the $\Gal(F/L)\times J$-decomposition (\lref{UpperTheory.V.ZK.Cl.cVal}):
\[
 \theta\bigl(\mu_*P_{g_t}((\rho\circ d_t)_*z_{q,e})\bigr)=\sum_{\sigma\in\Gal(F/L)}\sum_{x\in\Z/q}(\kappa(\sigma)x)^{e+1}f\bigl(\sigma(\delta_t(x))\bigr).
\]
For $t\ne t_0$ this vanishes since $f(1)=0$; for $t=t_0$ it equals $\sum_\sigma\kappa(\sigma)^{e+1}f(\sigma\beta)$. With
$\sigma=\gamma^j$, $\kappa(\sigma)^{e+1}\equiv a^j$, and $\sum_ja^jf(\gamma^j\gamma^{-i}\alpha)=a^i\sum_ma^mf(\gamma^m\alpha)$ because
$a^{\ord\gamma}\equiv1\pmod q$; hence $\sum_\sigma\kappa(\sigma)^{e+1}f(\sigma\beta)=P(a)f(\mathrm N_a(\alpha))$
(\lref{UpperTheory.V.ZK.Cl.cBetaVal}). As $P(a)=p_e$ and $f(\mathrm N_a(\alpha))=v(u)$, we get
$\theta(r_*y)=n_{t_0}^{-1}n_{t_0}\,p_e\,v(u)=p_e\,v(u)$.
\end{proof}

\subsubsection{Counting equivariant characters}

\begin{lemma}[Twisted divisor data]\label{pf:vwin:ctalg}
Let $Y$, $I$ be abelian groups, $I$ torsion-free, $T$ the torsion subgroup of $Y$, $d:Y\to I$ a homomorphism, $\sigma\in\End(Y)$ and
$\sigma_I\in\End(I)$ with $d\sigma=\sigma_Id$ and $\sigma^n=1$, $\sigma_I^n=1$ for some $n\ge1$, and $a\ge2$. Put $\delta=\sigma-a$,
$\delta_I=\sigma_I-a$, $\mathrm N=\sum_{j<n}a^j\sigma^{n-j}$, $U=d^{-1}(\delta_II)$, $K=\ker d$, and
\[
 W=U/(\delta Y+T),\quad \mathcal U_\delta=K/(\delta K+T),\quad \mathrm{Cl}_\delta=I/(dY+\delta_II),\quad Y_m=Y/(mY+T)\ (m\ge1).
\]
\begin{enumerate}[label=(\alph*)]
\item\label{pf:vwin:ctalg-a} If $I/dY$ and $\mathcal U_\delta$ are finite, then $\mathrm{Cl}_\delta$ and $W$ are finite and
$\#W=\#\mathcal U_\delta\cdot\#\mathrm{Cl}_\delta$.
\item\label{pf:vwin:ctalg-b} If $a^n-1=Qu$ with $Q\ge1$, $u\in\Z$, then for every $s\in\N$ the map $y\mapsto s\mathrm N(y)$ induces a
homomorphism $f:W\to Y_Q$, and $f(w)=0$ implies $usw=0$.
\item\label{pf:vwin:ctalg-c} If $\ell\ge1$ divides $a-1$, the number of $w\in W$ with $\ell w=0$ is at most the number of classes in
$Y_\ell$ of elements $z\in Y$ with $\sigma z-z\in\ell Y+T$ and $dz\in\ell I$.
\item\label{pf:vwin:ctalg-d} Let $A$ be a finite abelian group, $f:A\to Z$ a homomorphism, $u\in\Z$ prime to $\ell$ and $s,g\in\N$ such
that $f(w)=0$ implies $u\ell^sw=0$, and $\#A[\ell]\le\ell^g$. Then $\ell^{v_\ell(\#A)}\le\#f(A)\cdot\ell^{sg}$.
\end{enumerate}
\formallabel{UpperTheory.V.ZK.Ct.Dat}
\formallabel{UpperTheory.V.ZK.Ct.cardW}
\formallabel{UpperTheory.V.ZK.Ct.kerF}
\formallabel{UpperTheory.V.ZK.Ct.selInj}
\formallabel{UpperTheory.V.ZK.Ct.pCount}
\end{lemma}

\begin{proof}
Two preliminary remarks. First, $\delta_I$ is injective: $\sigma_Ix=ax$ gives $x=\sigma_I^nx=a^nx$, so $x=0$ as $I$ is torsion-free.
Second, $\mathrm N\delta=\delta\mathrm N=(1-a^n)\sigma$ (the sum telescopes, using $\sigma^{n+1}=\sigma$), so $\delta y\in T$ implies
$y\in T$. Note also $\delta Y\subset U$ and $T\subset K\subset U$.

\ref{pf:vwin:ctalg-a} For $u\in U$ let $\psi(u)\in I$ be the element with $\delta_I\psi(u)=du$. The map $u\mapsto\psi(u)\bmod dY$ sends
$U$ onto the kernel of the endomorphism $\bar\delta$ induced by $\delta_I$ on $\mathrm{Cl}=I/dY$ (if $\delta_Ii=du$, then $u\in U$
and $\psi(u)=i$), with kernel $K+\delta Y$ (if $\psi(u)=dy$, then $du=d\delta y$). Moreover $K\cap(\delta Y+T)=\delta K+T$, since
$d(\delta y)=0$ forces $dy=0$. Hence $W$ maps onto $\ker\bar\delta$ with kernel $(K+\delta Y)/(\delta Y+T)\cong\mathcal U_\delta$, and
$\#\ker\bar\delta=\#\mathrm{coker}\,\bar\delta=\#\mathrm{Cl}_\delta$ as $\mathrm{Cl}$ is finite.

\ref{pf:vwin:ctalg-b} $\mathrm N(\delta Y+T)\subset(1-a^n)Y+T\subset QY+T$, so $f$ is defined. If $s\mathrm N(y)=Qz+t$ with $t\in T$,
apply $\delta$: $-Qus\,\sigma y=Q\delta z+\delta t$, so $Q(us\sigma y+\delta z)\in T$, hence $us\sigma y\in\delta Y+T$, and applying
$\sigma^{n-1}=\sigma^{-1}$, which commutes with $\delta$, gives $usy\in\delta Y+T$.

\ref{pf:vwin:ctalg-c} Let $w=[y]$ with $\ell y=\delta z+t$, $t\in T$. Then $\sigma z-z=\ell y-t+(a-1)z\in\ell Y+T$, and
$\delta_Idz=d\delta z=\ell dy=\ell\delta_I\psi(y)$ gives $dz=\ell\psi(y)$. The class of $z$ in $Y_\ell$ is well defined ($z$ is unique
modulo $T$ by the second remark, and replacing $y$ by $y+\delta z'+t'$ replaces $z$ by $z+\ell z'$ modulo $T$), and $w\mapsto[z]$ is
injective: if $z=\ell z''+t''$, then $\ell(y-\delta z'')\in T$, so $y\in\delta Y+T$ and $w=0$.

\ref{pf:vwin:ctalg-d} Let $A_\ell$ be the $\ell$-primary part of $A$, of order $\ell^{v_\ell(\#A)}$. As $u$ acts bijectively on $A_\ell$,
$\ker f\cap A_\ell\subset A[\ell^s]$, which has at most $\#A[\ell]^s\le\ell^{sg}$ elements (multiplication by $\ell^{i-1}$ embeds
$A[\ell^i]/A[\ell^{i-1}]$ into $A[\ell]$). Hence $\#A_\ell=\#f(A_\ell)\cdot\#(\ker f\cap A_\ell)\le\#f(A)\,\ell^{sg}$.
\end{proof}

\begin{lemma}\label{pf:vwin:ctarith}
Let $F/L$ be a layer with data $(\ell,c,\nu,\zeta,\gamma)$, $n=\ell^{\nu-c}$, $E\ge1$ and $a=(1+\ell^c)^E$. Apply the notation of
\cref{pf:vwin:ctalg} to $Y=F^\times$ (written additively), $I$ the group of divisors of $F$, $d=\mathrm{div}$, $\sigma=\gamma^{-1}$
on $F^\times$ and $(\sigma_I\mathfrak d)(P)=\mathfrak d(\gamma P)$.
\begin{enumerate}[label=(\alph*)]
\item\label{pf:vwin:ctarith-a} $a\ge2$, $v_\ell(a-1)=c+v_\ell(E)$ and $v_\ell(a^n-1)=\nu+v_\ell(E)$.
\item\label{pf:vwin:ctarith-b} $I/dY=\mathrm{Cl}(F)$ is finite, and $\#\mathrm{Eq}(\gamma,a,\ell^\nu)\le\ell^{v_\ell(\#\mathrm{Cl}_\delta)}$.
\item\label{pf:vwin:ctarith-c} $U$ is the set of elements satisfying the divisor relation for $(\gamma,a)$, and $\mathrm N=\mathrm N_a$.
\item\label{pf:vwin:ctarith-d} If $\mathrm{Few}(\gamma,a,s,\ell^\nu,B)$ holds and $W$ is finite, then the homomorphism $f$ of
\cref{pf:vwin:ctalg}\ref{pf:vwin:ctalg-b} (for $Q=\ell^\nu$ and $s$) has image of order at most $B$.
\end{enumerate}
\formallabel{UpperTheory.V.ZK.Ct.IsDat}
\formallabel{UpperTheory.V.ZK.Ct.datEx}
\formallabel{UpperTheory.V.ZK.Ct.lte}
\formallabel{UpperTheory.V.ZK.Ct.clFin}
\formallabel{UpperTheory.V.ZK.Ct.eqHomLe}
\formallabel{UpperTheory.V.ZK.Ct.fewLe}
\end{lemma}

\begin{proof}
The hypotheses of \cref{pf:vwin:ctalg} hold: $d\sigma=\sigma_Id$ because $\ord_P(\gamma^{-1}\alpha)=\ord_{\gamma P}(\alpha)$, and
$\gamma^n=1$.
\ref{pf:vwin:ctarith-a} $a\ge1+\ell^c\ge2$. Write $(1+x)^E-1=Ex+\sum_{k\ge2}\binom Ekx^k$ with $x=\ell^c$. Since $k\binom Ek=E\binom{E-1}{k-1}$,
$v_\ell(\binom Ekx^k)\ge v_\ell(E)-v_\ell(k)+kc>c+v_\ell(E)$ for $k\ge2$, because $(k-1)c>v_\ell(k)$ when $c\ge1$ and $\ell$ is odd, or
$c\ge2$. Hence $v_\ell(a-1)=c+v_\ell(E)$, and the same formula with $E\ell^{\nu-c}$ in place of $E$ gives $v_\ell(a^n-1)=\nu+v_\ell(E)$.
\ref{pf:vwin:ctarith-b} The class group is finite. A function $\varphi$ on the primes defines a homomorphism $I\to\Z/\ell^\nu$, which
vanishes on $dY$ iff $\varphi$ vanishes on principal divisors, and $\varphi(\sigma_I\mathfrak d)=\sum_Q\mathfrak d(Q)\varphi(\gamma^{-1}Q)$;
so $\mathrm{Eq}(\gamma,a,\ell^\nu)\cong\Hom(\mathrm{Cl}_\delta,\Z/\ell^\nu)$, whose order is at most that of the $\ell$-primary part of
$\mathrm{Cl}_\delta$.
\ref{pf:vwin:ctarith-c} $\alpha\in U$ means $\mathrm{div}(\alpha)=\sigma_I\mathfrak d-a\mathfrak d$ for some $\mathfrak d$, which is the
divisor relation; and $\sigma^{n-j}=\gamma^j$.
\ref{pf:vwin:ctarith-d} For each element of $f(W)$ choose $y\in U$ with $f([y])$ equal to it. If $y\ne y'$ are chosen for different
elements, then $s\mathrm N_a(y)-s\mathrm N_a(y')\notin\ell^\nu Y+T$, so $\mathrm N_a(y)^s\ne\mathrm N_a(y')^s\delta^{\ell^\nu}$ for all
$\delta$; the chosen elements form a set as in $\mathrm{Few}$, so $\#f(W)\le B$.
\end{proof}

\begin{proposition}[$\mathrm S_{\mathrm{count}}$]\label{pf:vwin:scount}
Let $L$ be a totally imaginary number field, $F/L$ a layer with data $(\ell,c,\nu,\zeta,\gamma)$, $E\ge1$, $b,B\in\N$ and
$a=(1+\ell^c)^E$, and assume $\mathrm{Few}(\gamma,a,\ell^b,\ell^\nu,B)$. Then $\mathrm{Eq}(\gamma,a,\ell^\nu)$ is finite and
\[
 \#\mathrm{Eq}(\gamma,a,\ell^\nu)\cdot(\ell^\nu)^{r_2(L)}\le B\cdot\ell^{(b+v_\ell(E))g_L+c+v_\ell(E)}.
\]
\formallabel{UpperTheory.V.ZK.sCount}
\formallabel{UpperTheory.V.ZK.SCount}
\end{proposition}

\begin{proof}
Use the notation of \cref{pf:vwin:ctarith}; write $v=v_\ell$. Here $K=\OO_F^\times$, $T=\mu(F)$ and $\mathcal U_\delta=\OO_F^\times/(\mu(F)\cdot
(\gamma^{-1}-a)\OO_F^\times)$. Since $L$ is totally imaginary and $F/L$ is cyclic of degree $n$, $\mathcal U_\delta$ is finite and
$\#\mathcal U_\delta\cdot(a-1)=(a^n-1)^{r_2(L)}$ (\lref{UpperTheory.V.ZK.Ct.cardEQ}: by Herbrand's theorem
$(\OO_F^\times\otimes\Q)\oplus\Q\cong\Q[\Gal(F/L)]^{r_2(L)}$, so the characteristic polynomial of $\gamma^{-1}$ on the lattice
$\OO_F^\times/\mu(F)$ is $(X^n-1)^{r_2(L)}/(X-1)$, and the cokernel of the injective endomorphism $\gamma^{-1}-a$ has order
$|\det(\gamma^{-1}-a)|=(a^n-1)^{r_2(L)}/(a-1)$). By \cref{pf:vwin:ctalg}\ref{pf:vwin:ctalg-a} and \cref{pf:vwin:ctarith}\ref{pf:vwin:ctarith-b},
$W$ is finite with $\#W=\#\mathcal U_\delta\cdot\#\mathrm{Cl}_\delta$, and $\mathrm{Eq}(\gamma,a,\ell^\nu)$ is finite of order at most $\ell^{v(\#\mathrm{Cl}_\delta)}$.
By \cref{pf:vwin:ctarith}\ref{pf:vwin:ctarith-a}, $v(\#\mathcal U_\delta)+c+v(E)=r_2(L)(\nu+v(E))\ge r_2(L)\nu$. Hence
\[
 \#\mathrm{Eq}(\gamma,a,\ell^\nu)\,(\ell^\nu)^{r_2(L)}\le\ell^{v(\#\mathrm{Cl}_\delta)+v(\#\mathcal U_\delta)+c+v(E)}=\ell^{v(\#W)}\,\ell^{c+v(E)}.
\]
Write $a^n-1=\ell^{\nu+v(E)}m$ with $\ell\nmid m$. By \cref{pf:vwin:ctalg}\ref{pf:vwin:ctalg-b} with $Q=\ell^\nu$, $u=\ell^{v(E)}m$ and
$s=\ell^b$ there is $f:W\to Y_{\ell^\nu}$ with $f(w)=0\Rightarrow m\ell^{b+v(E)}w=0$, and $\#f(W)\le B$ by
\cref{pf:vwin:ctarith}\ref{pf:vwin:ctarith-d}. Since $\ell\mid a-1$, \cref{pf:vwin:ctalg}\ref{pf:vwin:ctalg-c} bounds $\#W[\ell]$ by the
number of $\gamma$-invariant $\ell$-Selmer classes of $F$, which is at most $\#\mathrm{Sel}_\ell(L)\cdot\ell^2$
(\lref{UpperTheory.V.ZK.Ct.selInvLe}: such a class is represented by $z$ with $\gamma z=z\,w^\ell\rho^r$, $\rho$ a generator of
$\mu(F)$, $r<\ell$; two representatives with the same $r$ and the same norm of $w$ differ by a Selmer class of $L$, and these data take
at most $\ell^2$ values), where $\mathrm{Sel}_\ell(L)\subset L^\times/L^{\times\ell}$ is the group of classes of elements whose divisor
is divisible by $\ell$; and $\#\mathrm{Sel}_\ell(L)\le\ell^{2[L:\Q]+\log_2h_L}$ (\lref{UpperTheory.V.ZK.Ct.selLCard}: it is an
extension of a subgroup of $\mathrm{Cl}(L)[\ell]$ by $\OO_L^\times/\OO_L^{\times\ell}$). So $\#W[\ell]\le\ell^{g_L}$, and
\cref{pf:vwin:ctalg}\ref{pf:vwin:ctalg-d} (with $s=b+v(E)$, $u=m$, $g=g_L$) gives $\ell^{v(\#W)}\le B\,\ell^{(b+v(E))g_L}$.
Combining the two displayed bounds proves the claim.
\end{proof}

\subsubsection{From equivariant characters to \texorpdfstring{$\Sha$}{Sha}}

\begin{proposition}[$\mathrm S_\Sha$]\label{pf:vwin:ssha}
Let $L$ be a totally imaginary number field with top root $\ell^c$, $E\ge1$ and $B\in\N$. Assume that for every layer $F/L$ with data
$(\ell,c,\nu,\zeta,\gamma)$ the group $\mathrm{Eq}(\gamma,(1+\ell^c)^E,\ell^\nu)$ is finite of order at most $B$. Then
$\vwinsh(L,\ell,E)$ is finite and $\ell^{\vwinsh(L,\ell,E)}\le B$.
\formallabel{PositivityTheory.QuaternionShaBounds.sSha}
\formallabel{UpperTheory.V.ZK.SSha}
\formallabel{PositivityTheory.QuaternionShaBounds.ShaAsm_holds}
\formallabel{PositivityTheory.QuaternionShaBounds.ShaAsm}
\formallabel{PositivityTheory.QuaternionShaBounds.LayEx_holds}
\formallabel{PositivityTheory.QuaternionShaBounds.LayEx}
\formallabel{PositivityTheory.QuaternionShaBounds.ClInj_holds}
\formallabel{PositivityTheory.QuaternionShaBounds.ClInj}
\formallabel{PositivityTheory.QuaternionShaBounds.ShaAssumptions.card_le}
\formallabel{PositivityTheory.QuaternionShaBounds.ShaAssumptions.good_of}
\formallabel{PositivityTheory.QuaternionShaBounds.ShaAssumptions.Good}
\end{proposition}

\begin{proof}
Fix an algebraic closure $\Omega$ of $L$, an $\ell$-datum $D$, $G_L=\Gal(\Omega/L)$, $U_\infty=\ker\vwincyc$, and for $\nu\ge1$ let
$F_\nu=L(\mu_{\ell^\nu})\subset\Omega$ and $G(\nu)=\Gal(\Omega/F_\nu)$. As $\mu_{q_\ell}\subset L$, $L$ is procyclic at $\ell$, and by
\cref{pf:vwin-he:resinj,pf:vwin-he:resimage} (with $e=E\ne0$) restriction to $U_\infty$ is an injective $\Zl$-linear map from
$\Sha_D(E)$ onto the module $\mathcal W$ of locally constant homomorphisms $\varphi:U_\infty\to\vwinQZ$ vanishing on the inertia groups
and on the decomposition groups above $\ell$, with $\varphi(gug^{-1})=\vwincyc(g)^{-E}\varphi(u)$ for $g\in G_L$. It suffices to show that
every finite subset $T\subset\mathcal W$ has at most $B$ elements: then $\mathcal W$ is finite of order at most $B$, and a finite
$\Zl$-module of order $\ell^k$ has length $k$ (\lref{UpperTheory.V.Kum.lenOfCard}), so $\vwinsh(L,\ell,E)=\len_{\Zl}\mathcal W$ is
finite and $\ell^{\vwinsh(L,\ell,E)}=\#\mathcal W\le B$.

\emph{(1) Layers.} Let $\nu\ge c$. As $\vwincyc(G_L)=1+\ell^c\Zl$ (\cref{pf:vwin-he:layers}), there is $g\in G_L$ with
$\vwincyc(g)=1+\ell^c$. Let $\gamma=g|_{F_\nu}$ and $\zeta\in F_\nu$ a primitive $\ell^\nu$-th root of unity; then
$\gamma\zeta=\zeta^{1+\ell^c}$, $\gamma$ has the order $\ell^{\nu-c}$ of $1+\ell^c$ in $(\Z/\ell^\nu)^\times$ (as $c\ge1$, and $c\ge2$
if $\ell=2$), and $[F_\nu:L]\le\ell^{\nu-c}$ since $\zeta^{\ell^{\nu-c}}\in\mu_{\ell^c}\subset L$. So $\gamma$ generates
$\Gal(F_\nu/L)$, and $F_\nu/L$ is a layer with data $(\ell,c,\nu,\zeta,\gamma)$.

\emph{(2) Class groups.} Let $\mathcal C_\nu$ be the group of homomorphisms $\theta:\mathrm{Cl}(F_\nu)\to\vwinQZ$ with
$\theta(\tau x)=\vwincyc(h)^{-E}\theta(x)$ whenever $\tau\in\Gal(F_\nu/L)$ is the restriction of $h\in G_L$. For $\theta\in\mathcal C_\nu$
killed by $\ell^\nu$ let $\varphi_\theta(P)\in\Z/\ell^\nu$ correspond to $\theta([P])\in\ell^{-\nu}\Zl/\Zl$. Then
$\varphi_\theta\in\mathrm{Eq}(\gamma,(1+\ell^c)^E,\ell^\nu)$: it vanishes on principal divisors, and taking $h=g^{-1}$,
$\varphi_\theta(\gamma^{-1}P)=\vwincyc(g)^E\varphi_\theta(P)=(1+\ell^c)^E\varphi_\theta(P)$ in $\Z/\ell^\nu$. The map $\theta\mapsto\varphi_\theta$ is
injective, since the classes of primes generate $\mathrm{Cl}(F_\nu)$. Moreover, for $\nu\ge2$, by class field theory the locally
constant homomorphisms $\psi:G(\nu)\to\vwinQZ$ vanishing on all inertia groups with $\psi(huh^{-1})=\vwincyc(h)^{-E}\psi(u)$ ($h\in G_L$)
are exactly the compositions $\theta\circ\mathrm{Art}$ of the Artin map $G(\nu)\to\mathrm{Cl}(F_\nu)$ with $\theta\in\mathcal C_\nu$
(\lref{UpperTheory.V.Kum.pullEq}); as the Artin map is surjective, $\ell^\nu\psi=0$ iff $\ell^\nu\theta=0$.

\emph{(3) Finite level.} For every $\varphi\in\mathcal W$ there is $n_\varphi$ such that for all $\nu\ge n_\varphi$, $\varphi$ is the restriction to
$U_\infty$ of a homomorphism $\psi$ on $G(\nu)$ as in (2), vanishing moreover on the decomposition groups above $\ell$, with
$\ell^\nu\psi=0$: an equivariant homomorphism on $U_\infty$ extends equivariantly to some $G(n)$
(\lref{UpperTheory.V.Kum.eqExt}); by (2) the extension factors through the finite group $\mathrm{Cl}(F_{\max(n,2)})$, hence is killed
by some $\ell^r$, and its restrictions to $G(\nu)$, $\nu\ge\max(n,2,r)$, are as required.

\emph{(4) Conclusion.} Let $T\subset\mathcal W$ be finite and $\nu\ge\max(c,2)$ with $\nu\ge n_\varphi$ for $\varphi\in T$. For $\varphi\in T$ choose
$\psi_\varphi$ by (3) and $\theta_\varphi\in\mathcal C_\nu$ with $\psi_\varphi=\theta_\varphi\circ\mathrm{Art}$ and $\ell^\nu\theta_\varphi=0$ by (2). The
map $\varphi\mapsto\varphi_{\theta_\varphi}$ is injective ($\varphi$ is the restriction of $\psi_\varphi$, which is determined by $\theta_\varphi$),
so $\#T\le\#\mathrm{Eq}(\gamma,(1+\ell^c)^E,\ell^\nu)\le B$ by the hypothesis for the layer of (1).
\end{proof}


\subsubsection{The upper bound from the rank bound}

Recall from \ref{pf:vwin:K3} that a \emph{primitive family} with values in a $\Z/q$-module $B$ is a family of linear maps
$\Phi_n:C_{k+1}(\GL_n(\OO_L);\Z/q)\to B$ vanishing on boundaries, with $\Phi_{n'}(\mathrm{st}_*y)=\Phi_n(y)$ and
$\Phi_{dn}(\psi_{d*}y)=d\,\Phi_n(y)$ for cycles $y$. For a number field $L$ let $t(L,N,m)$ be the order of the torsion subgroup
of $H_m(\GL_N(\OO_L),\Z)$.

\begin{proposition}[$\mathrm S_{\mathrm{upper}}$]\label{pf:vwin:supper}
Assume $\mathrm S_{\mathrm{chain}}$ and \cref{pf:vwin:bodd}. Let $L$ be totally imaginary, $k$ even with $k+1<r_2(L)$, $q=\ell^\nu$, and
$\Phi$ a primitive family of degree $k+1$ with values in a $\Z/q$-module $B$. Then every finite set of elements of the form
$\ell^k\Phi_n(y)$, $y$ a cycle, has at most $q^{r_2(L)}\,t(L,2k+6,k+1)\,t(L,2k+6,k)$ elements.
\formallabel{UpperTheory.V.ZK.sUpperOf}
\formallabel{UpperTheory.V.ZK.SUpperOf}
\formallabel{UpperTheory.V.ZK.SUpper}
\formallabel{UpperTheory.V.ZK.IsPrimFam}
\formallabel{UpperTheory.V.ZK.tc}
\formallabel{UpperTheory.V.ZK.Up.zVal}
\formallabel{UpperTheory.V.ZK.Up.ZVal}
\formallabel{UpperTheory.V.ZK.Up.bock}
\formallabel{UpperTheory.V.ZK.Up.Bock}
\formallabel{UpperTheory.V.ZK.Up.fibCount}
\formallabel{UpperTheory.V.ZK.Up.FibCount}
\formallabel{UpperTheory.V.ZK.Up.rankCount}
\formallabel{UpperTheory.V.ZK.Up.RankCount}
\formallabel{UpperTheory.V.ZK.Up.IsPrimH}
\end{proposition}

\begin{proof}
\emph{(1) Descent to integral homology.} For an integral cycle $z$ let $\bar z$ be its reduction modulo $q$; then
$\varphi_n([z])=\Phi_n(\bar z)$ defines homomorphisms $\varphi_n:H_{k+1}(\GL_n(\OO_L),\Z)\to B$ (as $\Phi_n$ kills boundaries), compatible
with stabilization and multiplied by $d$ under $\psi_d$.

\emph{(2) Values on the elements $z_t$.} Let $z_t(x)$ be as in \cref{pf:vwin:bodd}. Since $\varphi(\psi_\ell y)=\ell\varphi(y)$ and
$\varphi(\mathrm{st}\,y)=\varphi(y)$, $\varphi(z_{t+1}(x))=\ell(1-\ell^{t+1})\varphi(z_t(x))$, so $\varphi(z_t(x))=\ell^tw_t\varphi(x)$ with
$w_t=\prod_{i<t}(1-\ell^{i+1})\equiv1\pmod\ell$. With $u\in\Z$ inverse to $w_t$ modulo $q$, $\ell^t\varphi_n(x)=\varphi_{\ell^tn}(u\,z_t(x))$
for all $x$.

\emph{(3) Bockstein.} For a group $G$ and a cycle $y\in C_{k+1}(G;\Z/q)$ let $\tilde y$ be its lift with coefficients in
$\{0,\dots,q-1\}$; then $\partial\tilde y=qw$ with $w$ an integral cycle, and $\beta(y)=[w]$ is a torsion class in $H_k(G,\Z)$ ($q[w]=0$).
$\beta$ is additive on cycles, and if $\beta(y)=0$, say $w=\partial v$, then $\tilde y-qv$ is an integral cycle reducing to $y$.

\emph{(4) Counting.} Two elementary facts: (i) if $\beta:Y\to T$ and $g:Y\to B$ are homomorphisms, $T$ finite, and every set of values
of $g$ on $\ker\beta$ has at most $c$ elements, then every set of values of $g$ has at most $c\,\#T$ elements (on each fibre of $\beta$
the values of $g$ form a translate of a subset of $g(\ker\beta)$); (ii) if $H$ is an abelian group with finite torsion, $J\subset H$ a
subgroup of rank at most $r$ and $f:H\to B$ a homomorphism into a group killed by $q$, then $\#f(J)\le\#(J/qJ)\le q^r\,\#H_{\mathrm{tors}}$.

\emph{(5) Conclusion.} Let $S$ be the finite set, and stabilize the cycles $y_s$ ($s\in S$) into one $\GL_n(\OO_L)$ with $n\ge2k+6$. By
homological stability for the torsion (\lref{UpperTheory.V.ZK.Up.torsStab}: for $n'\ge n\ge2m+4$ the torsion of
$H_m(\GL_{n'}(\OO_L),\Z)$ has order $t(L,n,m)$), $T=H_k(\GL_n(\OO_L),\Z)_{\mathrm{tors}}$ has order $t(L,2k+6,k)$. Apply (4)(i) to the
group $Y$ of $(k+1)$-cycles of $\GL_n(\OO_L)$ with coefficients in $\Z/q$, $\beta$ of (3) and $g(y)=\ell^k\Phi_n(y)$. If $\beta(y)=0$,
then $y=\bar z$ for an integral cycle $z$, and by (1), (2), $g(y)=\ell^k\varphi_n([z])=\varphi_{\ell^kn}(u\,z_k([z]))$; so the values of $g$ on
$\ker\beta$ lie in $\varphi_{\ell^kn}(J)$, where $J\subset H=H_{k+1}(\GL_{\ell^kn}(\OO_L),\Z)$ is the subgroup generated by the $z_k(x)$.
By \cref{pf:vwin:bodd}, $J$ has rank at most $r_2(L)$, and $\#H_{\mathrm{tors}}=t(L,2k+6,k+1)$ by stability; so (4)(ii) gives
$c=q^{r_2(L)}t(L,2k+6,k+1)$, and (4)(i) the claim.
\end{proof}

\subsubsection{Detection of powers and few classes}

\begin{proposition}[$\mathrm S_{\mathrm{detect}}$]\label{pf:vwin:sdetect}
Let $F/L$ be an extension of number fields, $q=\ell^\nu$, $\zeta\in\OO_F^\times$ a primitive $q$-th root of unity, $\beta\in F^\times$ and
$0\ne b\in\OO_F$. If for every prime $P$ of $\OO_F$ with $b\notin P$ and $\OO_L\to\OO_F/P$ surjective the residue of $\beta$ is a $q$-th power
in $(\OO_F/P)^\times$, then $\beta$ is a $q$-th power in $F$.
\formallabel{LateralTheory.Fin.SDetect.intersectionNumber_baseChange}
\formallabel{UpperTheory.V.ZK.SDetect}
\formallabel{LateralTheory.Fin.SDetect.detection_statement_holds}
\formallabel{LateralTheory.Q.SDetectCut.goal}
\formallabel{LateralTheory.Q.SDetectCut.frob_fix}
\end{proposition}

\begin{proof}
Let $\Omega$ be an algebraic closure of $F$, $x\in\Omega$ with $x^q=\beta$, and $g\in\Gal(\Omega/F)$. The set of $\sigma$ with
$\sigma x=gx$ is open and non-empty. By Chebotarev's density theorem for primes of absolute degree one, it contains a Frobenius
element $\sigma$ at a place $B$ of $\Omega$ above a prime $v$ of $\OO_F$ of absolute degree one with $bq\notin v$; then $\OO_F/v=\F_p$
is the image of $\Z$, so $\OO_L\to\OO_F/v$ is surjective, and by hypothesis $\beta\equiv t^q$ modulo $v$ for a $v$-unit $t\in\OO_F$
(after clearing a denominator prime to $v$). As $x^q-t^q=\prod_{i<q}(x-\zeta^it)$ has positive $B$-valuation, $x\equiv c:=\zeta^it\pmod B$
for some $i$; in particular $x$ is a $B$-unit. Since $\sigma y\equiv y^p\pmod B$ for integral $y$ and $c^p\equiv c\pmod v$,
$\sigma x\equiv x^p\equiv c^p\equiv c\equiv x\pmod B$. But $\sigma x=\zeta^jx$ for some $j$, and $\zeta^j-1$ is a $B$-unit unless $\zeta^j=1$,
because $q\notin v$; hence $\sigma x=x$, i.e.\ $gx=x$. So $x$ is fixed by $\Gal(\Omega/F)$, $x\in F$, and $\beta=x^q$.
\end{proof}

\begin{proposition}[$\mathrm S_{\mathrm{few}}$]\label{pf:vwin:sfew}
Assume $\mathrm S_\theta$, $\mathrm S_\phi$, $\mathrm S_{\mathrm{class}}$, $\mathrm S_{\mathrm{upper}}$ and $\mathrm S_{\mathrm{detect}}$. Let
$L$ be totally imaginary, $F/L$ a layer with data $(\ell,c,\nu,\zeta,\gamma)$, $e\in\N$ with $2e+3<\ell^c$ and $2e+3<r_2(L)$, and
$a=(1+\ell^c)^{e+1}$. Then $\mathrm{Few}(\gamma,a,\ell^{b},\ell^\nu,B)$ holds with $b=2e+2+(2e+3)(c+2e+3)$ and
$B=(\ell^\nu)^{r_2(L)}\,t(L,4e+10,2e+3)\,t(L,4e+10,2e+2)$.
\formallabel{LateralTheory.Fin.SFewOf.intersectionNumber_baseChange}
\formallabel{UpperTheory.V.ZK.SFewOf}
\formallabel{UpperTheory.V.ZK.SFew}
\formallabel{LateralTheory.Fin.SFewOf.few_of_statement_holds}
\formallabel{LateralTheory.Q.SFewOfCut.sFewOf}
\formallabel{LateralTheory.Q.SFewOfCut.pe_dvd}
\formallabel{LateralTheory.Q.SFewOfCut.exists_index}
\formallabel{LateralTheory.Q.SFewOfCut.prim_map}
\end{proposition}

\begin{proof}
Let $q=\ell^\nu$ and $S$ a finite set as in the definition of $\mathrm{Few}$. For $\alpha\in S$ let $N_\alpha$, $y_\alpha$, $b_\alpha$ be as in
$\mathrm S_{\mathrm{class}}$, and let $I$ be the set of primes $P$ of $\OO_F$ not dividing $\ell\prod_\alpha b_\alpha$ with $\OO_L\to\OO_F/P$
surjective. For $P\in I$ let $k_P=\OO_F/P$; the residue of $\zeta$ is a primitive $q$-th root of unity, so $q$ divides $\#k_P^\times$ and there
is $v_P:k_P^\times\to\Z/q$ with kernel the $q$-th powers ($k_P^\times$ is cyclic). By $\mathrm S_\theta$ there is an additive family $\theta_P$ normalized
for $(\zeta\bmod P,v_P)$, and by $\mathrm S_\phi$ its pull-back along $\OO_L\to k_P$ is a primitive family $\Phi_P$. By $\mathrm S_{\mathrm{class}}$,
$\Phi_P(y_\alpha)=p_e\,v_P(u_{\alpha,P})$, $u_{\alpha,P}$ the residue of $\mathrm N_a(\alpha)$.

Each factor $(1+\ell^c)^{e+1}-(1+\ell^c)^j$ of $p_e$ has $\ell$-adic valuation $c+v_\ell(e+1-j)<c+2e+3$ (as $|e+1-j|<2e+3<\ell^c$), so
$\ell^{(2e+3)(c+2e+3)}$ is a multiple of $p_e$ in $\Z/q$. Let $f(\alpha)=(\ell^{2e+2}\Phi_P(y_\alpha))_{P\in I}\in(\Z/q)^I$. If $f(\alpha)=f(\alpha')$,
then for every $P\in I$, $\ell^{2e+2}p_e(v_P(u_{\alpha,P})-v_P(u_{\alpha',P}))=0$, hence $v_P\bigl((u_{\alpha,P}/u_{\alpha',P})^{\ell^b}\bigr)=0$, i.e.\ the
residue of $\mathrm N_a(\alpha)^{\ell^b}\mathrm N_a(\alpha')^{-\ell^b}$ is a $q$-th power. By $\mathrm S_{\mathrm{detect}}$ this element is a $q$-th power
in $F$, so $\alpha=\alpha'$ by the defining property of $S$. Thus $f$ is injective. The family $(\Phi_P)_{P\in I}$ is primitive with values in
$(\Z/q)^I$, and $f(S)$ consists of elements $\ell^{2e+2}\Phi(y)$; $\mathrm S_{\mathrm{upper}}$ with $k=2e+2$ gives $\#S=\#f(S)\le B$, as
$2(2e+3)+4=4e+10$.
\end{proof}

\subsubsection{Passing to a deeper layer, and torsion}

\begin{proposition}[$\mathrm S_{\mathrm{layer}}$]\label{pf:vwin:slayer}
Let $L$ be a number field containing a primitive $q_\ell$-th root of unity and $E\ge1$. There are a totally imaginary extension $L'$ of $L$ of
degree $d$ and $c\ge1$ such that $L'$ has top root $\ell^c$, $2E+1<\ell^c\le2[L':\Q]$, $2E+1<r_2(L')$, $d\le([L:\Q]+1)(4E+4)$,
$|d_{L'}|\le|d_L|^d(\ell^c)^{[L':\Q]}$ and $\vwinsh(L,\ell,E)\le\vwinsh(L',\ell,E)$.
\formallabel{LateralTheory.Fin.SLayer.intersectionNumber_baseChange}
\formallabel{UpperTheory.V.ZK.SLayer}
\formallabel{LateralTheory.Fin.SLayer.layer_statement_holds}
\formallabel{LateralTheory.Q.SLayerCut.goal_of}
\end{proposition}

\begin{proof}
Let $m$ be least with $\ell^m\ge4E+4$, so $\ell^m\le\ell(4E+4)$, and let $L'$ be the layer of degree $d=\ell^m$ over $L$ in the cyclotomic
$\ell$-tower, i.e.\ $L'=L(\mu_{\ell^{c_0+m}})$ with $c_0$ maximal such that $\mu_{\ell^{c_0}}\subset L$ (\cref{pf:vwin-he:layers}). Then $L'$ is
totally imaginary, contains $\mu_{q_\ell}$, and has top root $\ell^c$ with $c\ge c_0+m$, so $\ell^c\ge4E+4>2E+1$; $\ell^c\le2[L':\Q]$
since $[\Q(\mu_{\ell^c}):\Q]=\ell^{c-1}(\ell-1)\ge\ell^c/2$ (\lref{UpperTheory.V.Win.rootdeg}); $r_2(L')=[L':\Q]/2\ge\ell^m>2E+1$ (as $[L:\Q]\ge2$);
and $d=\ell^m\le\ell(4E+4)\le([L:\Q]+1)(4E+4)$, because $\ell-1=[\Q(\mu_\ell):\Q]\le[L:\Q]$ for odd $\ell$. The discriminant bound follows
from $L'=L(\zeta)$ with $\zeta$ a root of $X^{\ell^{c_0+m}}-1$, whose discriminant divides a power of $\ell^c$ of exponent at most the degree
(\lref{LateralTheory.Q.SLayerCut.disc_cyc}). Finally, by the descent isomorphism \cref{pf:vwin-he:desc} for $L_1=L$ and $L_1=L'$ with the
same group $U=G_L$, $\Sha(L,E)$ and $\Sha(L',E)$ are the modules of invariants of $G_L$, resp.\ of its subgroup $G_{L'}$, in the same module;
so $\vwinsh(L,\ell,E)\le\vwinsh(L',\ell,E)$ (\lref{LateralTheory.Q.SLayerCut.sha_le}, \lref{UpperTheory.V.Tower.exists_descends}).
\end{proof}

\begin{proposition}[$\mathrm S_{\mathrm{tors}}$]\label{pf:vwin:stors}
There is $C_0$ such that $\log t(L,4E+6,m)\le\kappa(C_0,E,[L:\Q],|d_L|)$ for all totally imaginary $L$, $E\ge1$ and $m\le2E+1$.
\formallabel{UpperTheory.V.ZK.sTors}
\formallabel{UpperTheory.V.ZK.STors}
\formallabel{UpperTheory.V.ZK.torsAll}
\end{proposition}

\begin{proof}
We show that $C_0=24$ works, for all $1\le N\le4E+6$. Let $n=[L:\Q]$ and $D=|d_L|$. The group $\GL_N(\OO_L)$ acts on an acyclic
complex of reduced $\OO_L$-lattices (Hermite normal forms), and a general bound for the torsion in the homology of a group acting on such
a complex gives $\log t(L,N,m)\le\exp(a+(m+1)\sigma+2m)$ whenever the number $A$ of orbits of the relevant cells satisfies $A\le e^a$ and
the constant $s=3^{nN^2}$ satisfies $s\le e^\sigma$ (\lref{UpperTheory.V.KT.tors_le},
\lref{UpperTheory.V.KH.TorsLe.log_le_exp}). Here $A\le(2DnN)^{12(m+1)(nN)^3N}$ (\lref{UpperTheory.V.KT.orbnum}) and
$s\le\exp(2nN^2)$. Inserting $N\le4E+6$ and $m\le2E+1$, the resulting bound is at most $\exp((2En)^{24}\log(2D))=\kappa(24,E,n,D)$
(\lref{UpperTheory.V.ZK.TorsionExponentBound.exponent_le}, an elementary estimate).\end{proof}


% ---------------------------------------------------------------------------
\subsection{A patching theorem for ordinary determinants}\label{sec:pf:vwin:patch}

The last root of this part of the development is a patching theorem of Taylor--Wiles type
\cite{TaylorWiles1995,Kisin2009,ClozelHarrisTaylor2008} for deformation rings of residually reducible ordinary
determinants in the sense of Chenevier \cite{Chenevier2014}, close in spirit to \cite{Thorne2012}. It is used, outside
this section, in the proofs of two automorphy statements for unitary groups in three variables
(\lref{UpperTheory.V.HV.Core.modU_of_printsH},
\lref{UpperTheory.V.Z.J.D.modularity_n_of_num_j_printed_mult_comp}). We state it and describe the structure of its
formal proof, whose six ingredients we do not reproduce.

\paragraph{Setting.} Let $\ell$ be an odd prime, $L$ a CM field with maximal totally real subfield $L^+$, $\Omega$ an algebraic closure of
$L$, and $\OO$ a coefficient ring for $\ell$. A \emph{Hida setting} fixes a complex conjugation, a finite set $S$ of places of $L^+$, and places of
$\Omega$ above the places of $L^+$, under the condition that the inertia groups of $L^+$ are contained in the image of $G_L$ (so that $L/L^+$ is
unramified at all finite places). A \emph{level} is a finite set $S'\supset S$ of places of $L^+$ which split in $L$, with the quotient
$\Gamma=G_{L^+,S'}$, the character $\eta:\Gamma\to\{\pm1\}$ cutting out $L$ (so $\eta(c)=-1$ for complex conjugation $c$), its kernel $G_0$, and the
cyclotomic character $\varepsilon:\Gamma\to\Zl^\times$. A \emph{residual datum} $r$ of dimension $n\ge1$ consists of $n$ distinct characters
$\bar\chi_1,\dots,\bar\chi_n$ of $G_0$ with values in the residue field of $\OO$, whose sum $\bar D$ is polarized, together with local characters at
the places of $S'$. For a complete local $\OO$-algebra $R$ with the same residue field, a \emph{pair} over $R$ is a continuous $n$-dimensional determinant
$D$ of $G_0$ with values in $R$ lifting $\bar D$, polarized, together with local characters lifting those of $r$, such that $D$ is ordinary at the places
above $\ell$ with respect to these characters and satisfies a condition at the inertia groups of $r$. The \emph{universal ordinary determinant ring} $B$ is
an $\OO$-algebra with a pair which is universal for maps to such $R$. A \emph{point} of $B$ is, roughly, a representation
$\rho:G_0\to\GL_n(K[[T]])$ ($K$ a finite extension of the residue field) with a symmetric polarization, ordinary flags at the places above $\ell$ and
the residual determinant $\bar D$, together with the homomorphism $f:B\to K[[T]]$ classifying its determinant; points come with niceness conditions of
two types (``type I'' and ``type E''). A prime $\mathfrak q$ of $B$ is \emph{stable} if the determinant at $\mathfrak q$ is that of an absolutely
irreducible representation, and \emph{thick} if it is not stable and a certain local algebra at $\mathfrak q$ is not semisimple.

\begin{theorem}[Patching]\label{pf:vwin:pfromaut}
Assume the Galois cohomology statements of group (c) of $\sfname{LiteratureRankNHypothesis}$ (local finiteness, local Tate duality, local Euler
characteristic, global finiteness, global Euler characteristic, the Greenberg--Wiles formula and normalized Poitou--Tate duality). Let
$\ell,L,\Omega$, a Hida setting, a level with $S'=S$, $\OO$, $n\ge1$, a residual datum $r$, the universal ring $B$ and an ideal $J\subset B$ be given.
Assume that the local degrees satisfy the class field theory condition $\sum_vd_v=[L^+:\Q]$ for suitable local abelian data at the places above $\ell$,
that $\varepsilon\bmod\ell$ is non-trivial on $G_0$, and that \emph{$J$ is automorphic at all Taylor--Wiles levels}: for every Taylor--Wiles datum (a
choice of auxiliary primes and of Frobenius elements, with the associated places) there is an automorphic patching datum for $B$ and $J$. Then:
\begin{enumerate}[label=(\alph*)]
\item if $x$ is a point of type I with $J\subset\ker f_x$, then every prime $\mathfrak q\subset\ker f_x$ of $B$ contains $J$;
\item if $x_0,\dots,x_h$ are points of type E such that consecutive points are linked by thick primes ($\mathfrak e_i\subset\ker f_{x_{i-1}}\cap\ker
f_{x_i}$), and some stable prime $\mathfrak q_1$ with $J\subset\mathfrak q_1\subset\ker f_{x_0}$ exists, then every stable or thick prime
$\mathfrak q\subset\ker f_{x_h}$ contains $J$.
\end{enumerate}
\formallabel{UpperTheory.V.Pa.pFromAut}
\formallabel{UpperTheory.V.Pa.pFromAut_of}
\end{theorem}

In other words, automorphy (membership in the closed subset defined by $J$) propagates from a nice automorphic point to all irreducible components of
$\Spec B$ through it, and along chains of components linked by thick primes. This is the conclusion of the Taylor--Wiles--Kisin method: the patched
module is supported on full components.

\begin{proof}[Structure of the proof]
The theorem is deduced (\lref{UpperTheory.V.Pa.chain_of}) from six statements of the development:
\begin{enumerate}[label=(P\arabic*)]
\item (\lref{UpperTheory.V.Pa.norm}.) Normalization: every nice point can be replaced, after a finite local injective extension of its coefficient
ring, by a point in normal form with the same determinant and local characters.
\item (\lref{UpperTheory.V.Pa.choice}.) Choice of Taylor--Wiles elements: for finitely many points there are elements $\sigma_1,\dots,\sigma_q$ of
the subgroup $G'$ (on which the residual characters are distinct) at which the characteristic polynomials of the points have distinct roots, and which
\emph{detect} the points (\lref{UpperTheory.V.Pa.detection}: a closed subset of $G'$ whose closure contains $G'$ detects, when $\varepsilon\bmod\ell$ is
non-trivial).
\item (\lref{UpperTheory.V.Pa.arPl}.) Existence of Taylor--Wiles places: for any finite family of elements of $G'$ there is a Taylor--Wiles datum
whose Frobenius elements are the given ones (Chebotarev's theorem).
\item (\lref{UpperTheory.V.Pa.b1Ar}.) A bound for the first-order dual (the number of generators of the relevant dual Selmer groups) at Taylor--Wiles
levels, by $n[L^+:\Q]+nq$, from the Galois cohomology statements and the class field theory condition.
\item (\lref{UpperTheory.V.Pa.b2Ar}.) A bound for the obstructions at Taylor--Wiles levels, from Poitou--Tate duality and the Euler characteristic
formulas.
\item (\lref{UpperTheory.V.Pa.raw}.) The patching argument proper: given automorphic patching data at all Taylor--Wiles levels and the two
bounds, the conclusion of (a) and (b) holds for the points of the patched ring.
\end{enumerate}
Combining (P1)--(P6), together with the comparison between stable primes of $B$ and of the patched rings
(\lref{UpperTheory.V.Pa.Det.stableIff}), gives (a) and (b).
\end{proof}


% ---------------------------------------------------------------------------
\subsection{Plain local components}\label{sec:pf:vwin:plain}

This subsection contains a local lemma on cuspidal representations which is used in the proof of the Kato-type
statements elsewhere in Part~II. Let $K$ be a totally imaginary number field, $n\ge1$, $\pi$ a cuspidal automorphic
subrepresentation of $\GL_n(\A_K)$ and $v$ a finite place of $K$ with residue field of order $q_v$ and uniformizer
$\varpi_v$. A finite multiset $s$ of characters of $K_v^\times$ is \emph{unlinked} if $(\xi/\xi')(\varpi_v)\ne q_v$ for all
$\xi,\xi'\in s$ such that $\xi/\xi'$ is unramified. An irreducible smooth representation of $\GL_n(K_v)$ is \emph{plain} if it
has a supercuspidal support consisting of characters of $\GL_1(K_v)=K_v^\times$ which is unlinked; $\pi$ is \emph{plain at
$v$} if all its local components at $v$ are plain. For $t\in\R$, $\pi$ \emph{satisfies the window at $(t,v)$} if every local
component of $\pi$ at $v$ has a supercuspidal support consisting of characters of $K_v^\times$, and every character $\xi$
occurring in such a support satisfies $q_v^{t-1/2}<|\xi(\varpi_v)|<q_v^{t+1/2}$.

\begin{lemma}\label{pf:vwin:plainof}
Let $K$, $n$, $\pi$ and $v$ be as above.
\begin{enumerate}[label=(\alph*)]
\item If $\pi$ satisfies the window at $(t,v)$ for some $t\in\R$, then $\pi$ is plain at $v$.
\item If $\pi$ has Satake parameters at $v$, then $\pi$ is plain at $v$.
\end{enumerate}
\formallabel{UpperTheory.V.Z.Win.plainOf}
\formallabel{UpperTheory.V.Z.Win.PlainOf}
\formallabel{UpperTheory.V.Z.Win.PlainAt}
\end{lemma}

\begin{proof}
Since supercuspidal supports are unique up to reordering, $\pi$ is plain at $v$ if and only if every local component
of $\pi$ at $v$ has a supercuspidal support consisting of characters of $K_v^\times$ and every such support is unlinked
(\lref{UpperTheory.V.Z.ImaginaryQuadraticEtaCharacter.plainIff}). (a) The first condition is part of the window.
If $\xi,\xi'$ occur in a support, then $|\xi(\varpi_v)|<q_v^{t+1/2}$ and $|\xi'(\varpi_v)|>q_v^{t-1/2}$, so
$|(\xi/\xi')(\varpi_v)|<q_v$, and in particular $(\xi/\xi')(\varpi_v)\ne q_v$
(\lref{UpperTheory.V.Z.ImaginaryQuadraticEtaCharacter.winUnl}). (b) If $\pi$ has Satake parameters at $v$, every local
component at $v$ is a subquotient of a principal series induced from unramified characters whose values at $\varpi_v$ are
the Satake parameters, so it has a support of unramified characters; and the Satake parameters of a cuspidal
representation are unlinked, since its local components are generic and an unramified generic representation is the
full, irreducible, induced representation (\lref{UpperTheory.V.Z.ImaginaryQuadraticEtaCharacter.sphUnl}). In both
cases the criterion above shows that $\pi$ is plain at $v$.
\end{proof}
